% Lesson 20 — Grammar: Expressing preferences — Anglais Tle A — Cours signé OnBuch+
% Programme officiel MINESEC. Compiler avec : tectonic EN20-grammar-expressing-preferences.tex
\newcommand{\DOCMATIERE}{Anglais}
\newcommand{\DOCNIVEAU}{Tle A}
\newcommand{\DOCMODULE}{Chapter 2 : Economic Life and Occupations}
\newcommand{\DOCLECON}{EN20}
\newcommand{\DOCTITRE}{Lesson 20 — Grammar: Expressing preferences}
% OnBuch+ — préambule « cours signés » ENRICHI (compilé avec tectonic / XeLaTeX)
% Système visuel « Pop » d'OnBuch+ : fond crème (jamais blanc), bordures encre
% épaisses, ombres « dures » décalées, titres Archivo Black en capitales.
% Version MATHÉMATIQUES : graphes (pgfplots/tikz), boîtes pédagogiques
% (définition, propriété, méthode, attention, le savais-tu, exercice résolu,
% exercice, TP, bilan), page de garde et signature OnBuch+ ancrée en bas.
%
% Macros attendues dans main.tex AVANT \input{preamble} :
%   \DOCMATIERE  \DOCNIVEAU  \DOCMODULE  \DOCLECON (numéro)  \DOCTITRE
\documentclass[11pt]{article}

\usepackage{fontspec}
\usepackage[english,french]{babel} % français = langue principale ; l'anglais via \eng{..} / textebox
\usepackage[a4paper,margin=2cm,top=3.4cm,bottom=2.8cm,headheight=1.4cm,footskip=1.3cm]{geometry}
\usepackage{amsmath,amssymb,amsfonts}
\usepackage{mathtools} % \xmapsto, \coloneqq... souvent utilisés par les modèles
\usepackage{cancel} % simplifications barrées (bilans de forces)
% \abs{z} (module, valeur absolue) et \norm{u} (norme), utilisés par les modèles
\providecommand{\abs}{}\let\abs\relax\DeclarePairedDelimiter{\abs}{\lvert}{\rvert}
\providecommand{\norm}{}\let\norm\relax\DeclarePairedDelimiter{\norm}{\lVert}{\rVert}
\usepackage[table]{xcolor} % [table] : fournit aussi \rowcolors (alternance de lignes)
\usepackage{pagecolor}
\usepackage{tikz}
\usetikzlibrary{arrows.meta,positioning,calc,shapes.geometric,shapes.misc,shapes.arrows,decorations.pathmorphing,decorations.pathreplacing,decorations.markings,patterns,shadows,mindmap,trees,fit,backgrounds,matrix,intersections,angles,quotes}
\usepackage{pgfplots}
\usepackage[version=4]{mhchem} % équations nucléaires \ce{^{235}_{92}U}
\usepackage[european,siunitx]{circuitikz} % schémas électriques
\pgfplotsset{compat=1.18}
\usepgfplotslibrary{fillbetween,groupplots} % aires entre courbes (calcul intégral)
\usepackage{enumitem}
\usepackage{fancyhdr}
% [most] charge toutes les bibliothèques tcolorbox (listings, minted, raster,
% documentation...) et gonfle inutilement la mémoire TeX sur un document long
% (~300 boîtes) : on ne charge que ce qui est réellement utilisé.
\usepackage[skins,breakable]{tcolorbox}
\usepackage{graphicx}
\usepackage{colortbl}
\usepackage{booktabs}
\usepackage{tabularx}
\usepackage{diagbox} % case d'en-tête diagonale des tableaux à double entrée
% Colonnes extensibles centrée / gauche / droite (C, L, R), souvent utilisées par les modèles
\newcolumntype{C}{>{\centering\arraybackslash}X}
\newcolumntype{L}{>{\raggedright\arraybackslash}X}
\newcolumntype{R}{>{\raggedleft\arraybackslash}X}
\usepackage{array}
\usepackage{multicol}
\usepackage{multirow}
\usepackage{siunitx}
% parse-numbers=false : accepte aussi \SI{2,5 \times 10^{-3}}{...}
\sisetup{locale=FR,output-decimal-marker={,},group-minimum-digits=4,parse-numbers=false}
\DeclareSIUnit\ohm{\ensuremath{\symup{\Omega}}} % U+2126 absent de la police
\DeclareSIUnit\division{div} % réglages d'oscilloscope (ms/div, V/div)
\DeclareSIUnit\tr{tr} % tours (vitesse de rotation, ex. \SI{1500}{\tr\per\minute})
\DeclareSIUnit\molar{M} % concentration molaire (ex. \SI{3}{\molar}, \SI{1}{\micro\molar})
\DeclareSIUnit\an{an} % année (le modèle confond parfois avec \year, un compteur TeX) — ex. \SI{5}{\kilo\meter\per\an}
\DeclareSIUnit\FCFA{FCFA} % franc CFA (ex. \SI{500}{\FCFA\per\kilo\gram})
\DeclareSIUnit\degreeBrix{\textdegree Brix} % teneur en sucre d'un sirop/jus (réfractométrie)
\DeclareSIUnit\month{mois} % le modèle utilise parfois \month, un compteur TeX, comme unité de durée
\usepackage{qrcode}
% \degree : commande fréquemment utilisée par les modèles pour un angle ou une
% température en dehors de \SI{}{} (non fournie par les paquets chargés).
\providecommand{\degree}{^{\circ}}
% \qed (fin de démonstration), utilisé par les modèles sans amsthm
\providecommand{\qed}{\hfill$\square$}
% \barre : double barre des valeurs interdites dans un tableau de variations (tkz-tab)
\providecommand{\barre}{\ensuremath{\|}}
% environnement proof (amsthm non chargé) utilisé par les modèles
\ifdefined\proof\else\newenvironment{proof}[1][Démonstration]{\par\noindent\textit{#1.}\ }{\qed\par}\fi
\usepackage{lastpage}
\usepackage{hyperref}
\hypersetup{hidelinks}

% ============================================================
% Palette « Pop » OnBuch+ — mêmes hex que la classe P (plus_ui.dart)
% ============================================================
\definecolor{poporange}{HTML}{F59321}
\definecolor{popdark}{HTML}{E07A0C}
\definecolor{popcream}{HTML}{FFF6E8}
\definecolor{popink}{HTML}{1C1714}
\definecolor{poppurple}{HTML}{7A5AE0}
\definecolor{popgreen}{HTML}{2E9E6A}
\definecolor{poppink}{HTML}{E0457B}
\definecolor{popblue}{HTML}{2F7DE1}
\definecolor{popgold}{HTML}{C98A12}
\definecolor{popmuted}{HTML}{8A7F76}
\definecolor{popline}{HTML}{EADFD2}
\definecolor{popgray}{HTML}{8A7F76}
\colorlet{popgrey}{popgray}
% Alias tolérants (le modèle nomme parfois les couleurs différemment)
\colorlet{popwhite}{white}
\colorlet{popblack}{popink}
\colorlet{popred}{poppink}
\colorlet{popyellow}{popgold}
% Teintes claires (fonds de tableaux/encadrés) — le modèle nomme parfois
% les couleurs avec un suffixe L (ex. popblueL) sans les avoir définies.
\colorlet{poporangeL}{poporange!20}
\colorlet{popdarkL}{popdark!20}
\colorlet{poppurpleL}{poppurple!20}
\colorlet{popgreenL}{popgreen!20}
\colorlet{poppinkL}{poppink!20}
\colorlet{popblueL}{popblue!20}
\colorlet{popgoldL}{popgold!20}
\colorlet{popredL}{popred!20}
\colorlet{popgrayL}{popgray!20}
% Teintes claires pour fonds de tableaux / zones
\colorlet{popblueL}{popblue!10!white}
\colorlet{poporangeL}{poporange!14!white}
\colorlet{popgreenL}{popgreen!12!white}
\colorlet{poppurpleL}{poppurple!12!white}
\colorlet{poppinkL}{poppink!12!white}
\colorlet{popgoldL}{popgold!14!white}

\pagecolor{popcream}

% ============================================================
% Typographie
% ============================================================
\defaultfontfeatures{Ligatures=TeX}
\setmainfont{PlusJakartaSans-Medium.ttf}[Path=../../../fonts/,BoldFont=PlusJakartaSans-Bold.ttf]
% Symboles, API et emojis absents de la police principale : repli sur DejaVu Sans (caractères actifs)
\newfontface\symfont{DejaVuSans.ttf}[Path=../../../fonts/]
\makeatletter
\newcommand\symactive[1]{\catcode#1=13 \begingroup\lccode`\~=#1 \lowercase{\endgroup\def~}{{\symfont\char#1}}}
\newcommand\symblank[1]{\catcode#1=13 \begingroup\lccode`\~=#1 \lowercase{\endgroup\def~}{}}
\newcommand\symhyph[1]{\catcode#1=13 \begingroup\lccode`\~=#1 \lowercase{\endgroup\def~}{\mbox{-}}}
\newcount\symc
\newcommand\symrange[2]{\symc=#1 \@whilenum\symc<#2 \do{\expandafter\symactive\expandafter{\the\symc}\advance\symc by 1 }}
\newcommand\symblankrange[2]{\symc=#1 \@whilenum\symc<#2 \do{\expandafter\symblank\expandafter{\the\symc}\advance\symc by 1 }}
\makeatother
\symrange{"0250}{"02FF}\symactive{"2713}\symactive{"2714}\symactive{"2717}\symactive{"2718}\symactive{"2610}\symactive{"2611}\symactive{"2612}
\symhyph{"2011}\symhyph{"2E1A}\symblankrange{"1F300}{"1FAFF}\symblank{"FE0F}
% Maths en sans-serif (Fira Math) avec chiffres et lettres droites de Plus
% Jakarta, pour que les formules se fondent dans le texte.
\usepackage[math-style=ISO,bold-style=ISO]{unicode-math}
\setmathfont{FiraMath-Regular.otf}[Path=../../../fonts/]
\setmathfont{PlusJakartaSans-Medium.ttf}[Path=../../../fonts/,range={up/num,up/{Latin,latin},\mathup}]
\usepackage{icomma} % virgule décimale sans espace en mode math
% Caractères mathématiques Unicode absents des polices de texte (Plus Jakarta,
% Archivo Black) : rendus par leur équivalent mathématique, en texte comme en
% formule — sinon ils s'affichent en carré vide (ex. « Arithmétique dans ℤ »).
\usepackage{newunicodechar}
\newunicodechar{ℝ}{\ensuremath{\mathbb{R}}}
\newunicodechar{ℤ}{\ensuremath{\mathbb{Z}}}
\newunicodechar{ℂ}{\ensuremath{\mathbb{C}}}
\newunicodechar{ℕ}{\ensuremath{\mathbb{N}}}
\newunicodechar{ℚ}{\ensuremath{\mathbb{Q}}}
\newunicodechar{𝔻}{\ensuremath{\mathbb{D}}}
\newunicodechar{α}{\ensuremath{\alpha}}
\newunicodechar{β}{\ensuremath{\beta}}
\newunicodechar{γ}{\ensuremath{\gamma}}
\newunicodechar{δ}{\ensuremath{\delta}}
\newunicodechar{ε}{\ensuremath{\varepsilon}}
\newunicodechar{θ}{\ensuremath{\theta}}
\newunicodechar{λ}{\ensuremath{\lambda}}
\newunicodechar{μ}{\ensuremath{\mu}}
\newunicodechar{σ}{\ensuremath{\sigma}}
\newunicodechar{τ}{\ensuremath{\tau}}
\newunicodechar{φ}{\ensuremath{\varphi}}
\newunicodechar{ψ}{\ensuremath{\psi}}
\newunicodechar{ω}{\ensuremath{\omega}}
\newunicodechar{Σ}{\ensuremath{\Sigma}}
% majuscules produites par \MakeUppercase dans les titres (α-aminés -> Α)
\newunicodechar{Α}{\ensuremath{\alpha}}
\newunicodechar{Β}{\ensuremath{\beta}}
\newunicodechar{Γ}{\ensuremath{\Gamma}}
\newunicodechar{Δ}{\ensuremath{\Delta}}
\newunicodechar{Ε}{\ensuremath{\varepsilon}}
\newunicodechar{Θ}{\ensuremath{\Theta}}
\newunicodechar{Λ}{\ensuremath{\Lambda}}
\newunicodechar{Μ}{\ensuremath{\mu}}
\newunicodechar{Τ}{\ensuremath{\tau}}
\newunicodechar{Φ}{\ensuremath{\Phi}}
\newunicodechar{Ψ}{\ensuremath{\Psi}}
\newunicodechar{Ω}{\ensuremath{\Omega}}
\newunicodechar{↦}{\ensuremath{\mapsto}}
\newunicodechar{∈}{\ensuremath{\in}}
\newunicodechar{∉}{\ensuremath{\notin}}
\newunicodechar{∧}{\ensuremath{\wedge}}
\newunicodechar{⇔}{\ensuremath{\Leftrightarrow}}
\newunicodechar{⟺}{\ensuremath{\Longleftrightarrow}}
\newunicodechar{⇒}{\ensuremath{\Rightarrow}}
\newunicodechar{⟹}{\ensuremath{\Longrightarrow}}
\newunicodechar{∘}{\ensuremath{\circ}}
\newunicodechar{∪}{\ensuremath{\cup}}
\newunicodechar{∩}{\ensuremath{\cap}}
\newunicodechar{≡}{\ensuremath{\equiv}}
\newunicodechar{⊂}{\ensuremath{\subset}}
\newunicodechar{⊥}{\ensuremath{\perp}}
\newunicodechar{∃}{\ensuremath{\exists}}
\newunicodechar{∀}{\ensuremath{\forall}}
\newunicodechar{∛}{\ensuremath{\sqrt[3]{\,}}}
\newunicodechar{∜}{\ensuremath{\sqrt[4]{\,}}}
\newunicodechar{⃗}{\ensuremath{{}^{\to}}}
\newunicodechar{ⁿ}{\ifmmode^{n}\else\textsuperscript{\ensuremath{n}}\fi}
\newunicodechar{ˣ}{\ifmmode^{x}\else\textsuperscript{\ensuremath{x}}\fi}
\newunicodechar{ᵇ}{\ifmmode^{b}\else\textsuperscript{\ensuremath{b}}\fi}
\newunicodechar{ʳ}{\ifmmode^{r}\else\textsuperscript{\ensuremath{r}}\fi}
\newunicodechar{ᵘ}{\ifmmode^{u}\else\textsuperscript{\ensuremath{u}}\fi}
\newunicodechar{ʸ}{\ifmmode^{y}\else\textsuperscript{\ensuremath{y}}\fi}
\newunicodechar{ᵃ}{\ifmmode^{a}\else\textsuperscript{\ensuremath{a}}\fi}
\newunicodechar{ᵉ}{\ifmmode^{e}\else\textsuperscript{\ensuremath{e}}\fi}
\newunicodechar{ᵏ}{\ifmmode^{k}\else\textsuperscript{\ensuremath{k}}\fi}
\newunicodechar{ᵖ}{\ifmmode^{p}\else\textsuperscript{\ensuremath{p}}\fi}
\newunicodechar{ᴺ}{\ifmmode^{N}\else\textsuperscript{\ensuremath{N}}\fi}
\newunicodechar{ᶜ}{\ifmmode^{c}\else\textsuperscript{\ensuremath{c}}\fi}
\newunicodechar{ʰ}{\ifmmode^{h}\else\textsuperscript{\ensuremath{h}}\fi}
\newunicodechar{ᵠ}{\ifmmode^{q}\else\textsuperscript{\ensuremath{q}}\fi}
\newunicodechar{ₙ}{\ifmmode_{n}\else\textsubscript{\ensuremath{n}}\fi}
\newunicodechar{ₐ}{\ifmmode_{a}\else\textsubscript{\ensuremath{a}}\fi}
\newunicodechar{ᵢ}{\ifmmode_{i}\else\textsubscript{\ensuremath{i}}\fi}
\newunicodechar{ᵣ}{\ifmmode_{r}\else\textsubscript{\ensuremath{r}}\fi}
\newunicodechar{ₖ}{\ifmmode_{k}\else\textsubscript{\ensuremath{k}}\fi}
\newunicodechar{ᵦ}{\ifmmode_{\beta}\else\textsubscript{\ensuremath{\beta}}\fi}
\newunicodechar{ₓ}{\ifmmode_{x}\else\textsubscript{\ensuremath{x}}\fi}
\newfontfamily\popdisplay{ArchivoBlack-Regular.ttf}[Path=../../../fonts/]
\newfontfamily\poplogo{SpaceGrotesk-Bold.ttf}[Path=../../../fonts/]
\newfontfamily\popsemi{PlusJakartaSans-SemiBold.ttf}[Path=../../../fonts/]
\newfontfamily\popxbold{PlusJakartaSans-ExtraBold.ttf}[Path=../../../fonts/]
\newfontfamily\popxboldls{PlusJakartaSans-ExtraBold.ttf}[Path=../../../fonts/,LetterSpace=3.0]

\setlist[enumerate]{leftmargin=1.7em,itemsep=5pt,topsep=4pt}
\setlist[itemize]{leftmargin=1.7em,itemsep=4pt,topsep=4pt,label={\color{poporange}\textbullet}}
\setlength{\parindent}{0pt}
\setlength{\parskip}{5pt}
\renewcommand{\arraystretch}{1.25}

% pgfplots : style Pop par défaut
\pgfplotsset{
  every axis/.append style={axis line style={popink,line width=1pt},tick style={popink},
    grid style={popline,line width=0.5pt},label style={font=\small},tick label style={font=\footnotesize}},
  popcurve/.style={poporange,line width=1.8pt,no markers,smooth},
}
% Même style disponible dans un \draw TikZ simple (hors pgfplots)
\tikzset{popcurve/.style={poporange,line width=1.8pt}}
\tikzset{popgrid/.style={popline,very thin}}
\colorlet{popcurve}{poporange} % le nom du style est parfois utilisé comme couleur

% ============================================================
% Pill / squiggle
% ============================================================
\newcommand{\poppill}[3]{%
  \tikz[baseline=-0.55ex]\node[draw=popink,line width=1.2pt,rounded corners=2.6pt,%
    fill=#2,inner xsep=6.5pt,inner ysep=3pt]{%
    \popxboldls\fontsize{7}{7}\selectfont\color{#3}\MakeUppercase{#1}};%
}
\newcommand{\popsquiggle}[1]{%
  \tikz[baseline=-0.4ex]{
    \draw[#1,line width=2.6pt,line cap=round]
      plot [smooth] coordinates {(0,0.09) (0.34,0.01) (0.68,0.21) (1.02,0.01) (1.36,0.10)};
  }%
}
% Mot-clé surligné (marqueur orange sous le texte)
\newcommand{\cle}[1]{{\popxbold\color{popdark}#1}}

% ============================================================
% En-tête / pied
% ============================================================
\pagestyle{fancy}
\fancyhf{}
\renewcommand{\headrulewidth}{0pt}
\renewcommand{\footrulewidth}{0pt}
\lhead{\raisebox{-0.15ex}{\poplogo\fontsize{13}{13}\selectfont\color{popink}OnBuch\color{poporange}+}}
\rhead{\poppill{\DOCMATIERE\ \textbullet\ \DOCNIVEAU\ \textbullet\ Leçon \DOCLECON}{white}{popink}}
\lfoot{\popxbold\fontsize{7.6}{7.6}\selectfont\color{popmuted}\MakeUppercase{Cours signé OnBuch+}\ \textbullet\ {\popsemi\DOCTITRE}}
\rfoot{\tikz[baseline=-0.6ex]\node[rounded corners=8.5pt,fill=popink,minimum height=17pt,minimum width=17pt,inner xsep=5pt,inner sep=0]{\popxbold\fontsize{8}{8}\selectfont\color{popcream}\thepage\,/\,\pageref*{LastPage}};}

% ============================================================
% Titres
% ============================================================
\newcommand{\chaptitre}[2]{%
  \begin{tcolorbox}[enhanced,colback=poporange,colframe=popink,boxrule=2.6pt,arc=11pt,
      drop shadow=popink,left=15pt,right=15pt,top=12pt,bottom=11pt,before skip=3pt,after skip=18pt]
    \poppill{#2}{popink}{popcream}\par\vspace{9pt}
    {\raggedright\hyphenpenalty=10000\exhyphenpenalty=10000\popdisplay\fontsize{22}{24}\selectfont\color{popcream}\MakeUppercase{#1}\par}
  \end{tcolorbox}
}
% Section numérotée : pastille orange avec le numéro + titre Archivo
\newcounter{popsec}
\newcommand{\coursec}[1]{%
  \stepcounter{popsec}\setcounter{popsub}{0}%
  \par\vspace{16pt}\noindent
  \begin{minipage}{\linewidth}
  \tikz[baseline=-0.7ex]\node[draw=popink,line width=1.6pt,fill=poporange,rounded corners=4pt,minimum size=22pt,inner sep=0]{\popdisplay\fontsize{12}{12}\selectfont\color{popcream}\thepopsec};%
  \hspace{8pt}{\popdisplay\fontsize{15}{17}\selectfont\color{popink}\MakeUppercase{#1}}\par
  \vspace{4pt}\noindent{\color{poporange}\rule{36pt}{3.6pt}}
  \end{minipage}\par\nopagebreak\vspace{8pt}
}
\newcounter{popsub}
\newcommand{\courssub}[1]{%
  \stepcounter{popsub}%
  \par\vspace{9pt}\noindent{\popxbold\fontsize{11.5}{13}\selectfont\color{popink}{\color{poporange}\thepopsec.\thepopsub}\hspace{6pt}#1}\par\nopagebreak\vspace{4pt}
}

% ============================================================
% Boîtes pédagogiques — même recette : fond blanc, bordure épaisse colorée,
% ombre dure encre, badge plein. Titre optionnel [#1].
% ============================================================
\tcbset{popbase/.style={breakable,enhanced,colback=white,boxrule=2.3pt,arc=9pt,
  shadow={3pt}{-3pt}{0pt}{popink},left=14pt,right=14pt,top=11pt,bottom=11pt,before skip=11pt,after skip=15pt,
  fontupper=\popsemi\fontsize{10.6}{14.5}\selectfont\color{popink},
  fontlower=\popsemi\fontsize{10.6}{14.5}\selectfont\color{popink}}}
% Coche dessinée (le glyphe ✓ n'existe pas dans les polices du document)
\newcommand{\popcheck}[1]{\tikz[baseline=-0.5ex]\draw[#1,line width=1.6pt,line cap=round,line join=round](-0.13,0.0)--(-0.04,-0.09)--(0.13,0.1);}
\providecommand{\coche}{\popcheck{popgreen}}
\providecommand{\resultat}[1]{\par\smallskip\noindent\fcolorbox{popgreen}{popgreenL}{\parbox{\dimexpr\linewidth-2\fboxsep-2\fboxrule\relax}{\textbf{Résultat.} #1}}\par\smallskip}
\newcommand{\popboxhead}[3]{\poppill{#1}{#2}{white}\ifx\relax#3\relax\else\hspace{6pt}{\popxbold\fontsize{10.6}{12}\selectfont\color{popink}#3}\fi\par\vspace{7pt}}

\newtcolorbox{aretenir}[1][]{popbase,colframe=popgreen,before upper={\popboxhead{À retenir}{popgreen}{#1}}}
% Texte en anglais (document de lecture, dialogue, extrait) : cadre
% d'étude. Un titre optionnel (source, auteur) s'affiche dans l'en-tête.
\newtcolorbox{textebox}[1][]{popbase,colframe=popblue,before upper={\popboxhead{Text}{popblue}{#1}\begin{otherlanguage}{english}},after upper={\end{otherlanguage}}}
% Marques ✓ / ✗ (forme correcte / fautive) souvent utilisées par les modèles
\providecommand{\cross}{\textcolor{poppink}{\ensuremath{\times}}}
\providecommand{\xmark}{\cross}\providecommand{\crossmark}{\cross}\providecommand{\cmark}{\ensuremath{\checkmark}}
% Ligne de réponse / justification à compléter (inventée par les modèles)
\providecommand{\justification}[1]{\par\noindent\textit{Justification~:} \dotfill\par}
% \textipa (package tipa non chargé) : la police ne couvre pas l'API -> simple passage
\providecommand{\textipa}[1]{#1}
% Soulignement (ulem non chargé) : \uline, \ul
\providecommand{\uline}[1]{\underline{#1}}\providecommand{\ul}[1]{\underline{#1}}
% \glossaire{\item ...} inventée par les modèles : liste de glossaire
\providecommand{\glossaire}[1]{\begin{itemize}#1\end{itemize}}
% \alert (beamer) inventée par les modèles : texte mis en évidence
\providecommand{\alert}[1]{\textbf{\textcolor{poppink}{#1}}}
% variantes ASCII d'ordinaux écrites en macro
\providecommand{\iere}{\textsuperscript{re}}\providecommand{\ieme}{\textsuperscript{e}}\providecommand{\ere}{\textsuperscript{re}}\providecommand{\eme}{\textsuperscript{e}}\providecommand{\er}{\textsuperscript{er}}\providecommand{\ie}{\textsuperscript{e}}
% Ordinaux français écrits en macro par les modèles : 1\ière, 3\ième
\newcommand{\ière}{\textsuperscript{re}}\newcommand{\ième}{\textsuperscript{e}}
% \exemple (inventée par les modèles) : étiquette « Exemple : »
\providecommand{\exemple}{\textbf{Exemple~:}\ }
% \square utilisable aussi hors mode math (cases à cocher)
\AtBeginDocument{\let\squareMath\square\renewcommand{\square}{\ensuremath{\squareMath}}}
% Mot ou phrase en anglais à l'intérieur d'un paragraphe français
\newcommand{\eng}[1]{\ifmmode\text{\textit{#1}}\else\begingroup\selectlanguage{english}\itshape #1\endgroup\fi}
\let\esp\eng % alias toléré
% flèches d'intonation inventées par les modèles
\providecommand{\textsearrow}{}\renewcommand{\textsearrow}{\ensuremath{\searrow}}\providecommand{\textnearrow}{}\renewcommand{\textnearrow}{\ensuremath{\nearrow}}
% \fr{..} : mot français (inventé par les modèles)
\providecommand{\fr}[1]{\textit{#1}}
\newtcolorbox{exemplebox}[1][]{popbase,colframe=poporange,before upper={\popboxhead{Exemple}{poporange}{#1}}}
\newtcolorbox{definition}[1][]{popbase,colframe=popblue,before upper={\popboxhead{Définition}{popblue}{#1}}}
\newtcolorbox{propriete}[1][]{popbase,colframe=poppurple,before upper={\popboxhead{Propriété}{poppurple}{#1}}}
\newtcolorbox{methode}[1][]{popbase,colframe=popgold,before upper={\popboxhead{Méthode}{popgold}{#1}}}
\newtcolorbox{attention}[1][]{popbase,colframe=poppink,before upper={\popboxhead{Attention}{poppink}{#1}}}
\newtcolorbox{experience}[1][]{popbase,colframe=popblue,colback=popblueL,before upper={\popboxhead{Expérience}{popblue}{#1}}}
% « Le savais-tu ? » : carte violette pleine (contexte, culture, Cameroun)
\newtcolorbox{savaistu}[1][]{popbase,colback=poppurple,colframe=popink,
  fontupper=\popsemi\fontsize{10.4}{14.2}\selectfont\color{white},
  before upper={\poppill{Le savais-tu\,?}{white}{poppurple}\ifx\relax#1\relax\else\hspace{6pt}{\popxbold\color{white}#1}\fi\par\vspace{7pt}}}
% Exercice résolu : énoncé en haut, \tcblower puis la solution sur fond crème
\newcounter{popexo}\newcounter{popexr}
\newtcolorbox{exoresolu}[1][]{popbase,colframe=poporange,colbacklower=popcream,
  segmentation style={popink,line width=1.2pt,dashed},
  before upper={\stepcounter{popexr}\popboxhead{Exercice résolu \thepopexr}{poporange}{#1}},
  before lower={\poppill{Solution}{popink}{popcream}\par\vspace{6pt}}}
% Exercice (non résolu en place) — la correction est dans la partie Corrigés
\newtcolorbox{exercice}[1][]{popbase,colframe=popink,
  before upper={\stepcounter{popexo}\popboxhead{Exercice \thepopexo}{popink}{#1}}}
% Corrigé
\newtcolorbox{corrige}[1][]{popbase,colframe=popgreen,colback=popgreenL,
  before upper={\popboxhead{Corrigé}{popgreen}{#1}}}
% Objectifs / prérequis / bilan
\newtcolorbox{objectifs}{popbase,colframe=popink,colback=poporangeL,
  before upper={\popboxhead{Objectifs de la leçon}{popink}{}}}
\newtcolorbox{prerequis}{popbase,colframe=popink,colback=popblueL,
  before upper={\popboxhead{Prérequis}{popblue}{}}}
\newtcolorbox{bilan}{popbase,colframe=popink,colback=white,boxrule=2.8pt,
  before upper={\popboxhead{Fiche bilan}{poporange}{L'essentiel à connaître pour l'examen}}}
% Figure encadrée avec légende
\newtcolorbox{popfigure}[1][]{enhanced,breakable=false,colback=white,colframe=popline,boxrule=1.4pt,arc=8pt,
  left=10pt,right=10pt,top=10pt,bottom=8pt,before skip=10pt,after skip=12pt,halign=center,
  lower separated=false,halign lower=center,
  fontlower=\popsemi\fontsize{9}{11.5}\selectfont\color{popmuted},
  #1}
\newcommand{\legende}[1]{\par\vspace{6pt}{\popsemi\fontsize{9}{11.5}\selectfont\color{popmuted}#1}}

% ============================================================
% Signature OnBuch+
% ============================================================
\newcommand{\poplogoinline}[1]{{\poplogo\fontsize{#1}{#1}\selectfont\color{popink}OnBuch\color{poporange}+}}
\newcommand{\signatureonbuch}{%
  \begin{tcolorbox}[enhanced,colback=popink,colframe=popink,boxrule=2.2pt,arc=10pt,
      drop shadow=poporange,left=14pt,right=14pt,top=10pt,bottom=10pt,before skip=0pt,after skip=0pt]
    \tikz[baseline=-0.7ex]\node[circle,fill=popgreen,draw=popcream,line width=1.3pt,minimum size=22pt,inner sep=0]{\popcheck{white}};%
    \hspace{9pt}\begin{minipage}[c]{0.62\linewidth}
      {\popxbold\fontsize{12}{13}\selectfont\color{popcream}Signé\ \textbullet\ {\poplogo OnBuch\color{poporange}+}}\par\vspace{2pt}
      {\raggedright\popsemi\fontsize{8}{10.5}\selectfont\color{popline}\MakeUppercase{Équipe pédagogique OnBuch+}\\\MakeUppercase{Conforme au programme officiel MINESEC}\par}
    \end{minipage}\hfill
    {\poplogo\fontsize{22}{22}\selectfont\color{popcream}OnBuch\color{poporange}+}
  \end{tcolorbox}%
}

% ============================================================
% Page de garde
% #1 titre · #2 matière · #3 niveau · #4 module · #5 n° leçon · #6 durée ·
% #7 description / objectifs (texte ou itemize)
% ============================================================
\newcommand{\pagegarde}[7]{%
  \thispagestyle{empty}
  \newgeometry{a4paper,margin=1.7cm,top=1.6cm,bottom=1.4cm}%
  \begin{tikzpicture}[remember picture,overlay]
    \fill[poporange!18] ([xshift=-3.2cm,yshift=-3.6cm]current page.north east) circle (4.6cm);
    \fill[poppurple!14] ([xshift=2.4cm,yshift=7.6cm]current page.south west) circle (3.8cm);
  \end{tikzpicture}%
  {\poplogo\fontsize{30}{30}\selectfont\color{popink}OnBuch\color{poporange}+}\hfill
  \raisebox{0.6ex}{\poppill{Cours signé \textbullet\ Premium}{popink}{popcream}}\par
  \vspace{1pt}\popsquiggle{poppurple}\par
  \vspace{8pt}
  \begin{tcolorbox}[enhanced,colback=poporange,colframe=popink,boxrule=2.8pt,arc=13pt,
      drop shadow=popink,left=18pt,right=18pt,top=12pt,bottom=13pt,after skip=10pt]
    \poppill{#2 \textbullet\ #3}{popink}{popcream}\hspace{5pt}\poppill{#4}{white}{popink}\par\vspace{8pt}
    {\popdisplay\fontsize{14}{14}\selectfont\color{popink}LEÇON #5}\par\vspace{4pt}
    {\raggedright\hyphenpenalty=10000\exhyphenpenalty=10000\popdisplay\fontsize{25}{27}\selectfont\color{popcream}\MakeUppercase{#1}\par}
    \vspace{8pt}
    {\popxbold\fontsize{9.5}{9.5}\selectfont\color{popink}\MakeUppercase{Durée conseillée : #6}}
  \end{tcolorbox}
  \begin{tcolorbox}[enhanced,colback=white,colframe=popink,boxrule=2.2pt,arc=10pt,
      drop shadow=popink,left=16pt,right=16pt,top=10pt,bottom=10pt,after skip=8pt]
    \poppill{Dans cette leçon}{poppurple}{white}\par\vspace{5pt}
    \popsemi\fontsize{9.3}{12.6}\selectfont\color{popink}#7
  \end{tcolorbox}
  \noindent\begin{minipage}[t]{0.487\linewidth}
    \begin{tcolorbox}[enhanced,colback=popgreen,colframe=popink,boxrule=2.2pt,arc=10pt,
        drop shadow=popink,left=11pt,right=11pt,top=10pt,bottom=10pt,height=3.1cm,valign=center]
      \tcbox[colback=white,colframe=popink,boxrule=1.3pt,arc=5pt,boxsep=0pt,nobeforeafter,
        left=3pt,right=3pt,top=3pt,bottom=3pt,valign=center]{\qrcode[height=1.9cm]{https://play.google.com/store/apps/details?id=com.luvvix.onbuch&hl=fr}}%
      \hspace{8pt}\begin{minipage}[c]{0.5\linewidth}
        {\popdisplay\fontsize{11}{12.5}\selectfont\color{white}\MakeUppercase{Télécharge OnBuch}}\par\vspace{4pt}
        {\raggedright\popsemi\fontsize{8.4}{11}\selectfont\color{white}Gratuit \textbullet\ Google Play\\Tuteur IA, annales, concours, résultats\par}
      \end{minipage}
    \end{tcolorbox}
  \end{minipage}\hfill
  \begin{minipage}[t]{0.487\linewidth}
    \begin{tcolorbox}[enhanced,colback=poppurple,colframe=popink,boxrule=2.2pt,arc=10pt,
        drop shadow=popink,left=11pt,right=11pt,top=10pt,bottom=10pt,height=3.1cm,valign=center]
      \tikz[baseline=-0.7ex]\node[circle,fill=white,draw=popink,line width=1.3pt,minimum size=1.6cm,inner sep=0]{\popdisplay\fontsize{24}{24}\selectfont\color{poppurple}+};%
      \hspace{8pt}\begin{minipage}[c]{0.6\linewidth}
        {\popdisplay\fontsize{11}{12.5}\selectfont\color{white}\MakeUppercase{Passe à OnBuch+}}\par\vspace{4pt}
        {\raggedright\popsemi\fontsize{8.4}{11}\selectfont\color{white}Tous les cours signés, Tuteur IA illimité, fiches PDF — onglet OnBuch+ de l'app\par}
      \end{minipage}
    \end{tcolorbox}
  \end{minipage}
  \vspace{8pt}
  \popcontenu
  \vfill
  \signatureonbuch
  \restoregeometry
  \newpage
}

% Rangée « Ce cours contient » (page de garde)
\newcommand{\poptile}[3]{% #1 couleur #2 gros texte #3 légende
  \begin{tcolorbox}[enhanced,colback=#1,colframe=popink,boxrule=2pt,arc=9pt,shadow={2.5pt}{-2.5pt}{0pt}{popink},
      left=6pt,right=6pt,top=9pt,bottom=9pt,halign=center,valign=center,height=1.75cm,width=\linewidth,nobeforeafter]
    {\popdisplay\fontsize{12.5}{13}\selectfont\color{white}\MakeUppercase{#2}}\par\vspace{3pt}
    {\centering\popsemi\fontsize{7.8}{9.5}\selectfont\color{white}#3\par}
  \end{tcolorbox}}
\newcommand{\popcontenu}{%
  \par\noindent{\popxboldls\fontsize{7.5}{7.5}\selectfont\color{popmuted}CE COURS SIGNÉ CONTIENT}\par\vspace{7pt}
  \noindent\begin{minipage}[t]{0.235\linewidth}\poptile{popblue}{Cours}{complet, illustré, pas à pas}\end{minipage}\hfill
  \begin{minipage}[t]{0.235\linewidth}\poptile{poporange}{Méthodes}{et exercices résolus}\end{minipage}\hfill
  \begin{minipage}[t]{0.235\linewidth}\poptile{poppink}{Exercices}{type examen + corrigés}\end{minipage}\hfill
  \begin{minipage}[t]{0.235\linewidth}\poptile{popgreen}{Fiche bilan}{l'essentiel à retenir}\end{minipage}\par}

% Sommaire
\newtcolorbox{sommaire}{enhanced,colback=white,colframe=popink,boxrule=2.2pt,arc=10pt,
  drop shadow=popink,left=18pt,right=18pt,top=13pt,bottom=13pt,before skip=6pt,after skip=20pt,
  before upper={\poppill{Sommaire}{poppurple}{white}\par\vspace{9pt}\popsemi\fontsize{10.6}{14.5}\selectfont\color{popink}}}

% Signature de fin de cours — poussée tout en bas de la dernière page
\newcommand{\findecours}{%
  \par\vfill\nopagebreak
  \begin{center}\popsquiggle{poporange}\end{center}\vspace{4pt}
  \signatureonbuch
}
\AtBeginDocument{\def\llbracket{\mathopen{[\mkern-3.5mu[}}\def\rrbracket{\mathclose{]\mkern-3.5mu]}}} % intervalles d'entiers (glyphe ⟦ absent de la police)
% Glyphes absents de Fira Math : redéfinis à partir de glyphes disponibles
\AtBeginDocument{%
  \def\setminus{\mathbin{\backslash}}\let\smallsetminus\setminus
  \def\vdots{\mathord{\vbox{\baselineskip3.2pt\lineskiplimit0pt\kern4pt\hbox{.}\hbox{.}\hbox{.}}}}%
  \def\ddots{\mathinner{\mkern1mu\raise6.5pt\hbox{.}\mkern2mu\raise3.6pt\hbox{.}\mkern2mu\raise0.7pt\hbox{.}\mkern1mu}}%
  \def\triangle{\mathord{\scalebox{0.9}{$\symup{\Delta}$}}}%
  \def\nabla{\mathord{\raisebox{\depth}{\scalebox{1}[-1]{$\symup{\Delta}$}}}}%
  \def\checkmark{\popcheck{popdark}}%
  \def\star{\mathbin{\ast}}\def\bigstar{\mathord{\ast}}%
  \def\diamond{\mathbin{\scalebox{0.6}{$\Diamond$}}}%
  \def\bigcup{\mathop{\mathchoice{\raisebox{-0.3em}{\scalebox{1.7}{$\cup$}}}{\scalebox{1.25}{$\cup$}}{\cup}{\cup}}\displaylimits}%
  \def\bigcap{\mathop{\mathchoice{\raisebox{-0.3em}{\scalebox{1.7}{$\cap$}}}{\scalebox{1.25}{$\cap$}}{\cap}{\cap}}\displaylimits}%
  \def\lesseqqgtr{\mathrel{\lessgtr}}\def\bigoplus{\mathop{\oplus}\displaylimits}%
}
\providecommand{\ssi}{si et seulement si} % abréviation fréquente
\AtBeginDocument{\providecommand{\wideparen}{\overparen}} % arc de cercle

%% FIGFIX-BEGIN v5 — correctifs globaux des figures (docs/onbuch-generation/tools/figfix)
\usepackage{adjustbox}\usepackage{etoolbox}\tcbuselibrary{raster}
% toute figure TikZ/pgfplots plus large que la ligne est réduite à la largeur disponible
\BeforeBeginEnvironment{tikzpicture}{\begin{adjustbox}{max width=\linewidth}}
\AfterEndEnvironment{tikzpicture}{\end{adjustbox}}
% légendes pgfplots lisibles par-dessus les courbes
\pgfplotsset{every axis/.append style={legend style={fill=white,fill opacity=0.92,text opacity=1,draw=black!15,inner sep=3pt}}}
% étiquettes d'arêtes (syntaxe edge["..."]) avec fond blanc
\tikzset{every edge quotes/.append style={fill=white,inner sep=1.2pt}}
%% FIGFIX-END
\begin{document}

\pagegarde{Lesson 20 — Grammar: Expressing preferences}{Anglais}{Tle A}{Chapter 2 : Economic Life and Occupations}{EN20}{2 h}{Cette leçon de grammaire présente les différentes structures anglaises pour exprimer ses préférences : prefer, would rather, would prefer, like... better than et les comparatifs de préférence. À travers l'observation d'exemples, des règles illustrées et des exercices progressifs, les élèves apprennent à comparer des activités, des métiers et des choix de vie sans tomber dans les pièges des francophones.
\vspace{4pt}
\begin{multicols}{2}\small
\begin{itemize}
\item À la fin de cette leçon, je sais identifier les structures d'expression des préférences dans un texte ou un dialogue
\item À la fin de cette leçon, je sais construire correctement prefer + nom/gérondif et prefer X to Y
\item À la fin de cette leçon, je sais utiliser would rather + base verbale et would rather X than Y
\item À la fin de cette leçon, je sais employer would prefer + to-infinitive et le distinguer de prefer + -ing
\item À la fin de cette leçon, je sais exprimer une préférence avec like... better than et les comparatifs
\item À la fin de cette leçon, je sais éviter les erreurs fréquentes des francophones (prefer... than, to prefer do, rather to do)
\end{itemize}\end{multicols}}

\begin{sommaire}
\begin{multicols}{2}
\begin{enumerate}
\item Observation et découverte
\item Prefer : règle et emplois
\item Would rather et would prefer
\item Like... better than et autres structures
\item Comparaison français/anglais et pièges
\item Entraînement et production guidée
\item Activité d'intégration
\item Méthodes et exercices résolus
\item Exercices
\item Corrigés des exercices
\item Fiche bilan
\end{enumerate}
\end{multicols}
\end{sommaire}

\begin{objectifs}
\begin{itemize}
\item À la fin de cette leçon, je sais identifier les structures d'expression des préférences dans un texte ou un dialogue
\item À la fin de cette leçon, je sais construire correctement prefer + nom/gérondif et prefer X to Y
\item À la fin de cette leçon, je sais utiliser would rather + base verbale et would rather X than Y
\item À la fin de cette leçon, je sais employer would prefer + to-infinitive et le distinguer de prefer + -ing
\item À la fin de cette leçon, je sais exprimer une préférence avec like... better than et les comparatifs
\item À la fin de cette leçon, je sais éviter les erreurs fréquentes des francophones (prefer... than, to prefer do, rather to do)
\item À la fin de cette leçon, je sais exprimer mes préférences professionnelles et économiques dans des phrases et de courts textes
\item À la fin de cette leçon, je sais transformer une phrase d'une structure à une autre sans changer le sens
\end{itemize}
\end{objectifs}

\begin{prerequis}
\begin{itemize}
\item Le gérondif (-ing) et l'infinitif avec to (Première)
\item Les comparatifs d'adjectifs et d'adverbes (better, more... than)
\item Les auxiliaires modaux de base : would, can, must
\item Le vocabulaire des métiers et de la vie économique (Chapter 2, leçons précédentes)
\item La conjugaison des verbes au présent simple et au conditionnel avec would
\end{itemize}
\end{prerequis}

\newpage

\coursec{Observation et découverte}

\courssub{Situation d'accroche : la pause à l'école de Yaoundé}

C'est la récréation de 10 heures au lycée de Yaoundé. Sous le grand manguier de la cour, \textbf{Amina} et \textbf{Peter}, tous deux en Terminale A, parlent de leur avenir. Le bruit des bendskins et les cris des vendeurs de beignets au coin de la rue ne les distraient pas : le Bac approche, et le choix de carrière devient urgent.

\begin{textebox}[Dialogue — Career plans after the Bac]
Amina: “What job would you like after the Bac, Peter?”
Peter: “I'd rather work in a bank than teach. I prefer dealing with figures to marking essays.”
Amina: “Really? I like teaching better than banking. I'd prefer to work with young people.”
Peter: “My father says I should choose a stable career. He'd rather I studied economics than law.”
Amina: “My parents would prefer me to become a nurse, but I enjoy explaining things to others.”
Peter: “Well, we still have time to decide. I'd rather not rush my choice.”
Amina: “Neither would I. I prefer thinking carefully to making a quick mistake.”
\end{textebox}

\cle{Glossaire}
\begin{itemize}
    \item \eng{the Bac} (« le bæk ») : le Baccalauréat
    \item \eng{deal with} (« diil wiiz ») : manipuler, traiter, s'occuper de
    \item \eng{figures} (« fi-gyurz ») : les chiffres, les nombres
    \item \eng{marking essays} (« mar-king e-seiz ») : corriger des dissertations / copies
    \item \eng{stable career} (« stei-bl ca-ri-yur ») : carrière stable
    \item \eng{rush} (« rach ») : se précipiter, faire à la hâte
    \item \eng{neither would I} (« ni-der woud ai ») : moi non plus (accord négatif)
\end{itemize}

\courssub{Problématique}

Comment exprimer clairement ses goûts et ses choix en anglais ? Dans la vie économique et professionnelle — entretiens d'embauche, choix de carrière, achats au marché de Mokolo — savoir exprimer une préférence est indispensable. Cette leçon vous donne toutes les structures pour le faire correctement.

\courssub{Méthode d'observation du dialogue}

\begin{methode}[Comment repérer les structures de préférence]
\begin{enumerate}
    \item \textbf{Soulignez} le verbe ou l'expression qui marque la préférence (ex. \eng{prefer}, \eng{'d rather}, \eng{like ... better}).
    \item \textbf{Encadrez} l'élément \textbf{préféré} (ce que le locuteur veut).
    \item \textbf{Entourez} l'élément \textbf{rejeté} ou mis en second (ce que le locuteur refuse ou aime moins).
    \item \textbf{Notez} le mot de liaison utilisé : \eng{to}, \eng{than}, \eng{over}, \eng{to} + infinitif, \eng{-ing}...
    \item \textbf{Classez} chaque occurrence dans l'une des trois familles : \cle{PREFER}, \cle{WOULD RATHER}, \cle{LIKE ... BETTER THAN}.
\end{enumerate}
\end{methode}

\courssub{Application de la méthode sur le dialogue}

\begin{exemplebox}[Analyse guidée des premières répliques]
\textbf{1.} \eng{Peter: “I'd rather \underline{work} in a bank \framebox{than teach}.”}
\begin{itemize}
    \item Verbe de préférence : \eng{'d rather} (= \eng{would rather})
    \item Préféré : \eng{work in a bank}
    \item Rejeté : \eng{teach}
    \item Liaison : \eng{than} + \textbf{base verbale} (sans \eng{to})
\end{itemize}

\textbf{2.} \eng{Peter: “I \underline{prefer} \framebox{dealing with figures} \framebox{to} \framebox{marking essays}.”}
\begin{itemize}
    \item Verbe de préférence : \eng{prefer} (présent simple)
    \item Préféré : \eng{dealing with figures} (gérondif)
    \item Rejeté : \eng{marking essays} (gérondif)
    \item Liaison : \eng{to} + \textbf{gérondif}
\end{itemize}

\textbf{3.} \eng{Amina: “I \underline{like} \framebox{teaching} \underline{better than} \framebox{banking}.”}
\begin{itemize}
    \item Structure : \eng{like ... better than}
    \item Préféré : \eng{teaching} (gérondif)
    \item Rejeté : \eng{banking} (gérondif)
    \item Liaison : \eng{better than} + \textbf{gérondif}
\end{itemize}

\textbf{4.} \eng{Amina: “I'd \underline{prefer} \framebox{to work} with young people.”}
\begin{itemize}
    \item Structure : \eng{would prefer} + \textbf{infinitif avec \eng{to}}
    \item Préféré : \eng{to work with young people}
    \item (Pas d'élément rejeté explicite ici : préférence absolue)
\end{itemize}
\end{exemplebox}

\courssub{Classification des structures relevées}

\begin{center}
\begin{tabularx}{\linewidth}{|X|X|X|}
\hline
\rowcolor{popblueL}
\textbf{Famille PREFER} & \textbf{Famille WOULD RATHER} & \textbf{Famille LIKE ... BETTER THAN} \\
\hline
\eng{prefer dealing ... to marking} (gérondif) & \eng{'d rather work ... than teach} (base verbale) & \eng{like teaching better than banking} (gérondif) \\
\eng{'d prefer to work} (infinitif) & \eng{'d rather I studied ... than ...} (subjonctif / prétérit modal) & \eng{enjoy explaining} (verbe + gérondif, sens proche) \\
\eng{prefer thinking ... to making} (gérondif) & \eng{'d rather not rush} (négation) &  \\
\hline
\end{tabularx}
\end{center}

\courssub{Idée générale dégagée}

\begin{aretenir}[Ce que l'observation nous apprend]
L'anglais possède \textbf{plusieurs structures} pour exprimer la préférence, selon :
\begin{itemize}
    \item le \textbf{degré} de la préférence (générale / habituelle vs ponctuelle / situationnelle) ;
    \item le \textbf{contexte} (choix entre deux alternatives précises vs goût général) ;
    \item la \textbf{grammaire} qui suit (gérondif, infinitif, base verbale, subordonnée).
\end{itemize}
Les trois familles principales sont :
\begin{enumerate}
    \item \cle{PREFER} (présent simple ou \eng{would prefer}) — préférence générale ou polie.
    \item \cle{WOULD RATHER} (\eng{'d rather}) — choix précis, souvent entre deux options, plus fort.
    \item \cle{LIKE ... BETTER THAN} — goût personnel, comparatif de supériorité.
\end{enumerate}
\end{aretenir}

\courssub{Premières hypothèses de règles (formulées par les élèves)}

\begin{experience}[Activité de classe : « Hypothèses avant leçon »]
En binôme, complétez le tableau ci-dessous à partir \textbf{uniquement} des exemples du dialogue. Ne consultez pas encore la règle officielle.

\begin{center}
\begin{tabularx}{\linewidth}{|X|X|X|}
\hline
\rowcolor{poporangeL}
\textbf{Structure} & \textbf{Forme qui suit (préféré)} & \textbf{Mot de liaison + forme (rejeté)} \\
\hline
\eng{prefer} + V-\eng{ing} & \eng{dealing with figures} & \eng{to} + V-\eng{ing} (\eng{marking essays}) \\
\hline
\eng{would prefer} + \eng{to} + V & \eng{to work with young people} & (souvent absent) \\
\hline
\eng{would rather} + V (base) & \eng{work in a bank} & \eng{than} + V (base) (\eng{teach}) \\
\hline
\eng{like} + V-\eng{ing} + \eng{better than} & \eng{teaching} & \eng{better than} + V-\eng{ing} (\eng{banking}) \\
\hline
\end{tabularx}
\end{center}

\textbf{Questions guides pour la discussion :}
\begin{enumerate}
    \item Pourquoi \eng{prefer} est-il suivi de \eng{-ing} dans \eng{prefer dealing} mais de \eng{to} + infinitif dans \eng{'d prefer to work} ?
    \item Quelle différence de sens entre \eng{I prefer tea} et \eng{I'd prefer tea} ?
    \item Pourquoi \eng{would rather} s'emploie-t-il sans \eng{to} (\eng{rather work}) ?
    \item Que devient la structure quand le sujet change ? (\eng{He'd rather I studied}...)
\end{enumerate}
\end{experience}

\courssub{Carte mentale : les trois familles de la préférence}

\begin{popfigure}
\begin{tikzpicture}[
    mindmap,
    grow cyclic,
    every node/.style={concept, execute at begin node=\hskip0pt, align=center},
    concept color=popblue!70,
    level 1/.style={level distance=4.5cm, sibling angle=120},
    level 2/.style={level distance=3cm, sibling angle=45},
]
\node[concept, text width=3.2cm, font=\bfseries\large] {Expressing\\ Preferences}
    child[concept color=poporange!80] {
        node[concept, text width=2.8cm, font=\bfseries] {PREFER}
        child[concept color=poporange!50] { node[concept, text width=2.2cm] {prefer + V-ing} }
        child[concept color=poporange!50] { node[concept, text width=2.2cm] {prefer + to + V} }
        child[concept color=poporange!50] { node[concept, text width=2.2cm] {would prefer + to + V} }
    }
    child[concept color=popgreen!70] {
        node[concept, text width=2.8cm, font=\bfseries] {WOULD RATHER}
        child[concept color=popgreen!50] { node[concept, text width=2.2cm] {'d rather + V (base)} }
        child[concept color=popgreen!50] { node[concept, text width=2.2cm] {'d rather + not + V} }
        child[concept color=popgreen!50] { node[concept, text width=2.2cm] {subject + 'd rather + past} }
    }
    child[concept color=poppurple!70] {
        node[concept, text width=2.8cm, font=\bfseries] {LIKE ... BETTER}
        child[concept color=poppurple!50] { node[concept, text width=2.2cm] {like + V-ing + better than} }
        child[concept color=poppurple!50] { node[concept, text width=2.2cm] {prefer + noun + to + noun} }
        child[concept color=poppurple!50] { node[concept, text width=2.2cm] {enjoy + V-ing (synonyme)} }
    };
\end{tikzpicture}
\legende{Carte mentale — Les trois grandes familles pour exprimer la préférence en anglais}
\end{popfigure}

\courssub{Exercice résolu : identification et classification}

\begin{exoresolu}[Repérer et classer les structures de préférence]
\textbf{Énoncé :} Dans les phrases ci-dessous, issues d'un forum d'orientation au Cameroun, \textbf{soulignez} le marqueur de préférence, \textbf{encadrez} l'élément préféré, \textbf{entourez} l'élément rejeté (s'il y a lieu) et \textbf{indiquez la famille} (PREFER / WOULD RATHER / LIKE BETTER).

\begin{enumerate}
    \item \eng{I prefer studying economics to studying law.}
    \item \eng{She'd rather work in Douala than in Yaoundé.}
    \item \eng{We like working outdoors better than sitting in an office.}
    \item \eng{My brother would prefer to start his own business.}
    \item \eng{They'd rather not take a loan from the bank.}
    \item \eng{Do you prefer tea or coffee?}
\end{enumerate}

\tcblower
\textbf{Solution détaillée :}

\begin{enumerate}
    \item \eng{I \underline{prefer} \framebox{studying economics} \framebox{to} \framebox{studying law}.} $\rightarrow$ Famille \textbf{PREFER} (présent simple + V-ing / to + V-ing). \textit{Traduction : « Je préfère étudier l'économie plutôt que le droit. »}
    \item \eng{She'd \underline{rather} \framebox{work in Douala} \framebox{than} \framebox{work in Yaoundé}.} $\rightarrow$ Famille \textbf{WOULD RATHER} ('d rather + base verbale / than + base verbale). \textit{« Elle préférerait travailler à Douala plutôt qu'à Yaoundé. »}
    \item \eng{We \underline{like} \framebox{working outdoors} \underline{better than} \framebox{sitting in an office}.} $\rightarrow$ Famille \textbf{LIKE BETTER} (like + V-ing + better than + V-ing). \textit{« Nous aimons mieux travailler dehors que rester assis au bureau. »}
    \item \eng{My brother \underline{would prefer} \framebox{to start his own business}.} $\rightarrow$ Famille \textbf{PREFER} (would prefer + to + infinitif). Pas d'élément rejeté explicite. \textit{« Mon frère préférerait créer sa propre entreprise. »}
    \item \eng{They'd \underline{rather not} \framebox{take a loan from the bank}.} $\rightarrow$ Famille \textbf{WOULD RATHER} (négation : 'd rather not + base verbale). \textit{« Ils préféreraient ne pas contracter de prêt à la banque. »}
    \item \eng{Do you \underline{prefer} \framebox{tea} \framebox{or} \framebox{coffee}?} $\rightarrow$ Famille \textbf{PREFER} (présent simple + nom / or + nom). Structure nominale : \eng{prefer noun to noun} ou \eng{prefer noun or noun} à l'interrogatif. \textit{« Préférez-vous le thé ou le café ? »}
\end{enumerate}
\end{exoresolu}

\courssub{Premier piège à éviter : confusion \eng{to} / \eng{than}}

\begin{attention}[Erreur fréquente des francophones : « Je préfère X \textbf{à} Y »]
En français, on dit « Je préfère le thé \textbf{à} le café » (préposition \textbf{à}). En anglais, la préposition change selon la structure :
\begin{itemize}
    \item \eng{prefer} + V-\eng{ing} \textbf{to} + V-\eng{ing} (\eng{prefer reading to writing})
    \item \eng{prefer} + nom \textbf{to} + nom (\eng{prefer tea to coffee})
    \item \eng{would rather} + V (base) \textbf{than} + V (base) (\eng{'d rather read than write})
    \item \eng{like} + V-\eng{ing} \textbf{better than} + V-\eng{ing} (\eng{like reading better than writing})
\end{itemize}
\textbf{Ne dites jamais :} \eng{*I prefer reading than writing} $\times$ \quad \eng{*I'd rather read to write} $\times$
\end{attention}

\courssub{Synthèse de la phase d'observation}

\begin{aretenir}[Bilan avant la règle formelle]
\begin{itemize}
    \item Trois familles principales : \cle{PREFER}, \cle{WOULD RATHER}, \cle{LIKE ... BETTER THAN}.
    \item Chaque famille impose sa \textbf{grammaire propre} (gérondif, infinitif, base verbale, prétérit modal).
    \item Le choix dépend du \textbf{contexte} (général vs ponctuel) et de la \textbf{force} de la préférence.
    \item Les francophones confondent systématiquement \eng{to} et \eng{than} : \textbf{apprenez les couples inséparables} (\eng{prefer ... to}, \eng{rather ... than}, \eng{better than}).
\end{itemize}
\end{aretenir}

\coursec{Prefer : règle et emplois}

Dans la section précédente, nous avons observé un dialogue au marché de Mokolo et relevé plusieurs phrases exprimant des préférences : \eng{I prefer farming to trading.}, \eng{She prefers to sell in the morning.}, \eng{We prefer buying local products.} Toutes ces phrases tournent autour d'un même verbe : \cle{prefer}. Il est temps maintenant de dégager la règle précise de ce verbe, ses différentes constructions, et de comprendre pourquoi les francophones se trompent si souvent avec lui.

\courssub{Le verbe prefer : présentation générale}

\begin{definition}[Le verbe prefer]
\eng{Prefer} est un verbe régulier qui signifie « préférer », « aimer mieux ». Il exprime une \cle{préférence générale}, un goût durable, une habitude de choix. Il se conjugue comme un verbe régulier, mais attention à l'orthographe : le \cle{r} final \textbf{redouble} devant une terminaison commençant par une voyelle.
\end{definition}

\begin{propriete}[Conjugaison et orthographe de prefer]
\begin{center}
\begin{tabular}{|l|l|l|}
\hline
\rowcolor{popblueL} Forme & Exemple & Remarque \\
\hline
Présent & I prefer, she prefer\textbf{s} & 3\textsuperscript{e} personne : -s \\
\hline
Prétérit & she prefer\textbf{red} & redoublement du r + -ed \\
\hline
Participe passé & she has prefer\textbf{red} & redoublement du r \\
\hline
Forme en -ing & prefer\textbf{ring} & redoublement du r \\
\hline
\end{tabular}
\end{center}
Le r redouble parce que l'accent tonique tombe sur la dernière syllabe (pre-\textbf{FER}). Comparez : \eng{offer} (accent sur la première syllabe) donne \eng{offered} sans redoublement.
\end{propriete}

\begin{attention}[Orthographe]
N'écrivez jamais \eng{prefered} ou \eng{prefering} : c'est une faute classique en rédaction. La bonne forme est \eng{preferred}, \eng{preferring}. Prononciation : « pri-FEUR » (« pri-feur-d » au passé).
\end{attention}

\courssub{Structure 1 : prefer + nom}

C'est la construction la plus simple : on exprime une préférence générale pour une chose, une personne, une catégorie.

\begin{propriete}[Prefer + nom]
\textbf{Forme :} sujet + prefer + nom (ou groupe nominal).\\
\textbf{Emploi :} préférence générale, durable, pour une chose.
\end{propriete}

\begin{exemplebox}[Prefer + nom]
\eng{I prefer tea.} « Je préfère le thé. »\\
\eng{She prefers local food.} « Elle préfère la cuisine locale. »\\
\eng{Most traders in Douala prefer cash payments.} « La plupart des commerçants de Douala préfèrent les paiements en espèces. »\\
\eng{My father preferred the quiet life of the village.} « Mon père préférait la vie tranquille du village. »
\end{exemplebox}

Remarquez qu'en anglais le nom suit directement le verbe, sans préposition : \eng{I prefer tea}, alors qu'en français on dit « je préfère \textbf{le} thé » avec l'article défini. L'anglais n'emploie généralement pas d'article quand on parle d'une catégorie en général.

\courssub{Structure 2 : prefer + gérondif (-ing)}

Pour parler d'une \cle{activité} qu'on aime faire en général, on utilise le gérondif (la base verbale en \cle{-ing}).

\begin{propriete}[Prefer + V-ing]
\textbf{Forme :} sujet + prefer + verbe en -ing.\\
\textbf{Emploi :} préférence générale pour une activité, une habitude de goût.
\end{propriete}

\begin{exemplebox}[Prefer + V-ing]
\eng{I prefer trading.} « Je préfère faire du commerce. »\\
\eng{She prefers selling at the market.} « Elle préfère vendre au marché. »\\
\eng{Many young people in Bamenda prefer working for themselves.} « Beaucoup de jeunes à Bamenda préfèrent travailler à leur compte. »\\
\eng{We preferred travelling by night bus.} « Nous préférions voyager en bus de nuit. »
\end{exemplebox}

\courssub{Structure 3 : prefer + to-infinitive}

On peut aussi utiliser l'infinitif avec \cle{to} après \eng{prefer}. Le sens est très proche du gérondif ; la nuance est subtile : l'infinitif évoque parfois un choix un peu plus précis ou délibéré, mais dans la pratique les deux formes sont souvent interchangeables.

\begin{propriete}[Prefer to + base verbale]
\textbf{Forme :} sujet + prefer + to + base verbale.\\
\textbf{Emploi :} préférence générale ou choix délibéré ; sens proche de prefer + V-ing.
\end{propriete}

\begin{exemplebox}[Prefer to + V]
\eng{I prefer to trade.} « Je préfère commercer. »\\
\eng{She prefers to sell in the morning.} « Elle préfère vendre le matin. »\\
\eng{They prefer to save their money in a credit union.} « Ils préfèrent épargner leur argent dans une coopérative de crédit. »
\end{exemplebox}

\begin{center}
\begin{tabularx}{\linewidth}{|X|X|X|}
\hline
\rowcolor{popblueL} Structure & Exemple & Sens \\
\hline
prefer + nom & I prefer coffee. & préférence pour une chose \\
\hline
prefer + V-ing & I prefer farming. & activité en général \\
\hline
prefer to + V & I prefer to farm. & sens proche, choix délibéré \\
\hline
\end{tabularx}
\end{center}

\courssub{Structure 4 : prefer X to Y}

C'est la construction la plus importante — et la plus dangereuse pour un francophone. Pour comparer deux éléments et dire lequel on préfère, l'anglais utilise la préposition \cle{to}, jamais \eng{than}.

\begin{propriete}[Règle de prefer X to Y]
\textbf{Forme :} sujet + prefer + X + \textbf{to} + Y.\\
X est l'élément préféré ; \textbf{to} introduit l'élément \cle{rejeté}.\\
\textbf{Condition obligatoire :} les deux éléments doivent être de \textbf{même nature grammaticale} : deux noms, ou deux gérondifs.\\
\eng{I prefer farming to fishing.} (deux gérondifs)\\
\eng{I prefer coffee to cocoa.} (deux noms)\\
Jamais \eng{than} avec \eng{prefer} !
\end{propriete}

\begin{exemplebox}[Prefer X to Y — exemples traduits]
\eng{I prefer coffee to cocoa.} « Je préfère le café au cacao. »\\
\eng{She prefers selling at the market to working in an office.} « Elle préfère vendre au marché plutôt que travailler dans un bureau. »\\
\eng{My uncle prefers teaching to trading.} « Mon oncle préfère l'enseignement au commerce. »\\
\eng{We prefer travelling by train to travelling by road.} « Nous préférons voyager en train plutôt que par la route. »\\
\eng{He preferred a government job to a private business.} « Il préférait un emploi public à une entreprise privée. »
\end{exemplebox}

\begin{attention}[Erreur fréquente des francophones]
En français on dit « je préfère le thé \textbf{au} café » ou « plutôt \textbf{que} le café ». Par calque, beaucoup d'élèves écrivent :\\
\eng{I prefer tea than coffee.} — FAUX.\\
La seule forme correcte est : \eng{I prefer tea to coffee.}\\
Retenez : \textbf{prefer ... to ...}, comme dans \eng{I prefer Cameroon to any other country.}
\end{attention}

\courssub{Structure 5 : prefer to do X rather than (do) Y}

Quand on veut comparer deux \cle{actions} avec des infinitifs, l'anglais utilise une construction spéciale : \eng{prefer to do X rather than (do) Y}. Ici, c'est bien \eng{rather than} qui apparaît — mais jamais \eng{than} seul juste après \eng{prefer}.

\begin{propriete}[Prefer to V ... rather than (V)]
\textbf{Forme :} sujet + prefer + \textbf{to} + base verbale (X) + \textbf{rather than} + base verbale (Y), sans to.\\
\eng{I prefer to work rather than stay idle.} « Je préfère travailler plutôt que de rester oisif. »\\
Le second verbe est à la base verbale, sans \eng{to} (le \eng{to} est sous-entendu).
\end{propriete}

\begin{exemplebox}[Prefer to ... rather than ...]
\eng{I prefer to walk rather than take a bendskin.} « Je préfère marcher plutôt que de prendre une moto-taxi. »\\
\eng{She prefers to save rather than spend all her salary.} « Elle préfère épargner plutôt que de dépenser tout son salaire. »\\
\eng{They prefer to grow cocoa rather than buy it from other farmers.} « Ils préfèrent cultiver le cacao plutôt que de l'acheter à d'autres planteurs. »
\end{exemplebox}

\courssub{Tableau récapitulatif des cinq structures}

\begin{popfigure}
\begin{tikzpicture}[
  box/.style={draw=popblue, fill=popblueL, rounded corners=2pt, align=center, font=\small, inner sep=4pt},
  ex/.style={draw=poporange, fill=poporangeL, rounded corners=2pt, align=center, font=\small, inner sep=4pt}
]
\node[box] (s1) at (0,0) {\textbf{prefer + N}};
\node[ex, right=0.4cm of s1] {I prefer tea.};
\node[box] (s2) at (0,-1.1) {\textbf{prefer + V-ing}};
\node[ex, right=0.4cm of s2] {I prefer trading.};
\node[box] (s3) at (0,-2.2) {\textbf{prefer to + V}};
\node[ex, right=0.4cm of s3] {I prefer to trade.};
\node[box] (s4) at (0,-3.3) {\textbf{prefer X to Y}};
\node[ex, right=0.4cm of s4] {I prefer farming to fishing.};
\node[box] (s5) at (0,-4.4) {\textbf{prefer to V rather than V}};
\node[ex, right=0.4cm of s5] {I prefer to work rather than stay idle.};
\end{tikzpicture}
\legende{Les cinq constructions du verbe prefer}
\end{popfigure}

\begin{center}
\begin{tabularx}{\linewidth}{|X|X|X|}
\hline
\rowcolor{popblueL} Français & English & Remarque \\
\hline
Je préfère le thé. & I prefer tea. & pas d'article en anglais \\
\hline
Je préfère le thé au café. & I prefer tea to coffee. & « to », jamais « than » \\
\hline
Elle préfère vendre au marché. & She prefers selling at the market. & gérondif après prefer \\
\hline
Il préfère travailler plutôt que rester oisif. & He prefers to work rather than stay idle. & rather than + base verbale \\
\hline
Nous préférions voyager de nuit. & We preferred travelling at night. & r redoublé au passé \\
\hline
\end{tabularx}
\end{center}

\courssub{Préférence générale ou choix ponctuel : la nuance essentielle}

\begin{attention}[Prefer + -ing ou would prefer to + V ?]
Ne confondez pas la \cle{préférence générale} et le \cle{choix ponctuel} :\\
\eng{I prefer walking.} = « Je préfère marcher » (en général, c'est mon habitude).\\
\eng{I'd prefer to walk today.} = « Je préférerais marcher aujourd'hui » (ce coup-ci, dans cette situation précise).\\
Le conditionnel \eng{would prefer} suivi de l'infinitif exprime un choix particulier, immédiat ; \eng{prefer} au présent suivi du gérondif exprime un goût durable. Nous étudierons \eng{would rather} et \eng{would prefer} en détail dans la section suivante.
\end{attention}

\courssub{Exercice résolu}

\begin{exoresolu}[Mettez les verbes à la bonne forme]
Énoncé — Complétez avec la forme correcte du verbe entre parenthèses :\\
1. I prefer (farm) \dots\ to (fish) \dots.\\
2. She prefers (sell) \dots\ groundnuts at the market.\\
3. They prefer (travel) \dots\ by train rather than (go) \dots\ by road.\\
4. Last year, my brother (prefer) \dots\ teaching to trading.
\tcblower
Solution détaillée :\\
1. \eng{I prefer farming to fishing.} — Structure 4 : prefer X to Y exige deux éléments de même nature ; ici deux gérondifs : \eng{farming} et \eng{fishing}.\\
2. \eng{She prefers selling groundnuts at the market.} — Structure 2 : gérondif après prefer pour une activité habituelle ; attention au -s de la 3\textsuperscript{e} personne.\\
3. \eng{They prefer to travel by train rather than go by road.} — Structure 5 : \eng{prefer to + V} puis \eng{rather than} + base verbale sans \eng{to}.\\
4. \eng{Last year, my brother preferred teaching to trading.} — Passé : redoublement du r (\eng{preferred}) ; la structure X to Y reste inchangée au passé.
\end{exoresolu}

\begin{exercice}[À vous]
Traduisez en anglais :\\
1. Je préfère le cacao au café.\\
2. Elle préfère travailler à son compte.\\
3. Nous préférons épargner plutôt que dépenser.\\
4. Mon père préférait la vie du village à la vie en ville.
\end{exercice}

\begin{corrige}[Exercice « À vous »]
1. \eng{I prefer cocoa to coffee.}\\
2. \eng{She prefers working for herself.} (ou \eng{being self-employed})\\
3. \eng{We prefer to save rather than spend.} (ou \eng{We prefer saving to spending.})\\
4. \eng{My father preferred village life to city life.}
\end{corrige}

\begin{aretenir}[L'essentiel sur prefer]
\begin{itemize}
\item \eng{prefer} + nom, + V-ing ou + to-V : préférence générale.
\item Comparaison de deux éléments : \eng{prefer X \textbf{to} Y} — jamais \eng{than}.
\item X et Y doivent être de même nature : deux noms ou deux gérondifs.
\item Comparaison d'actions à l'infinitif : \eng{prefer to do X rather than (do) Y}.
\item Orthographe : \eng{preferred}, \eng{preferring} (r redoublé).
\item Choix ponctuel : \eng{would prefer to + V} — à ne pas confondre avec la préférence générale.
\end{itemize}
\end{aretenir}

\coursec{Would rather et would prefer}

Dans la section précédente, nous avons étudié le verbe \eng{prefer} et ses constructions : \eng{prefer + -ing}, \eng{prefer + to-infinitive}, \eng{prefer A to B}. Nous avons vu que \eng{prefer} exprime une préférence \cle{générale}, une habitude de goût : \eng{I prefer tea to coffee.} « Je préfère le thé au café. » Mais que se passe-t-il quand le choix est \cle{ponctuel}, lié à une situation précise, ici et maintenant ? Imaginez un jeune diplômé de Yaoundé à qui l'on propose deux emplois : l'un dans une banque à Douala, l'autre dans une coopérative de cacao à Buea. Il ne parle pas de ses goûts en général : il doit décider \emph{maintenant}. C'est exactement le rôle de \eng{would rather} et \eng{would prefer}. C'est ce que nous allons découvrir dans cette section.

\textbf{Observation : un dialogue au bureau}\par

\begin{textebox}[At the job agency in Yaounde]
\textbf{Agent :} “Good morning, Mr Etoundi. We have two offers for you. The first one is a position as an accountant in a bank in Douala. The second one is a job as a manager in a cocoa cooperative near Buea. Would you rather work in a big city or in a smaller town?”

\textbf{Etoundi :} “Well, I'd rather stay near my family, so I'd prefer to take the job in Buea. Life in Douala is very expensive.”

\textbf{Agent :} “I understand. But the salary in Douala is higher.”

\textbf{Etoundi :} “I know, but I'd rather earn less money and live comfortably. Besides, I'd rather not spend two hours in traffic jams every morning.”

\textbf{Agent :} “What about your wife? Does she agree?”

\textbf{Etoundi :} “Honestly, I'd rather she found a job there too before we move. She is a nurse, so it should be possible.”

\textbf{Agent :} “Fine. Would you prefer to start in January or in March?”

\textbf{Etoundi :} “I'd prefer to start in January rather than wait until March. I need an income quickly.”
\end{textebox}

\textbf{Glossaire :} \emph{offer} (n.) : offre ; \emph{accountant} (n.) : comptable ; \emph{cooperative} (n.) : coopérative ; \emph{salary} (n.) : salaire ; \emph{to earn} (v.) : gagner (de l'argent) ; \emph{traffic jam} (n.) : embouteillage ; \emph{to move} (v.) : déménager ; \emph{income} (n.) : revenu.

\textbf{Questions de compréhension :}
\begin{enumerate}
\item What are the two job offers?
\item Why does Etoundi prefer the job in Buea?
\item What does he want his wife to do before they move?
\item When does he prefer to start, and why?
\end{enumerate}

Relevons maintenant les structures du dialogue : \eng{Would you rather work...?}, \eng{I'd rather stay...}, \eng{I'd rather not spend...}, \eng{I'd prefer to take...}, \eng{I'd rather she found...}, \eng{I'd prefer to start in January rather than wait...}. Deux familles se dessinent : \cle{would rather} et \cle{would prefer}. Observons-les en détail.

\textbf{Would rather + base verbale}\par

\begin{propriete}[Règle de would rather]
\eng{Would rather} est suivi \textbf{directement de la base verbale}, \textbf{sans} \eng{to} :
\begin{center}
\textbf{sujet + would rather + base verbale}
\end{center}
\begin{itemize}
\item \eng{I'd rather stay.} « Je préférerais rester. »
\item \eng{He'd rather work in a cooperative.} « Il préférerait travailler dans une coopérative. »
\end{itemize}
Points essentiels :
\begin{itemize}
\item \eng{would} se contracte en \textbf{'d} : \eng{I'd rather}, \eng{she'd rather}, \eng{we'd rather}.
\item \eng{would rather} est \textbf{invariable} : jamais de \eng{-s} à la troisième personne (\eng{She'd rather stay}, et non \emph{She'd rathers}).
\item La \textbf{négation} se fait avec \eng{not} placé \textbf{avant le verbe} : \eng{would rather \textbf{not} + V}.
\item L'\textbf{interrogation} s'obtient par inversion de \eng{would} : \eng{Would you rather...?}
\end{itemize}
\end{propriete}

\begin{exemplebox}[Exemples traduits]
\begin{itemize}
\item \eng{I'd rather invest in cocoa than in cotton.} « Je préférerais investir dans le cacao plutôt que dans le coton. »
\item \eng{I'd rather not discuss salaries during the interview.} « Je préférerais ne pas discuter des salaires pendant l'entretien. »
\item \eng{Would you rather be an employer or an employee?} « Préférerais-tu être employeur ou employé ? »
\item \eng{They'd rather trade than borrow money.} « Ils préféreraient commercer plutôt qu'emprunter de l'argent. »
\item \eng{We'd rather not take a bendskin in this rain.} « Nous préférerions ne pas prendre de moto-taxi sous cette pluie. »
\end{itemize}
\end{exemplebox}

\begin{attention}[Pas de « to » après would rather !]
L'erreur la plus fréquente des francophones est d'ajouter \eng{to} après \eng{would rather}, par analogie avec \eng{I'd prefer to go}.
\begin{itemize}
\item FAUX : \emph{I'd rather to go home.}
\item CORRECT : \eng{I'd rather go home.}
\end{itemize}
Astuce : \eng{rather} « avale » le verbe directement ; \eng{prefer}, lui, exige \eng{to}. Retenez : \textbf{rather = base verbale ; prefer = to + base verbale}.
\end{attention}

\textbf{Would rather + V ... than + V}\par

Pour opposer deux actions, on utilise \eng{than} suivi lui aussi de la \textbf{base verbale} (sans \eng{to}) :

\begin{center}
\textbf{sujet + would rather + V ... than + V}
\end{center}

\begin{exemplebox}[Exemples traduits]
\begin{itemize}
\item \eng{I'd rather trade than borrow money.} « Je préférerais commercer plutôt qu'emprunter de l'argent. »
\item \eng{She'd rather sell vegetables at the market than work in an office.} « Elle préférerait vendre des légumes au marché plutôt que de travailler dans un bureau. »
\item \eng{Would you rather work in Douala or in Bamenda?} « Préférerais-tu travailler à Douala ou à Bamenda ? » (à la question, on peut aussi utiliser \eng{or} pour proposer un choix)
\end{itemize}
\end{exemplebox}

\begin{attention}[Parallélisme des formes]
Les deux verbes autour de \eng{than} doivent avoir la \textbf{même forme} : la base verbale des deux côtés.
\begin{itemize}
\item FAUX : \emph{I'd rather stay than to leave.}
\item CORRECT : \eng{I'd rather stay than leave.}
\end{itemize}
\end{attention}

\textbf{Would prefer + to-infinitive}\par

\begin{propriete}[Règle de would prefer]
\eng{Would prefer} se construit avec la \textbf{to-infinitive} :
\begin{center}
\textbf{sujet + would prefer + to + base verbale}
\end{center}
Pour opposer deux actions, on ajoute \eng{rather than} suivi d'un \textbf{élément parallèle} (base verbale sans \eng{to}, gérondif, nom ou groupe prépositionnel) :
\begin{center}
\textbf{would prefer + to V ... rather than + élément parallèle}
\end{center}
\begin{itemize}
\item \eng{I'd prefer to start my own business.} « Je préférerais créer ma propre entreprise. »
\item \eng{I'd prefer to work in Buea rather than live in Douala.} « Je préférerais travailler à Buea plutôt que de vivre à Douala. »
\item \eng{I'd prefer to invest in cocoa rather than in cotton.} « Je préférerais investir dans le cacao plutôt que dans le coton. »
\item Négation : \eng{I'd prefer \textbf{not} to wait.} « Je préférerais ne pas attendre. » (le \eng{not} se place avant \eng{to})
\end{itemize}
\end{propriete}

\begin{exemplebox}[Exemples traduits]
\begin{itemize}
\item \eng{I'd prefer to be self-employed rather than depend on an employer.} « Je préférerais être à mon compte plutôt que de dépendre d'un employeur. »
\item \eng{Would you prefer to pay cash or to use mobile money?} « Préféreriez-vous payer en espèces ou par mobile money ? »
\item \eng{She'd prefer not to borrow from the bank.} « Elle préférerait ne pas emprunter à la banque. »
\end{itemize}
\end{exemplebox}

\textbf{Choix ponctuel ou préférence générale ?}\par

C'est la distinction capitale entre les structures de cette section et le \eng{prefer} de la section précédente :

\begin{itemize}
\item \eng{prefer + -ing} exprime une préférence \textbf{générale}, une habitude : \eng{I prefer walking to work.} « Je préfère aller au travail à pied » (en général, tous les jours).
\item \eng{would rather} et \eng{would prefer + to V} expriment un choix \textbf{ponctuel}, dans une situation précise : \eng{It's raining; I'd rather take a taxi today.} « Il pleut ; je préférerais prendre un taxi aujourd'hui » (cette fois-ci).
\end{itemize}

\begin{popfigure}
\begin{tikzpicture}[
node distance=0.9cm,
q/.style={rectangle, rounded corners, draw=popblue, fill=popblueL, align=center, text width=4.2cm, font=\small\bfseries},
yes/.style={rectangle, rounded corners, draw=popgreen, fill=popgreenL, align=center, text width=4.6cm, font=\small},
no/.style={rectangle, rounded corners, draw=poporange, fill=poporangeL, align=center, text width=4.6cm, font=\small},
lbl/.style={font=\small\itshape, text=popdark}
]
\node[q] (q) {Choix ponctuel, situation précise ?};
\node[yes, below left=1.1cm and 1.2cm of q, align=center] (y){OUI : \textbf{would rather + V}\\ \textbf{would prefer + to V}\\ \emph{I'd rather take a taxi.}};
\node[no, below right=1.1cm and 1.2cm of q, align=center] (n){NON (habitude, goût général) :\\ \textbf{prefer + -ing}\\ \emph{I prefer walking.}};
\draw[-{Stealth}, thick, popblue] (q.south) -- ++(0,-0.45) -| (y.north) node[lbl, pos=0.25, left]{oui};
\draw[-{Stealth}, thick, popblue] (q.south) -- ++(0,-0.45) -| (n.north) node[lbl, pos=0.25, right]{non};
\end{tikzpicture}
\legende{Organigramme de choix : would rather / would prefer (ponctuel) contre prefer + -ing (général)}
\end{popfigure}

\textbf{Would rather + sujet + prétérit : parler d'une autre personne}\par

\begin{propriete}[Structure soutenue : would rather + sujet + prétérit]
Quand la personne qui exprime la préférence \textbf{n'est pas} celle qui fait l'action, on change de construction :
\begin{center}
\textbf{sujet 1 + would rather + sujet 2 + prétérit}
\end{center}
\begin{itemize}
\item \eng{I'd rather you stayed here.} « Je préférerais que tu restes ici. »
\item \eng{She'd rather her son studied medicine.} « Elle préférerait que son fils fasse médecine. »
\item Négation : \eng{I'd rather you \textbf{didn't} smoke.} « Je préférerais que tu ne fumes pas. »
\end{itemize}
Le prétérit n'exprime pas le passé ici : c'est un \textbf{prétérit irréel} (comme après \eng{I wish}), qui marque l'éloignement entre le souhait et la réalité.
\end{propriete}

\begin{savaistu}[Une structure qui impressionne au Bac]
\eng{I'd rather you didn't smoke} utilise le prétérit après \eng{would rather} parce que le sujet change : ce n'est plus « je » qui agit, mais « tu ». Cette structure, proche du subjonctif français (« je préférerais que tu ne fumes pas »), est du registre soutenu et est très appréciée dans les compositions du Baccalauréat. Attention toutefois à ne pas dire \emph{I'd rather you to stay} : après le nouveau sujet, c'est le prétérit simple qui s'impose, jamais l'infinitif.
\end{savaistu}

\textbf{Tableaux récapitulatifs}\par

\begin{center}
\begin{tabular}{|l|l|l|}
\hline
\rowcolor{popblueL} \textbf{Structure} & \textbf{Exemple} & \textbf{Traduction} \\
\hline
would rather + V & I'd rather stay. & Je préférerais rester. \\
\hline
would rather + V + than + V & I'd rather trade than borrow. & Je préférerais commercer plutôt qu'emprunter. \\
\hline
would rather not + V & I'd rather not argue. & Je préférerais ne pas discuter. \\
\hline
Would + sujet + rather + V ? & Would you rather wait? & Préférerais-tu attendre ? \\
\hline
would prefer + to V & I'd prefer to leave now. & Je préférerais partir maintenant. \\
\hline
would prefer + to V + rather than + V & I'd prefer to save rather than spend. & Je préférerais épargner plutôt que dépenser. \\
\hline
would rather + sujet + prétérit & I'd rather you came. & Je préférerais que tu viennes. \\
\hline
\end{tabular}
\end{center}

\begin{center}
\begin{tabularx}{\linewidth}{|X|X|X|}
\hline
\rowcolor{popblueL} \textbf{Français} & \textbf{English} & \textbf{Remarque} \\
\hline
Je préférerais partir. & I'd rather go. / I'd prefer to go. & Pas de \eng{to} après \eng{rather}. \\
\hline
Je préférerais que tu partes. & I'd rather you left. & Prétérit après le nouveau sujet. \\
\hline
Je préférerais ne pas payer. & I'd rather not pay. & \eng{not} avant le verbe. \\
\hline
En général, je préfère le thé. & I prefer tea. & \eng{prefer} simple, habitude. \\
\hline
Ce soir, je préférerais du café. & Tonight, I'd rather have coffee. & Choix ponctuel. \\
\hline
\end{tabularx}
\end{center}

\textbf{Entraînement}\par

\begin{exoresolu}[Complétez avec la forme correcte]
Complétez chaque phrase avec \eng{would rather} ou \eng{would prefer} et le verbe entre parenthèses à la forme correcte.

1. I .................. (stay) in Bamenda these holidays.
2. She .................. (not / discuss) her salary with strangers.
3. .................. you .................. (work) in a bank or in a shop?
4. We .................. (invest) in cocoa rather than in cotton.
5. I'd rather you .................. (not / smoke) in the office.
\tcblower
\textbf{Solution détaillée :}

1. \eng{I'd rather stay in Bamenda these holidays.} — Choix ponctuel ; après \eng{would rather}, base verbale \eng{stay} sans \eng{to}.

2. \eng{She'd rather not discuss her salary with strangers.} (ou \eng{She'd prefer not to discuss...}) — Négation : \eng{not} avant la base verbale avec \eng{rather}, avant \eng{to} avec \eng{prefer}.

3. \eng{Would you rather work in a bank or in a shop?} — Interrogation : inversion de \eng{would} ; base verbale \eng{work}.

4. \eng{We'd prefer to invest in cocoa rather than in cotton.} — Avec \eng{would prefer}, la to-infinitive \eng{to invest} est obligatoire ; structure parallèle \eng{rather than in cotton} (groupe prépositionnel). On aurait aussi pu dire \eng{We'd rather invest in cocoa than in cotton.}

5. \eng{I'd rather you didn't smoke in the office.} — Le sujet change (\eng{you}), donc prétérit irréel : \eng{didn't smoke}.
\end{exoresolu}

\begin{exercice}[À vous de jouer]
Traduisez en anglais :
\begin{enumerate}
\item Je préférerais créer ma propre entreprise.
\item Préférerais-tu travailler à Douala ou à Buea ?
\item Il préférerait ne pas emprunter d'argent à la banque.
\item Je préférerais que tu m'accompagnes au marché.
\item En général, je préfère épargner, mais ce mois-ci je préférerais acheter une moto.
\end{enumerate}
\end{exercice}

\begin{corrige}[Exercice « À vous de jouer »]
\begin{enumerate}
\item \eng{I'd prefer to start my own business.} (ou \eng{I'd rather start my own business.})
\item \eng{Would you rather work in Douala or in Buea?}
\item \eng{He'd rather not borrow money from the bank.} (ou \eng{He'd prefer not to borrow...})
\item \eng{I'd rather you came with me to the market.}
\item \eng{In general, I prefer saving money, but this month I'd rather buy a motorbike.}
\end{enumerate}
Notez la différence dans la phrase 5 : \eng{prefer saving} pour l'habitude, \eng{would rather buy} pour le choix ponctuel.
\end{corrige}

\begin{aretenir}[L'essentiel de la section]
\begin{itemize}
\item \eng{would rather + base verbale} (sans \eng{to}) ; négation : \eng{would rather not + V} ; question : \eng{Would you rather...?}
\item Opposition de deux actions : \eng{would rather V than V} ; \eng{would prefer to V rather than + élément parallèle}.
\item \eng{would prefer} exige la \textbf{to-infinitive} : \eng{I'd prefer to go.}
\item Choix \textbf{ponctuel} : \eng{would rather / would prefer} ; préférence \textbf{générale} : \eng{prefer + -ing}.
\item Sujet différent : \eng{would rather + sujet + prétérit} : \eng{I'd rather you stayed.}
\end{itemize}
\end{aretenir}

\coursec{Like... better than et autres structures}

Dans la section précédente, nous avons étudié \eng{would rather} et \eng{would prefer}, deux structures qui expriment un choix ponctuel, souvent lié à une situation précise : \eng{I would rather stay in Douala tonight.} Mais l'anglais possède d'autres outils pour exprimer des goûts généraux, des comparaisons entre deux activités, ou encore des préférences nuancées selon le registre de langue. C'est l'objet de cette section : la structure \eng{like X better than Y}, l'expression de l'équivalence avec \eng{as much as}, et un petit lexique d'expressions idiomatiques (\eng{be fond of}, \eng{be keen on}, \eng{have a preference for}, \eng{fancy}) qui vous permettront d'adapter votre langue au contexte, de la conversation familière à la lettre formelle.

\courssub{Observation : un dialogue au lycée}

\begin{textebox}[Career talk at Government High School, Bamenda]
\textbf{Guidance counsellor:} So, Ngong, have you thought about your future career?

\textbf{Ngong:} Well, sir, I like trading better than farming. My father is a cocoa farmer, and I respect his work, but business attracts me more.

\textbf{Counsellor:} Interesting. And what about your sister? I heard she is keen on accounting.

\textbf{Ngong:} That's right. She is fond of numbers. She says accounting is more rewarding than trading because it offers a stable salary.

\textbf{Counsellor:} Do you fancy working in a bank one day, then?

\textbf{Ngong:} Not really. I like the market as much as the office, but I would prefer to run my own shop. Honestly, I have a preference for independent work.

\textbf{Counsellor:} I see. You like being your own boss better than taking orders!

\textbf{Ngong:} Exactly, sir. Freedom is better than routine, at least for me.
\end{textebox}

\textbf{Glossaire :} \emph{guidance counsellor} (n.) conseiller d'orientation ; \emph{rewarding} (adj.) gratifiant, enrichissant ; \emph{stable} (adj.) stable ; \emph{to run a shop} (v.) gérer une boutique ; \emph{to be your own boss} (expr.) être son propre patron ; \emph{to take orders} (expr.) recevoir des ordres ; \emph{at least} (loc. adv.) du moins.

\textbf{Questions de compréhension :}
\begin{enumerate}
\item Which activity does Ngong prefer: trading or farming?
\item Why does his sister prefer accounting to trading?
\item What does Ngong mean when he says he likes the market \emph{as much as} the office?
\item Find two expressions in the dialogue that mean “to like very much”.
\end{enumerate}

\textbf{Observation grammaticale.} Relevez les structures du dialogue : \eng{I like trading better than farming} ; \eng{accounting is more rewarding than trading} ; \eng{I like the market as much as the office} ; \eng{she is keen on accounting} ; \eng{do you fancy working in a bank?} ; \eng{I have a preference for independent work}. Toutes expriment une préférence, mais elles n'appartiennent pas au même registre et n'ont pas la même construction. Analysons-les une par une.

\courssub{La structure like X better than Y}

\begin{propriete}[Règle : like ... better than]
Pour comparer deux choses ou deux activités et dire que l'on préfère la première, l'anglais utilise :

\begin{center}
\textbf{sujet + like(s) + X + better than + Y}
\end{center}

\begin{itemize}
\item \cle{better} est le comparatif de \eng{well} (et de \eng{good}) : il signifie ici « mieux / plus ».
\item \cle{than} est obligatoire : c'est la particule de comparaison.
\item X et Y doivent être \textbf{parallèles} : deux noms, deux pronoms ou deux gérondifs (\eng{-ing}).
\end{itemize}

\eng{I like teaching better than banking.} « J'aime l'enseignement plus que la banque. »

\eng{She likes working in Buea better than working in Douala.} « Elle aime travailler à Buea plus qu'à Douala. »
\end{propriete}

\begin{exemplebox}[Exemples traduits]
\eng{I like trading better than farming.} « J'aime le commerce plus que l'agriculture. »

\eng{We like the new market better than the old one.} « Nous aimons le nouveau marché plus que l'ancien. »

\eng{My uncle likes driving better than taking a bendskin.} « Mon oncle aime conduire plus que prendre une moto-taxi. »

\eng{Do you like tea better than coffee?} « Aimes-tu le thé plus que le café ? »
\end{exemplebox}

\begin{attention}[Ne mélangez pas les structures !]
\eng{like ... better than} garde \cle{than} (comparatif), contrairement à \eng{prefer ... to} qui exige \cle{to}. Ne mélangez jamais les deux :

\begin{itemize}
\item Correct : \eng{I like trading better than farming.}
\item Correct : \eng{I prefer trading to farming.}
\item FAUX : \emph{I like trading better to farming.} (mélange des deux structures)
\item FAUX : \emph{I prefer trading than farming.} (\eng{prefer} ne prend jamais \eng{than} seul)
\end{itemize}

Autre erreur fréquente : oublier \cle{better}. \emph{I like trading than farming} est incorrect : sans \eng{better}, il n'y a pas de comparaison.
\end{attention}

\courssub{La variante avec les comparatifs d'adjectifs}

Une autre façon naturelle de comparer deux options consiste à utiliser un \textbf{adjectif au comparatif} avec \eng{than} :

\begin{propriete}[Règle : X + be + comparatif + than + Y]
\begin{center}
\textbf{X + be + adjectif comparatif + than + Y}
\end{center}

\begin{itemize}
\item Adjectifs courts : \eng{-er} + \eng{than} : \eng{Trading is safer than borrowing.} « Le commerce est plus sûr que l'emprunt. »
\item Adjectifs longs : \eng{more} + adjectif + \eng{than} : \eng{Teaching is more rewarding than banking.} « L'enseignement est plus gratifiant que la banque. »
\item Comparatifs irréguliers : \eng{good} $\rightarrow$ \cle{better} ; \eng{bad} $\rightarrow$ \cle{worse}.
\end{itemize}

\eng{Trading is better than borrowing.} « Commercer vaut mieux qu'emprunter. »

\eng{Working in town is worse than working in the village for my health.} « Travailler en ville est pire que travailler au village pour ma santé. »
\end{propriete}

Notez que le sujet X est souvent un \textbf{gérondif} (\eng{teaching}, \eng{trading}) : en anglais, le verbe en \eng{-ing} joue le rôle de nom, là où le français utilise l'infinitif (« commercer », « enseigner »).

\courssub{Exprimer l'équivalence : like X as much as Y}

\begin{propriete}[Règle : as much as]
Pour dire que deux choses vous plaisent \textbf{également}, utilisez :

\begin{center}
\textbf{sujet + like(s) + X + as much as + Y}
\end{center}

\eng{I like cocoa farming as much as coffee farming.} « J'aime la culture du cacao autant que celle du café. »

\eng{She likes Yaoundé as much as Bamenda.} « Elle aime Yaoundé autant que Bamenda. »

La négation inverse le sens : \eng{I don't like borrowing as much as saving.} « Je n'aime pas emprunter autant qu'épargner » (= je préfère épargner).
\end{propriete}

\begin{attention}[as ... as : ne pas confondre avec than]
La structure d'égalité exige \cle{as ... as}, jamais \eng{than} :

\begin{itemize}
\item Correct : \eng{I like coffee as much as tea.}
\item FAUX : \emph{I like coffee as much than tea.}
\end{itemize}

Retenez le contraste : \cle{as ... as} = égalité (« autant ... que ») ; \cle{more/-er ... than} = supériorité (« plus ... que »).
\end{attention}

\courssub{Expressions idiomatiques : be fond of, be keen on, have a preference for, fancy}

L'anglais dispose d'un petit lexique très vivant pour dire « aimer », chacun avec sa nuance et sa construction :

\begin{definition}[be fond of et be keen on]
\cle{be fond of} + nom/gérondif : « aimer beaucoup, avoir de l'affection pour » (ton chaleureux, neutre).

\eng{She is fond of numbers.} « Elle aime beaucoup les chiffres. »

\cle{be keen on} + nom/gérondif : « être passionné de, être très motivé par » (très courant en Grande-Bretagne).

\eng{She is keen on accounting.} « Elle est passionnée de comptabilité. »

\eng{The students are keen on learning computer skills.} « Les élèves sont très motivés pour apprendre l'informatique. »

Attention : dans les deux cas, la préposition est \cle{of} ou \cle{on}, et le verbe qui suit est au gérondif.
\end{definition}

\begin{definition}[have a preference for (formel)]
\cle{have a preference for} + nom : « avoir une préférence pour ». Registre formel, typique de l'écrit administratif, des lettres de motivation et des questionnaires.

\eng{I have a preference for independent work.} « J'ai une préférence pour le travail indépendant. »

\eng{Candidates with a preference for evening classes should tick box B.} « Les candidats qui préfèrent les cours du soir doivent cocher la case B. »
\end{definition}

\begin{definition}[fancy (familier, britannique)]
\cle{fancy} est un verbe familier britannique signifiant « avoir envie de ». Il est suivi d'un nom ou d'un gérondif.

\eng{Do you fancy a cup of tea?} « Tu veux une tasse de thé ? » (prononciation : « fan-si »)

\eng{Do you fancy working in a bank?} « Ça te dit de travailler dans une banque ? »

À réserver aux conversations entre amis : ne l'utilisez jamais dans une lettre officielle ou à l'examen dans un texte formel.
\end{definition}

\courssub{Nuances de registre : choisir la bonne expression}

\begin{center}
\begin{tabularx}{\linewidth}{|X|X|X|}
\hline
\rowcolor{popblueL} Expression & Registre & Exemple \\
\hline
fancy + nom/-ing & familier (GB) & \eng{Do you fancy a drink?} \\
\hline
be keen on + -ing & courant, oral & \eng{I'm keen on trading.} \\
\hline
be fond of + nom/-ing & neutre, chaleureux & \eng{She is fond of teaching.} \\
\hline
like X better than Y & neutre & \eng{I like teaching better than banking.} \\
\hline
prefer X to Y & neutre, écrit et oral & \eng{I prefer teaching to banking.} \\
\hline
would prefer + to-V & neutre soutenu, choix ponctuel & \eng{I would prefer to stay.} \\
\hline
have a preference for + nom & formel, écrit & \eng{I have a preference for office work.} \\
\hline
\end{tabularx}
\end{center}

\begin{popfigure}
\begin{tikzpicture}[font=\small]
\draw[-{Stealth}, thick, popline] (0,0) -- (12,0);
\node[above, popmuted] at (0.4,0.1) {familier};
\node[above, popmuted] at (11.6,0.1) {formel};
\fill[popgreen] (1.5,0) circle (0.12);
\node[above=0.25cm, align=center, popgreen] at (1.5,0) {\textbf{fancy}\\ \eng{Do you fancy a tea?}};
\fill[popblue] (4.5,0) circle (0.12);
\node[below=0.25cm, align=center, popblue] at (4.5,0) {\textbf{be keen on / be fond of}\\ \eng{I'm keen on trading.}};
\fill[poporange] (7,0) circle (0.12);
\node[above=0.25cm, align=center, poporange] at (7,0) {\textbf{like better than / prefer}\\ \eng{I like trading better than farming.}};
\fill[poppurple] (10.2,0) circle (0.12);
\node[below=0.25cm, align=center, poppurple] at (10.2,0) {\textbf{have a preference for / would prefer}\\ \eng{I have a preference for office work.}};
\end{tikzpicture}
\legende{Frise des registres : du familier britannique (\eng{fancy}) au formel écrit (\eng{have a preference for}).}
\end{popfigure}

\courssub{Comparaison avec le français}

\begin{center}
\begin{tabularx}{\linewidth}{|X|X|X|}
\hline
\rowcolor{popblueL} Français & English & Remarque \\
\hline
J'aime le commerce plus que l'agriculture. & \eng{I like trading better than farming.} & \eng{better than}, pas \eng{more than} avec \eng{like} \\
\hline
Je préfère le commerce à l'agriculture. & \eng{I prefer trading to farming.} & \eng{prefer ... to}, jamais \eng{than} \\
\hline
J'aime le cacao autant que le café. & \eng{I like cocoa as much as coffee.} & égalité : \eng{as ... as} \\
\hline
Elle est passionnée de comptabilité. & \eng{She is keen on accounting.} & préposition \eng{on} + gérondif \\
\hline
Tu veux un thé ? (familier) & \eng{Do you fancy a tea?} & familier britannique uniquement \\
\hline
Enseigner est plus gratifiant. & \eng{Teaching is more rewarding.} & infinitif français = gérondif anglais \\
\hline
\end{tabularx}
\end{center}

\begin{attention}[Pièges fréquents des francophones]
\begin{enumerate}
\item Traduire « plus que » mot à mot : \emph{I like trading more than farming} se comprend, mais la forme idiomatique avec \eng{like} est \cle{better than}. Réservez \eng{more than} aux adjectifs longs : \eng{more interesting than}.
\item Oublier le gérondif après les prépositions : \emph{She is keen on accountancy} est correct (nom), mais avec un verbe il faut \eng{-ing} : \eng{She is keen on doing accounts.} (jamais \emph{keen on to do}).
\item Utiliser \eng{fancy} à l'écrit formel : c'est un marqueur familier britannique ; à l'examen, préférez \eng{would like} ou \eng{prefer}.
\item Confondre \eng{be fond of} et \eng{be found of} : attention à l'orthographe, \cle{fond} (« affection ») ne prend pas de \emph{u}.
\end{enumerate}
\end{attention}

\courssub{Exercice résolu et entraînement}

\begin{exoresolu}[Complétez avec la structure correcte]
Énoncé — Complétez chaque phrase avec \eng{better than}, \eng{as much as}, \eng{keen on}, \eng{fond of}, \eng{fancy} ou \eng{have a preference for}, selon le registre indiqué.

1. I like sewing \dots\ cooking. (comparaison)
2. My brother is \dots\ mechanics; he spends all his time in the garage. (passion)
3. (familier) \dots\ you \dots\ a walk to the market?
4. (formel, lettre de motivation) I \dots\ morning shifts.
5. She likes Douala \dots\ Yaoundé; both cities suit her. (égalité)

\tcblower
Solution détaillée :

1. \eng{I like sewing better than cooking.} — comparaison de deux gérondifs : \cle{better than} (jamais \eng{to} avec \eng{like}).
2. \eng{My brother is keen on mechanics\dots} — \cle{keen on} + nom exprime la passion ; \eng{fond of} serait possible mais plus faible.
3. \eng{Do you fancy a walk to the market?} — registre familier britannique : \cle{fancy} + nom.
4. \eng{I have a preference for morning shifts.} — registre formel de la lettre de motivation : \cle{have a preference for} + nom.
5. \eng{She likes Douala as much as Yaoundé.} — égalité : \cle{as much as} ; le contexte « both cities suit her » confirme qu'il n'y a pas de préférence.
\end{exoresolu}

\begin{exercice}[À vous de jouer]
1. Traduisez en anglais :
\begin{enumerate}
\item J'aime vendre au marché plus que travailler au bureau.
\item Elle aime la comptabilité autant que le commerce.
\item Tu as envie d'un jus de mangue ? (familier)
\item Il est passionné d'agriculture.
\end{enumerate}
2. Corrigez les erreurs : \emph{I prefer farming than trading.} / \emph{She is keen on to sew.} / \emph{I like coffee as much than tea.}
3. Rédigez trois phrases sur vos propres goûts professionnels en utilisant \eng{like ... better than}, \eng{as much as} et \eng{have a preference for}.
\end{exercice}

\begin{corrige}[Exercice « À vous de jouer »]
1. a) \eng{I like selling at the market better than working in an office.} b) \eng{She likes accounting as much as trading.} c) \eng{Do you fancy a mango juice?} d) \eng{He is keen on farming.}

2. \eng{I prefer farming to trading.} (\eng{prefer ... to}) ; \eng{She is keen on sewing.} (gérondif après \eng{on}) ; \eng{I like coffee as much as tea.} (\eng{as ... as} pour l'égalité).

3. Modèle : \eng{I like teaching better than banking. I like English as much as economics. For my future career, I have a preference for a job in Buea.}
\end{corrige}

\begin{aretenir}[L'essentiel de la section]
\begin{itemize}
\item Comparaison de goûts : \cle{like X better than Y} — avec \eng{than}, jamais \eng{to}.
\item Avec un adjectif : \cle{X is more rewarding / better / worse than Y}.
\item Équivalence : \cle{like X as much as Y} (« autant ... que »).
\item Lexique nuancé : \eng{be fond of} (affection), \eng{be keen on} (passion), \eng{have a preference for} (formel), \eng{fancy} (familier britannique).
\item Après \eng{of} et \eng{on}, le verbe est au gérondif : \eng{keen on sewing}.
\end{itemize}
\end{aretenir}

\coursec{Comparaison français/anglais et pièges}

Dans les sections précédentes, nous avons étudié \eng{prefer}, \eng{would rather}, \eng{would prefer} et \eng{like... better than}. Il est temps maintenant de confronter systématiquement les deux langues. C'est ici que tout se joue pour un élève francophone : le français et l'anglais expriment tous deux la préférence, mais \textbf{jamais avec les mêmes mots-outils}. Traduire mot à mot « préférer... à », « plutôt que » ou « j'aime mieux » conduit presque toujours à une faute. Cette section vous donne les correspondances exactes, puis un tableau des erreurs-types à éliminer définitivement.

\courssub{Observation : deux langues, deux logiques}

Lisez attentivement ces paires de phrases. Soulignez mentalement ce qui change entre le français et l'anglais.

\begin{exemplebox}[Observer avant de généraliser]
\begin{enumerate}
\item \eng{I prefer tea to coffee.} « Je préfère le thé au café. »
\item \eng{I'd rather stay home tonight.} « Je préférerais rester à la maison ce soir. »
\item \eng{She likes rice better than yams.} « Elle aime mieux le riz que les ignames. »
\item \eng{I'd rather you worked here.} « Je préférerais que tu travailles ici. »
\item \eng{He preferred trading to teaching.} « Il préférait commercer à enseigner. »
\end{enumerate}
\end{exemplebox}

Que remarque-t-on ? Le français utilise la préposition \textbf{« à »} (préférer X \emph{à} Y), le mot \textbf{« que »} (plutôt \emph{que}) et l'adverbe \textbf{« mieux »} (aimer \emph{mieux}). L'anglais, lui, utilise trois outils différents : \cle{to}, \cle{than} et \cle{better}. Aucun ne se traduit littéralement. Retenir ces correspondances, c'est gagner la moitié de la bataille.

\begin{propriete}[Correspondances fondamentales français / anglais]
\begin{center}
\begin{tabularx}{\linewidth}{|X|X|X|}\hline
\rowcolor{popblueL} \textbf{Français} & \textbf{English} & \textbf{Remarque} \\ \hline
préférer X \textbf{à} Y & prefer X \textbf{to} Y & jamais \eng{than} après \eng{prefer} \\ \hline
je préférerais X \textbf{à} Y & would prefer X \textbf{to} Y & \eng{would prefer} suit la même règle que \eng{prefer} \\ \hline
\textbf{plutôt que} + infinitif & \textbf{rather than} + base verbale & \eng{rather than} joint deux actions (souvent après \eng{prefer to}) \\ \hline
aimer \textbf{mieux} & like... \textbf{better} & jamais \eng{like more} pour une préférence \\ \hline
je préférerais + infinitif & I'\textbf{d rather} + base verbale & base verbale sans \eng{to} \\ \hline
je préférerais \textbf{que} tu + subjonctif & I'd rather you + \textbf{prétérit} & pas de subjonctif en anglais ; prétérit « irréel » \\ \hline
\end{tabularx}
\end{center}
\end{propriete}

\courssub{« Préférer X à Y » : prefer X to Y, jamais « than »}

C'est l'erreur numéro un des candidats francophones. En français, on dit « je préfère le thé \textbf{au} café » ; la tentation est grande de traduire « à » par \eng{than}, parce que la phrase ressemble à une comparaison. Or \eng{prefer} contient déjà l'idée de comparaison (il vient du latin \emph{praeferre}, « porter devant ») : il n'a donc pas besoin de \eng{than}.

\begin{propriete}[La préposition après prefer et would prefer]
Après \cle{prefer} et \cle{would prefer}, l'élément rejeté est introduit par la préposition \cle{to} (et non par \eng{than}) :
\begin{center}
\eng{prefer / would prefer + nom / -ing + to + nom / -ing}
\end{center}
\eng{I prefer fufu to rice.} « Je préfère le fufu au riz. »\\
\eng{She prefers walking to taking a bendskin.} « Elle préfère marcher à prendre un bendskin. »\\
\eng{We would prefer cocoa to coffee.} « Nous préférerions le cacao au café. »\\
Exception : quand \eng{prefer} est suivi d'un \textbf{infinitif complet}, la seconde action est introduite par \eng{rather than} : \eng{I prefer to walk rather than take a taxi.} « Je préfère marcher plutôt que de prendre un taxi. »
\end{propriete}

\courssub{« J'aime mieux » : I prefer ou I like... better — jamais « I like more »}

Le français « j'aime mieux » possède deux traductions naturelles en anglais :
\begin{itemize}
\item \eng{I prefer...} : registre neutre, toujours correct ;
\item \eng{I like... better} : registre plus oral, très fréquent.
\end{itemize}
En revanche, \eng{I like more} est une \textbf{faute grave} dans ce contexte. \eng{More} exprime la \emph{quantité}, pas la \emph{préférence} : \eng{I like more rice} signifie « j'aime davantage de riz » (une plus grande quantité), ce qui n'est pas du tout « j'aime mieux le riz ».

\begin{exemplebox}[« Mieux » = better]
\eng{He likes teaching better than trading.} « Il aime mieux enseigner que commercer. »\\
\eng{Do you like Douala better than Yaoundé?} « Aimes-tu mieux Douala que Yaoundé ? »\\
\eng{I like this market better.} « J'aime mieux ce marché. »
\end{exemplebox}

\courssub{Pas de subjonctif en anglais : « I'd rather you stayed »}

Le français exige le subjonctif après « je préfère que » : « je préfère que tu \emph{restes} ». L'anglais n'a pas de subjonctif dans ce cas ; il utilise le \textbf{prétérit} après \eng{would rather} quand le sujet de la seconde action est différent :

\begin{propriete}[Would rather + sujet différent]
\begin{center}
\eng{would rather + sujet + prétérit}
\end{center}
\eng{I'd rather you stayed.} « Je préférerais que tu restes. »\\
\eng{I'd rather you worked here.} « Je préférerais que tu travailles ici. »\\
\eng{She'd rather we came early.} « Elle préférerait que nous venions tôt. »\\
Le prétérit n'exprime pas le passé ici : c'est un prétérit « irréel », comme dans \eng{I wish I knew}. Ne traduisez jamais « que tu restes » par \eng{that you stay} ou \eng{that you will stay}.
\end{propriete}

\begin{exemplebox}[Exemples traduits supplémentaires]
\eng{I'd rather you didn't sell the shop.} « Je préférerais que tu ne vendes pas la boutique. »\\
\eng{We'd rather the meeting started at nine.} « Nous préférerions que la réunion commence à neuf heures. »\\
\eng{He'd rather his son became a teacher.} « Il préférerait que son fils devienne enseignant. »
\end{exemplebox}

\courssub{Orthographe : le redoublement du r de prefer}

Le verbe \eng{prefer} se prononce « pri-FEUR » : l'accent tombe sur la \textbf{dernière syllabe}. Or la règle d'orthographe anglaise veut qu'une consonne finale précédée d'une voyelle unique se \textbf{redouble} quand l'accent porte sur la dernière syllabe.

\begin{propriete}[Prefer, preferred, preferring]
\begin{center}
\begin{tabular}{|l|l|l|}\hline
\rowcolor{popblueL} \textbf{Forme} & \textbf{Orthographe} & \textbf{Exemple} \\ \hline
base verbale & prefer & \eng{I prefer tea.} \\ \hline
prétérit / participe passé & prefe\textbf{rr}ed & \eng{He preferred farming.} \\ \hline
forme en -ing & prefe\textbf{rr}ing & \eng{Preferring safety, she stayed.} \\ \hline
nom & prefe\textbf{r}ence (un seul r) & \eng{My preference is tea.} \\ \hline
\end{tabular}
\end{center}
Attention au piège : l'accent de \eng{preference} (« PRÉ-fe-rense ») se déplace sur la première syllabe, donc \textbf{pas de redoublement}. Comparez : \eng{offer → offered} (accent sur la première syllabe, un seul r) mais \eng{prefer → preferred} (accent final, deux r).
\end{propriete}

\courssub{Tableau des erreurs-types des francophones}

\begin{attention}[Les cinq fautes à éliminer]
\begin{itemize}
\item \eng{I prefer tea than coffee.} → \eng{I prefer tea \textbf{to} coffee.} (préposition \eng{to})
\item \eng{I'd rather to stay.} → \eng{I'd rather \textbf{stay}.} (base verbale sans \eng{to})
\item \eng{I like more rice.} → \eng{I like rice \textbf{better}.} / \eng{I \textbf{prefer} rice.} (\eng{more} = quantité)
\item \eng{I'd rather you \textbf{stay} here.} → \eng{I'd rather you \textbf{stayed} here.} (prétérit irréel)
\item \eng{He prefered trading.} → \eng{He prefe\textbf{rr}ed trading.} (redoublement du r)
\end{itemize}
\end{attention}

\begin{popfigure}
\begin{tikzpicture}
\node[draw=popdark, fill=popblueL, rounded corners, font=\bfseries, minimum width=6.2cm, minimum height=0.8cm] (h1) at (0,0) {Erreur fréquente};
\node[draw=popdark, fill=popgreenL, rounded corners, font=\bfseries, minimum width=6.2cm, minimum height=0.8cm] (h2) at (6.8,0) {Forme correcte};
\node[draw=popline, fill=popcream, rounded corners, minimum width=6.2cm, minimum height=0.7cm, font=\small] at (0,-0.9) {I prefer tea than coffee};
\node[draw=popline, fill=white, rounded corners, minimum width=6.2cm, minimum height=0.7cm, font=\small] at (6.8,-0.9) {I prefer tea to coffee};
\node[draw=popline, fill=popcream, rounded corners, minimum width=6.2cm, minimum height=0.7cm, font=\small] at (0,-1.8) {I'd rather to stay home};
\node[draw=popline, fill=white, rounded corners, minimum width=6.2cm, minimum height=0.7cm, font=\small] at (6.8,-1.8) {I'd rather stay home};
\node[draw=popline, fill=popcream, rounded corners, minimum width=6.2cm, minimum height=0.7cm, font=\small] at (0,-2.7) {I like more rice};
\node[draw=popline, fill=white, rounded corners, minimum width=6.2cm, minimum height=0.7cm, font=\small] at (6.8,-2.7) {I like rice better};
\node[draw=popline, fill=popcream, rounded corners, minimum width=6.2cm, minimum height=0.7cm, font=\small] at (0,-3.6) {I'd rather you stay};
\node[draw=popline, fill=white, rounded corners, minimum width=6.2cm, minimum height=0.7cm, font=\small] at (6.8,-3.6) {I'd rather you stayed};
\node[draw=popline, fill=popcream, rounded corners, minimum width=6.2cm, minimum height=0.7cm, font=\small] at (0,-4.5) {She prefered teaching};
\node[draw=popline, fill=white, rounded corners, minimum width=6.2cm, minimum height=0.7cm, font=\small] at (6.8,-4.5) {She preferred teaching};
\draw[-{Stealth}, thick, poporange] (1.6,-0.9) -- (3.7,-0.9);
\draw[-{Stealth}, thick, poporange] (1.6,-1.8) -- (3.7,-1.8);
\draw[-{Stealth}, thick, poporange] (1.6,-2.7) -- (3.7,-2.7);
\draw[-{Stealth}, thick, poporange] (1.6,-3.6) -- (3.7,-3.6);
\draw[-{Stealth}, thick, poporange] (1.6,-4.5) -- (3.7,-4.5);
\end{tikzpicture}
\legende{De l'erreur à la forme correcte : les cinq corrections à automatiser.}
\end{popfigure}

\courssub{Exercice résolu}

\begin{exoresolu}[Corrigez les fautes]
Énoncé — Chaque phrase contient une erreur typique de francophone. Corrigez-la et justifiez.
\begin{enumerate}
\item \eng{I prefer trading than teaching.}
\item \eng{We'd rather to take the bus to Bamenda.}
\item \eng{My father likes more farming.}
\item \eng{I'd rather you come tomorrow.}
\item \eng{She prefered working in Douala.}
\end{enumerate}
\tcblower
Solution détaillée :
\begin{enumerate}
\item \eng{I prefer trading \textbf{to} teaching.} — Après \eng{prefer}, l'élément rejeté s'introduit par \eng{to}, jamais \eng{than}.
\item \eng{We'd rather \textbf{take} the bus to Bamenda.} — \eng{Would rather} est suivi de la base verbale, sans \eng{to}.
\item \eng{My father likes farming \textbf{better}.} (ou \eng{My father \textbf{prefers} farming.}) — « Aimer mieux » se dit \eng{like... better} ; \eng{more} exprime la quantité.
\item \eng{I'd rather you \textbf{came} tomorrow.} — Sujet différent (\eng{you}) : l'anglais utilise le prétérit irréel \eng{came}, là où le français prend le subjonctif (« que tu viennes »).
\item \eng{She prefe\textbf{rr}ed working in Douala.} — L'accent de \eng{prefer} porte sur la dernière syllabe : le r se redouble au prétérit.
\end{enumerate}
\end{exoresolu}

\begin{savaistu}[Prefer et would rather au Bac]
Au baccalauréat, les exercices de transformation de phrases demandent très souvent de passer d'une structure à l'autre : \eng{I prefer tea to coffee} doit devenir \eng{I would rather drink tea than coffee}, ou l'inverse. Le correcteur vérifie trois choses : la préposition (\eng{to} après \eng{prefer/would prefer}, \eng{than} après \eng{would rather}), la forme verbale (base verbale après \eng{would rather}) et le sens global qui doit rester identique. Un candidat qui maîtrise les deux formes — et qui sait que \eng{prefer} prend \eng{to} tandis que \eng{would rather} prend \eng{than} — sécurise des points faciles que beaucoup perdent par simple traduction mot à mot du français.
\end{savaistu}

\begin{aretenir}[L'essentiel de la section]
\begin{itemize}
\item « Préférer X à Y » = \cle{prefer X to Y} / \cle{would prefer X to Y} ; « plutôt que » = \cle{rather than}.
\item « J'aime mieux » = \cle{I prefer} ou \cle{I like... better} ; jamais \eng{I like more}.
\item « Je préfère que tu + subjonctif » = \cle{I'd rather you + prétérit} : \eng{I'd rather you stayed.}
\item Orthographe : \eng{prefer → preferred, preferring} (r redoublé), mais \eng{preference} (un seul r).
\end{itemize}
\end{aretenir}

\begin{exercice}[À vous de jouer]
Traduisez en anglais : 1. « Je préfère le cacao au café. » 2. « Elle aime mieux vendre au marché que travailler au bureau. » 3. « Je préférerais que tu m'attendes ici. » 4. « Nous préférerions partir plutôt qu'attendre. » Puis transformez : \eng{I prefer teaching to trading.} → avec \eng{would rather... than}.
\end{exercice}

\begin{corrige}[Exercice « À vous de jouer »]
1. \eng{I prefer cocoa to coffee.} 2. \eng{She likes selling at the market better than working in an office.} 3. \eng{I'd rather you waited for me here.} 4. \eng{We'd rather leave than wait.} Transformation : \eng{I would rather teach than trade.}
\end{corrige}

\coursec{Entraînement et production guidée}

Après avoir analysé les pièges de la traduction et les différences structurelles entre le français et l'anglais, il est temps de mettre en pratique les cinq structures majeures de la préférence. Cette section propose un entraînement progressif : de la manipulation morphosyntaxique à la production écrite autonome, en passant par la correction d'erreurs types et la transformation de structures. Chaque exercice est suivi d'une correction commentée pour ancrer les réflexes grammaticaux.

\courssub{Exercices d'application}

\begin{exercice}[Exercice 1 : Complétion — \eng{to}, \eng{than} ou \eng{rather}]
Complétez chaque phrase avec le mot manquant : \eng{to}, \eng{than} ou \eng{rather}. Attention : certaines phrases n'exigent aucun de ces trois mots (écrivez \eng{Ø}).
\begin{enumerate}
    \item I prefer coffee \_\_\_\_\_\_\_\_ tea.
    \item She would \_\_\_\_\_\_\_\_ stay at home tonight.
    \item He prefers reading \_\_\_\_\_\_\_\_ watching television.
    \item They would prefer \_\_\_\_\_\_\_\_ go by train.
    \item I'd \_\_\_\_\_\_\_\_ not go out if it rains.
    \item My brother likes football better \_\_\_\_\_\_\_\_ basketball.
    \item We prefer \_\_\_\_\_\_\_\_ travel light.
    \item Would you rather stay \_\_\_\_\_\_\_\_ leave now?
    \item She prefers the city \_\_\_\_\_\_\_\_ the village.
    \item I would prefer coffee \_\_\_\_\_\_\_\_ tea.
\end{enumerate}
\end{exercice}

\begin{exercice}[Exercice 2 : Correction d'erreurs — Pièges francophones]
Chaque phrase ci-dessous contient une erreur typique d'élèves francophones. Soulignez l'erreur, réécrivez la phrase correcte et justifiez la correction en une ligne (règle de grammaire).
\begin{enumerate}
    \item I prefer tea \textbf{than} coffee.
    \item He would rather \textbf{to} stay here.
    \item She prefers \textbf{to read than to write}.
    \item I \textbf{had} rather go now.
    \item They prefer fish \textbf{than} meat.
    \item Would you prefer \textbf{going} or \textbf{to stay}?
    \item I \textbf{like better} tea \textbf{than} coffee.
    \item He'd rather \textbf{prefers} coffee.
\end{enumerate}
\end{exercice}

\begin{exoresolu}[Exemple traité : Item 2 de l'exercice 2]
\textbf{Phrase fautive :} \eng{He would rather to stay here.}
\textbf{Correction :} \eng{He would rather stay here.}
\textbf{Justification :} Après \cle{would rather} (forme contractée \eng{'d rather}), le verbe qui suit est obligatoirement à la \textbf{base verbale} (sans \eng{to}). C'est un modal semi-auxiliaire. La présence de \eng{to} est une interférence du français « préférerais \textbf{de} rester » ou de la structure \eng{would prefer to}.
\end{exoresolu}

\begin{methode}[Transformer \eng{prefer} en \eng{would rather} : mode d'emploi]
Suivez ces 4 étapes pour réécrire une phrase avec \eng{prefer} en utilisant \eng{would rather} sans changer le sens :
\begin{enumerate}
    \item \textbf{Repérez les deux éléments comparés} (noms, gérondifs ou infinitifs) après \eng{prefer} et après \eng{to}.
    \item \textbf{Gardez l'ordre} : le premier élément reste le préféré, le second reste le moins préféré.
    \item \textbf{Remplacez la structure} :
    \begin{itemize}
        \item \eng{prefer + Noun + to + Noun} $\rightarrow$ \eng{would rather + have + Noun + than + have + Noun} (ou verbe d'action approprié : \eng{eat, drink, choose}).
        \item \eng{prefer + V-ing + to + V-ing} $\rightarrow$ \eng{would rather + Base Verbale + than + Base Verbale}.
        \item \eng{prefer + to V + to + to V} $\rightarrow$ \eng{would rather + Base Verbale + than + Base Verbale}.
    \end{itemize}
    \item \textbf{Vérifiez} : pas de \eng{to} après \eng{rather}, présence de \eng{than} (pas \eng{to}), verbe à la base verbale.
\end{enumerate}
\textbf{Exemple :} \eng{I prefer teaching to trading.} $\rightarrow$ \eng{I'd rather teach than trade.} (Français : « Je préférerais enseigner plutôt que commercer. »)
\end{methode}

\begin{exercice}[Exercice 3 : Transformation — \eng{prefer} $\rightarrow$ \eng{would rather}]
Réécrivez les phrases suivantes en utilisant \eng{would rather} (ou \eng{'d rather}). Conservez le sens exact.
\begin{enumerate}
    \item I prefer teaching to trading.
    \item She prefers to drive rather than walk.
    \item We prefer staying at home to going out.
    \item He prefers coffee to tea.
    \item They prefer working in an office to working in the field.
    \item My father prefers reading newspapers to watching TV.
\end{enumerate}
\end{exercice}

\begin{exercice}[Exercice 4 : QCM — Type Baccalauréat]
Choisissez la seule réponse correcte (a, b, c ou d).
\begin{enumerate}
    \item "Would you like tea or coffee?" \quad "I \_\_\_\_\_\_\_\_ coffee."
    \begin{tabular}{ll}
    a) would rather & b) prefer \\
    c) would prefer & d) like better
    \end{tabular}
    \item My sister \_\_\_\_\_\_\_\_ dancing \_\_\_\_\_\_\_\_ singing.
    \begin{tabular}{ll}
    a) prefers / than & b) would rather / than \\
    c) prefers / to & d) would prefer / to
    \end{tabular}
    \item \_\_\_\_\_\_\_\_ you \_\_\_\_\_\_\_\_ go to the cinema or stay at home?
    \begin{tabular}{ll}
    a) Would / rather & b) Do / prefer \\
    c) Would / prefer to & d) Are / liking better
    \end{tabular}
    \item He \_\_\_\_\_\_\_\_ not \_\_\_\_\_\_\_\_ mention it.
    \begin{tabular}{ll}
    a) would / rather & b) prefers / to \\
    c) would / prefer & d) had / better
    \end{tabular}
    \item I \_\_\_\_\_\_\_\_ live in Yaoundé \_\_\_\_\_\_\_\_ Douala.
    \begin{tabular}{ll}
    a) would rather / than & b) prefer / than \\
    c) like better / to & d) would prefer / than
    \end{tabular}
\end{enumerate}
\end{exercice}

\courssub{Production guidée : « My career preferences »}

\begin{experience}[Rédaction d'un paragraphe personnel]
\textbf{Consigne :} Rédigez un paragraphe de \textbf{6 à 8 lignes} (environ 80--100 mots) intitulé \eng{« My career preferences »}. Vous devez y exprimer vos préférences professionnelles futures (études, métier, lieu de travail, conditions) en utilisant \textbf{au moins trois structures différentes} parmi celles étudiées :
\begin{itemize}
    \item \eng{prefer + Noun/Gerund + to + Noun/Gerund}
    \item \eng{would rather + Base Verb + than + Base Verb}
    \item \eng{would prefer + to-Infinitive / Noun + to + Noun}
    \item \eng{like + Noun/Gerund + better than + Noun/Gerund}
\end{itemize}
\textbf{Contexte suggéré :} Vous êtes en Terminale à Yaoundé, Douala ou Buea. Vous hésitez entre plusieurs filières (médecine, droit, ingénierie, commerce, agriculture, informatique) ou types d'emploi (fonction publique, entreprise privée, entrepreneuriat, ONG).
\end{experience}

\begin{textebox}[Modèle de production attendue]
My career preferences

As a Terminale student in Yaoundé, I often think about my future. \textbf{I would rather study computer science than medicine} because I prefer working with data to working in hospitals. \textbf{I prefer programming to networking}, as I enjoy creating applications. However, \textbf{I would prefer to work for a tech startup rather than for the civil service}, because I like innovation better than routine. \textbf{I prefer living in Buea to living in Douala} for the climate, but I would rather travel abroad for a Master's degree than stay here. Ultimately, I'd rather be an entrepreneur than an employee.
\end{textebox}

\courssub{Correction collective et justifications}

\begin{corrige}[Exercice 1 : Complétion]
\begin{enumerate}
    \item \textbf{to} — Structure : \eng{prefer + Noun + to + Noun}.
    \item \textbf{rather} — Structure : \eng{would rather + Base Verb}. \eng{Would} est déjà présent.
    \item \textbf{to} — Structure : \eng{prefer + Gerund + to + Gerund}.
    \item \textbf{to} — Structure : \eng{would prefer + to-Infinitive}.
    \item \textbf{rather} — Structure : \eng{'d rather not + Base Verb} (négation : \eng{not} après \eng{rather}).
    \item \textbf{than} — Structure : \eng{like + Noun + better than + Noun}.
    \item \textbf{to} — Structure : \eng{prefer + to-Infinitive} (forme moins courante mais correcte : \eng{prefer to travel}).
    \item \textbf{than} — Structure : \eng{would rather + Base Verb + than + Base Verb}.
    \item \textbf{to} — Structure : \eng{prefer + Noun + to + Noun}.
    \item \textbf{to} — Structure : \eng{would prefer + Noun + to + Noun}.
\end{enumerate}
\end{corrige}

\begin{corrige}[Exercice 2 : Correction d'erreurs]
\begin{enumerate}
    \item \eng{I prefer tea \textbf{to} coffee.} — Règle : \eng{prefer} s'emploie avec \eng{to} (comparaison), jamais \eng{than}.
    \item \eng{He would rather \textbf{stay} here.} — Règle : \eng{would rather} + \textbf{base verbale} (sans \eng{to}).
    \item \eng{She prefers \textbf{reading to writing}.} — Règle : \eng{prefer + V-ing + to + V-ing} (parallélisme gérondif). \eng{Than} est interdit après \eng{prefer}.
    \item \eng{I \textbf{would} rather go now.} — Règle : \eng{Had rather} est archaïque/rare ; l'anglais moderne utilise \eng{would rather} (\eng{'d rather}).
    \item \eng{They prefer fish \textbf{to} meat.} — Règle : \eng{prefer} + \eng{to} pour comparer deux noms.
    \item \eng{Would you prefer \textbf{to go} or \textbf{to stay}?} (ou \eng{going or staying}). — Règle : Parallélisme obligatoire après \eng{would prefer} (soit deux infinitifs, soit deux gérondifs).
    \item \eng{I \textbf{prefer tea to coffee}.} (ou \eng{I like tea \textbf{better than} coffee}). — Règle : \eng{Like better than} existe mais l'ordre est \eng{like + Objet + better than + Objet}. \eng{I like better tea} est faux (ordre des mots).
    \item \eng{He'd rather \textbf{have} coffee.} (ou \eng{He prefers coffee}). — Règle : \eng{Would rather} ne peut pas être suivi d'un verbe conjugué (\eng{prefers}). Il faut un verbe d'action (\eng{have, take, drink}) à la base verbale.
\end{enumerate}
\end{corrige}

\begin{corrige}[Exercice 3 : Transformation]
\begin{enumerate}
    \item \eng{I'd rather teach than trade.} (Noms d'action $\rightarrow$ verbes à la base verbale).
    \item \eng{She'd rather drive than walk.} (Infinitifs \eng{to drive / walk} $\rightarrow$ base verbale ; \eng{rather than} dans l'original signale déjà l'opposition).
    \item \eng{We'd rather stay at home than go out.} (Gérondifs $\rightarrow$ base verbale).
    \item \eng{He'd rather have coffee than tea.} (Noms concrets $\rightarrow$ ajout du verbe \eng{have} pour actionnaliser).
    \item \eng{They'd rather work in an office than work in the field.} (Gérondifs $\rightarrow$ base verbale ; répétition du verbe ou ellipsis : \eng{than in the field} possible mais moins formel).
    \item \eng{My father would rather read newspapers than watch TV.} (Gérondifs $\rightarrow$ base verbale ; forme non contractée préférée à l'écrit formel).
\end{enumerate}
\end{corrige}

\begin{corrige}[Exercice 4 : QCM]
\begin{enumerate}
    \item \textbf{b) prefer} — Réponse courte à un choix \eng{or} : \eng{I prefer coffee.} (\eng{Would rather} attend \eng{rather than...} ; \eng{Would prefer} attend \eng{to}).
    \item \textbf{c) prefers / to} — Sujet 3e pers. sing. \eng{My sister prefers} + \eng{to} (comparaison noms/gérondifs).
    \item \textbf{a) Would / rather} — Question avec \eng{would rather} + base verbale (\eng{go, stay}). \eng{Would prefer} exige \eng{to}.
    \item \textbf{a) would / rather} — Négation de \eng{would rather} : \eng{would rather not + V\textsubscript{base}}. \eng{Had better} donne un conseil, pas une préférence. \textit{Note : le verbe doit être à la base verbale \eng{mention}, non \eng{mentioned}.}
    \item \textbf{a) would rather / than} — Structure unique \eng{would rather + V + than + V}. \eng{Prefer} prend \eng{to} ; \eng{Would prefer} prend \eng{to}.
\end{enumerate}
\end{corrige}

\begin{corrige}[Production guidée : Analyse du modèle]
Le modèle fourni utilise les 4 structures :
\begin{itemize}
    \item Ligne 2 : \eng{\textbf{would rather study} ... \textbf{than} \textbf{medicine}} (verbes base verbale ; \eng{medicine} = nom, ellipsis de \eng{study}).
    \item Ligne 2 : \eng{\textbf{prefer working} ... \textbf{to working}} (gérondifs + \eng{to}).
    \item Ligne 3 : \eng{\textbf{prefer programming to networking}} (gérondifs + \eng{to}).
    \item Ligne 3 : \eng{\textbf{would prefer to work} ... \textbf{rather than for}} (structure mixte \eng{would prefer + to V + rather than} — variante courante à l'oral, mais la forme standard est \eng{would prefer to work ... to working} ou \eng{would rather work ... than work}).
    \item Ligne 4 : \eng{\textbf{like innovation better than routine}} (\eng{like + N + better than + N}).
    \item Ligne 4 : \eng{\textbf{prefer living} ... \textbf{to living}} (gérondifs + \eng{to}).
    \item Ligne 5 : \eng{\textbf{'d rather be} ... \textbf{than be}} (base verbale \eng{be}).
\end{itemize}
\end{corrige}

\courssub{Synthèse visuelle : Les 5 structures de la préférence}

\begin{popfigure}
\begin{tikzpicture}[
    mindmap,
    concept color=popblue,
    level 1/.style={level distance=4.5cm, sibling angle=72},
    level 2/.style={level distance=3cm, sibling angle=45},
    every node/.style={concept, execute at begin node=\hskip0pt, font=\small},
    text=white,
    grow cyclic
]
\node[concept, scale=1.2, align=center]{Expressing\\ Preferences}
    child[concept color=poporange] { node {Prefer + N/to N} 
        child { node {I prefer tea \textbf{to} coffee.} }
        child { node {Préférer un nom à un autre} }
    }
    child[concept color=popgreen] { node {Prefer + V-ing/to V-ing} 
        child { node {She prefers reading \textbf{to} writing.} }
        child { node {Préférer une action (gérondif)} }
    }
    child[concept color=poppurple] { node {Would rather + V/than V} 
        child { node {We'd rather stay \textbf{than} go.} }
        child { node {Préférence immédiate/spécifique (base verbale)} }
    }
    child[concept color=poppink] { node {Would prefer + to V/to N} 
        child { node {I'd prefer \textbf{to} leave early.} }
        child { node {Préférence polie/conditionnelle (infinitif)} }
    }
    child[concept color=popgold] { node {Like ... better than} 
        child { node {They like fish \textbf{better than} meat.} }
        child { node {Goût personnel, comparatif (than)} }
    };
\end{tikzpicture}
\legende{Carte mentale récapitulative des 5 structures clés pour exprimer la préférence en anglais. Chaque branche donne la forme canonique, un exemple et la nuance d'emploi.}
\end{popfigure}

\begin{attention}[Check-list avant de rendre votre production écrite]
Avant de considérer votre paragraphe « My career preferences » comme fini, vérifiez \textbf{chaque phrase} contenant une structure de préférence :
\begin{itemize}
    \item \textbf{Prefer} : Ai-je mis \cle{to} (et non \eng{than}) après le premier élément ? \eng{Prefer V-ing to V-ing} ou \eng{Prefer N to N}.
    \item \textbf{Would rather} : Le verbe qui suit est-il à la \cle{base verbale} (sans \eng{to}) ? Ai-je mis \cle{than} (et non \eng{to}) avant le second verbe ? \eng{'d rather V than V}.
    \item \textbf{Would prefer} : Ai-je mis \cle{to} devant le verbe (\eng{to V}) ou \cle{to} entre les deux noms (\eng{N to N}) ?
    \item \textbf{Like better} : L'ordre est-il \eng{Like + Objet + better than + Objet} ? (Pas \eng{Like better + Objet}).
    \item \textbf{Négation} : \eng{would rather \textbf{not} V} (pas \eng{don't would rather}).
\end{itemize}
\textbf{Piège ultime :} Ne mélangez pas les structures dans une même phrase (ex. \eng{*I would prefer rather go}).
\end{attention}

\begin{center}
\begin{tabularx}{\linewidth}{|X|X|X|X|}
\hline
\rowcolor{popblueL} \textbf{Structure} & \textbf{Forme canonique} & \textbf{Exemple} & \textbf{Traduction / Nuance} \\
\hline
\cle{Prefer} (Noms) & prefer + N + \textbf{to} + N & I prefer coffee \textbf{to} tea. & Je préfère le café au thé. (Goût général) \\
\hline
\cle{Prefer} (Actions) & prefer + V-ing + \textbf{to} + V-ing & He prefers reading \textbf{to} writing. & Il préfère lire plutôt qu'écrire. \\
\hline
\cle{Would rather} & 'd rather + \textbf{V\textsubscript{base}} + \textbf{than} + V\textsubscript{base} & We'd rather walk \textbf{than} run. & Nous préférerions marcher plutôt que courir. (Choix précis, immédiat) \\
\hline
\cle{Would prefer} & 'd prefer + \textbf{to} V / N + \textbf{to} N & I'd prefer \textbf{to} stay. / I'd prefer coffee \textbf{to} tea. & Je préférerais rester. (Politesse, conditionnel) \\
\hline
\cle{Like better} & like + N/V-ing + \textbf{better than} + N/V-ing & She likes dancing \textbf{better than} singing. & Elle aime danser mieux que chanter. (Goût personnel, comparatif) \\
\hline
\end{tabularx}
\end{center}

\newpage

\coursec{Activité d'intégration}

\coursec{Lesson 20 — Grammar: Expressing preferences}

\courssub{Observation et découverte}

\begin{textebox}[National Youth Job Fair — Official Announcement]
\textbf{YAOUNDÉ NATIONAL YOUTH JOB FAIR — 2025 EDITION}

Are you about to finish secondary school? Do you dream of working in a bank, a mobile phone company, a hotel or an international NGO? Come and meet more than forty employers at the Yaoundé Conference Centre!

During the fair, every candidate will have a five-minute interview with a recruiter. Be ready to answer questions such as: “Would you rather work in an office or in the field?”, “Do you prefer working alone or in a team?”, “Which do you like better: a high salary or flexible working hours?”

Bring your CV, speak clearly, and do not forget to justify your choices. The best candidates will be offered a two-week internship.

\textit{Glossaire :} field (n.) : le terrain ; flexible (adj.) : souple ; internship (n.) : stage ; recruiter (n.) : recruteur ; salary (n.) : salaire.
\end{textebox}

\begin{experience}[Observation guidée]
\begin{enumerate}
\item Lis le texte ci-dessus et surligne toutes les expressions qui servent à \textbf{demander} ou à \textbf{exprimer une préférence}.
\item Classe-les dans trois colonnes selon la structure utilisée : \eng{Would you rather...?} / \eng{Do you prefer...?} / \eng{Which do you like better...?}
\item Quelle préposition suit \eng{prefer} dans l'exemple \eng{“Do you prefer working alone or in a team?”} ? (Indice : le texte utilise \eng{or}, mais la structure complète est \eng{prefer ... to}).
\item Après \eng{would rather}, le verbe est-il à l'infinitif avec \eng{to} ou à la base verbale (sans \eng{to}) ?
\end{enumerate}
\end{experience}

\courssub{Prefer : règle et emplois}

\begin{propriete}[Prefer — Formes et emplois]
\eng{Prefer} exprime une préférence générale, habituelle ou un choix entre deux options. Il se construit de trois façons principales :

\begin{center}
\begin{tabularx}{\linewidth}{|X|X|X|}
\hline
\rowcolor{popblueL}
\textbf{Structure} & \textbf{Exemple} & \textbf{Traduction / Note} \\
\hline
\eng{Prefer + Noun + to + Noun} & \eng{I prefer tea \textbf{to} coffee.} & « Je préfère le thé au café. » (Préférence nominale) \\
\hline
\eng{Prefer + V-ing + to + V-ing} & \eng{I prefer \textbf{working} in a team \textbf{to} \textbf{working} alone.} & « Je préfère travailler en équipe plutôt que seul. » (Action générale) \\
\hline
\eng{Prefer + to-inf + rather than + bare inf} & \eng{I prefer \textbf{to work} early \textbf{rather than} \textbf{work} late.} & « Je préfère travailler tôt plutôt que tard. » (Choix spécifique, plus formel) \\
\hline
\end{tabularx}
\end{center}

\textbf{Points clés :}
\begin{itemize}
\item La préposition après \eng{prefer} (noun/gerund) est toujours \textbf{\eng{to}}, jamais \eng{than}.
\item \eng{Rather than} s'emploie avec l'infinitif \eng{to} (structure 3) pour contraster deux infinitifs.
\item \eng{Prefer} n'est généralement pas utilisé au \eng{be + -ing} (pas de \eng{*I am preferring}).
\end{itemize}
\end{propriete}

\begin{exemplebox}[Exemples variés avec \eng{prefer}]
\begin{itemize}
\item \eng{Most students prefer reading novels to watching videos.} (La plupart des élèves préfèrent lire des romans à regarder des vidéos.)
\item \eng{She prefers to write her notes by hand rather than type them.} (Elle préfère écrire ses notes à la main plutôt que les taper.)
\item \eng{We prefer fish to meat.} (Nous préférons le poisson à la viande.)
\end{itemize}
\end{exemplebox}

\begin{attention}[Piège majeur : \eng{Prefer to} vs \eng{Prefer than}]
\begin{itemize}
\item \textbf{Incorrect} : \eng{*I prefer coffee than tea.} / \eng{*I prefer working than sleeping.}
\item \textbf{Correct} : \eng{I prefer coffee \textbf{to} tea.} / \eng{I prefer working \textbf{to} sleeping.}
\item \textit{Explication} : L'interférence du français « préférer \textbf{à} » (qui ressemble à « than » par la comparaison) pousse à l'erreur. En anglais, la comparaison de préférence avec \eng{prefer} se fait avec \textbf{\eng{to}}.
\end{itemize}
\end{attention}

\courssub{Would rather et would prefer}

\begin{propriete}[Would rather ('d rather) — Forme et emploi]
\eng{Would rather} (souvent contracté en \eng{'d rather}) exprime une \textbf{préférence spécifique, immédiate ou pour une action précise}. Il est suivi de la \textbf{base verbale (sans \eng{to})}.

\begin{center}
\begin{tabular}{|l|l|l|}
\hline
\rowcolor{poporangeL}
\textbf{Structure} & \textbf{Exemple} & \textbf{Traduction} \\
\hline
\eng{Subject + would rather + bare inf} & \eng{I'd rather \textbf{stay} here.} & Je préférerais rester ici. \\
\hline
\eng{Subject + would rather + bare inf + than + bare inf} & \eng{I'd rather \textbf{walk} \textbf{than} \textbf{take} a taxi.} & Je préférerais marcher plutôt que prendre un taxi. \\
\hline
\eng{Subject + would rather + not + bare inf} & \eng{I'd rather \textbf{not go} out tonight.} & Je préférerais ne pas sortir ce soir. \\
\hline
\end{tabular}
\end{center}

\textbf{Sujets différents (subjonctif) :} \eng{I'd rather \textbf{you came} tomorrow.} (Je préférerais que tu viennes demain.) $\rightarrow$ Prétérit modal (past subjunctive).
\end{propriete}

\begin{propriete}[Would prefer ('d prefer) — Forme et emploi]
\eng{Would prefer} est plus formel que \eng{would rather}. Il est suivi de l'\textbf{infinitif avec \eng{to}}.

\begin{center}
\begin{tabular}{|l|l|l|}
\hline
\rowcolor{popgreenL}
\textbf{Structure} & \textbf{Exemple} & \textbf{Traduction} \\
\hline
\eng{Subject + would prefer + to-inf} & \eng{I'd prefer \textbf{to leave} now.} & Je préférerais partir maintenant. \\
\hline
\eng{Subject + would prefer + to-inf + rather than + bare inf} & \eng{I'd prefer \textbf{to drive} \textbf{rather than} \textbf{fly}.} & Je préférerais conduire plutôt que voler. \\
\hline
\eng{Subject + would prefer + noun} & \eng{I'd prefer \textbf{coffee}.} & Je préférerais du café. \\
\hline
\end{tabular}
\end{center}
\end{propriete}

\begin{exemplebox}[Would rather vs Would prefer]
\begin{itemize}
\item \eng{A: “Would you like tea or coffee?” \quad B: “I'd \textbf{rather have} coffee.”} (Réponse rapide, courante)
\item \eng{A: “Shall we take the bus?” \quad B: “I'd \textbf{prefer to walk}.”} (Plus formel, choix réfléchi)
\item \eng{“I'd rather \textbf{not wait}.”} (Négation courante avec \eng{rather})
\item \eng{“I'd prefer \textbf{not to wait}.”} (Négation avec \eng{prefer})
\end{itemize}
\end{exemplebox}

\begin{attention}[Erreurs fréquentes : \eng{Would rather to} / \eng{Would prefer than}]
\begin{itemize}
\item \textbf{Incorrect} : \eng{*I would rather \textbf{to go}.} $\rightarrow$ \textbf{Correct} : \eng{I would rather \textbf{go}.} (Base verbale !)
\item \textbf{Incorrect} : \eng{*I would prefer \textbf{than} stay.} $\rightarrow$ \textbf{Correct} : \eng{I would prefer \textbf{to stay}.} / \eng{I would rather \textbf{stay}.}
\item \textbf{Incorrect} : \eng{*I would rather \textbf{not to go}.} $\rightarrow$ \textbf{Correct} : \eng{I would rather \textbf{not go}.} (Pas de \eng{to} après \eng{not} non plus).
\end{itemize}
\end{attention}

\courssub{Like... better than et autres structures}

\begin{propriete}[Like ... better than / Like ... more than]
Pour exprimer une préférence de manière moins formelle, on utilise \eng{like} + \textbf{\eng{better than}} (ou \eng{more than}) + \textbf{Noun / Pronoun / Gerund}.

\begin{center}
\begin{tabularx}{\linewidth}{|X|X|}
\hline
\rowcolor{poppurpleL}
\textbf{Structure} & \textbf{Exemple} \\
\hline
\eng{Like + Noun + better than + Noun} & \eng{I like \textbf{apples better than} oranges.} \\
\hline
\eng{Like + V-ing + better than + V-ing} & \eng{She likes \textbf{swimming better than} \textbf{running}.} \\
\hline
\eng{Like + Pronoun + better than + Pronoun} & \eng{We like \textbf{this one better than} \textbf{that one}.} \\
\hline
\end{tabularx}
\end{center}

\textbf{Variantes :} \eng{Enjoy ... more than} (apprécier davantage), \eng{Fancy + V-ing} (avoir envie de, familier UK), \eng{Be keen on + V-ing} (être enthousiaste à l'idée de).
\end{propriete}

\begin{exemplebox}[Autres structures de préférence]
\begin{itemize}
\item \eng{I \textbf{enjoy reading more than watching} TV.} (J'apprécie la lecture plus que la télé.)
\item \eng{He \textbf{fancies going} to the cinema tonight.} (Il a envie d'aller au ciné ce soir — UK, familier.)
\item \eng{They \textbf{are keen on learning} English.} (Ils sont motivés pour apprendre l'anglais.)
\end{itemize}
\end{exemplebox}

\courssub{Comparaison français/anglais et pièges}

\begin{center}
\begin{tabularx}{\linewidth}{|X|X|X|}
\hline
\rowcolor{popblueL}
\textbf{Français} & \textbf{English (Correct)} & \textbf{Erreur fréquente (Francophonisme)} \\
\hline
Je préfère le thé \textbf{au} café. & \eng{I prefer tea \textbf{to} coffee.} & \eng{*I prefer tea \textbf{than} coffee.} \\
\hline
Je préfère travailler \textbf{plutôt que} dormir. & \eng{I prefer \textbf{working to sleeping}.} / \eng{I prefer \textbf{to work rather than sleep}.} & \eng{*I prefer working \textbf{than} sleeping.} \\
\hline
Je \textbf{préférerais} partir. & \eng{I'd \textbf{rather leave}.} / \eng{I'd \textbf{prefer to leave}.} & \eng{*I would rather \textbf{to leave}.} \\
\hline
Je \textbf{préférerais ne pas} y aller. & \eng{I'd \textbf{rather not go}.} / \eng{I'd \textbf{prefer not to go}.} & \eng{*I would rather \textbf{not to go}.} \\
\hline
J'aime \textbf{mieux} le foot que le basket. & \eng{I \textbf{like football better than} basketball.} & \eng{*I \textbf{like more} football than basketball.} \\
\hline
\end{tabularx}
\end{center}

\begin{aretenir}[Récapitulatif des 3 structures clés pour le Bac]
\begin{enumerate}
\item \textbf{Prefer + V-ing + \cle{to} + V-ing} \quad (Préférence générale) \quad $\rightarrow$ \eng{I prefer \textbf{studying to playing}.}
\item \textbf{Would rather + \cle{base verbale} + \cle{than} + base verbale} \quad (Préférence spécifique) \quad $\rightarrow$ \eng{I'd rather \textbf{study than play}.}
\item \textbf{Would prefer + \cle{to} + infinitif} \quad (Préférence spécifique, formel) \quad $\rightarrow$ \eng{I'd prefer \textbf{to study}.}
\end{enumerate}
\textbf{Astuce mnémotechnique :} \eng{Prefer} $\rightarrow$ \textbf{T}O (comme \eng{T}Owards) ; \eng{Rather} $\rightarrow$ \textbf{R}apide, pas de \eng{to} (base verbale) ; \eng{Prefer} (would) $\rightarrow$ \eng{TO} (infinitif).
\end{aretenir}

\courssub{Entraînement et production guidée : Simulation « National Youth Job Fair »}

\courssub{Objectifs de l'activité}

À la fin de cette activité d'intégration, tu seras capable de :
\begin{itemize}
\item utiliser spontanément, à l'oral, les structures de préférence étudiées dans la leçon : \eng{prefer ... to}, \eng{would rather ... than}, \eng{would prefer to}, \eng{like ... better than} ;
\item poser des questions sur les préférences professionnelles d'un interlocuteur et y répondre en justifiant tes choix ;
\item rédiger un court paragraphe de synthèse (environ huit lignes) intitulé \eng{My ideal job and my preferences}, en réutilisant au moins trois structures différentes de la leçon ;
\item repérer et corriger les erreurs fréquentes des francophones (\eng{prefer than}, \eng{would rather to work}, etc.) dans une production écrite.
\end{itemize}

\courssub{Situation}

La ville de Yaoundé organise chaque année, au Palais des Congrès, un grand salon de l'emploi destiné aux jeunes diplômés : le \eng{National Youth Job Fair}. Des banques, des entreprises de téléphonie mobile, des ONG, des hôtels et des administrations publiques y tiennent des stands et rencontrent les candidats. Ta classe de Terminale A a été invitée à participer à une simulation de ce salon. Voici l'annonce officielle publiée par les organisateurs.

\begin{textebox}[National Youth Job Fair — Official Announcement]
\textbf{YAOUNDÉ NATIONAL YOUTH JOB FAIR — 2025 EDITION}

Are you about to finish secondary school? Do you dream of working in a bank, a mobile phone company, a hotel or an international NGO? Come and meet more than forty employers at the Yaoundé Conference Centre!

During the fair, every candidate will have a five-minute interview with a recruiter. Be ready to answer questions such as: “Would you rather work in an office or in the field?”, “Do you prefer working alone or in a team?”, “Which do you like better: a high salary or flexible working hours?”

Bring your CV, speak clearly, and do not forget to justify your choices. The best candidates will be offered a two-week internship.

\textit{Glossaire :} field (n.) : le terrain ; flexible (adj.) : souple ; internship (n.) : stage ; recruiter (n.) : recruteur ; salary (n.) : salaire.
\end{textebox}

Tu vas jouer deux rôles : d'abord celui du \cle{recruteur}, qui interroge le candidat sur ses goûts professionnels, puis celui du \cle{candidat}, qui doit exprimer et justifier ses préférences en anglais correct. À la fin de la simulation, chacun rédigera un paragraphe de synthèse qui pourra être lu et corrigé collectivement au tableau.

\courssub{Consignes}

\begin{enumerate}
\item \textbf{Préparation du lexique (5 min).} En silence, lis l'annonce du salon. Relève cinq mots ou expressions liés au monde du travail (\eng{employer}, \eng{salary}, \eng{internship}...) et écris leur traduction française dans ton cahier.
\item \textbf{Préparation des questions (10 min).} Par binômes, l'élève A (recruteur) rédige au moins cinq questions sur les préférences professionnelles, en utilisant des structures variées : \eng{Would you rather...?}, \eng{Do you prefer... or...?}, \eng{Which do you like better...?}. L'élève B (candidat) prépare mentalement ses réponses en notant trois structures différentes à employer.
\item \textbf{Simulation de l'entretien (10 min).} Jouez la scène. Le recruteur pose ses questions ; le candidat répond en phrases complètes et \textbf{justifie} chaque choix avec \eng{because}, \eng{as} ou \eng{since}. Le candidat doit utiliser au minimum : une fois \eng{prefer ... to}, une fois \eng{would rather ... than}, une fois \eng{like ... better than}.
\item \textbf{Inversion des rôles (10 min).} Échangez les rôles et recommencez l'entretien avec de nouvelles questions.
\item \textbf{Production écrite (15 min).} Chaque élève rédige individuellement un paragraphe d'environ huit lignes intitulé \eng{My ideal job and my preferences}. Le paragraphe doit contenir au moins trois structures de préférence différentes et deux justifications.
\item \textbf{Correction collective (10 min).} L'enseignant choisit deux ou trois paragraphes, les lit en classe et la classe corrige collectivement les erreurs au tableau, en particulier celles qui portent sur les prépositions (\eng{to}, pas \eng{than} après \eng{prefer}) et sur la base verbale après \eng{would rather}.
\end{enumerate}

\courssub{Ressources à mobiliser}

\begin{aretenir}[Boîte à outils de l'activité]
\textbf{Structures de préférence à réutiliser :}
\begin{itemize}
\item \eng{I prefer working in a team to working alone.} (prefer + V-ing + \textbf{to} + V-ing)
\item \eng{I would rather work in a bank than work in a factory.} (would rather + base verbale + \textbf{than} + base verbale)
\item \eng{I would prefer to work for an NGO.} (would prefer + to + base verbale)
\item \eng{I like teaching better than selling.} (like ... better than)
\end{itemize}
\textbf{Lexique utile :} \eng{employer} (employeur), \eng{employee} (employé), \eng{salary} (salaire), \eng{working hours} (horaires de travail), \eng{teamwork} (travail d'équipe), \eng{promotion} (promotion/avancement), \eng{internship} (stage), \eng{to apply for a job} (postuler à un emploi).\\
\textbf{Pour justifier :} \eng{because}, \eng{as}, \eng{since}, \eng{that is why}.
\end{aretenir}

\begin{attention}[Pièges à éviter pendant l'activité]
\begin{itemize}
\item Après \eng{prefer}, la préposition est \textbf{to}, jamais \eng{than} : \eng{I prefer tea \textbf{to} coffee.} (et non \eng{prefer tea than coffee}).
\item Après \eng{would rather}, utilise la \textbf{base verbale} sans \eng{to} : \eng{I would rather \textbf{stay} here.} (et non \eng{would rather to stay}).
\item Ne traduis pas mot à mot « j'aime mieux » par \eng{I like more} ; dis \eng{I like ... better} ou \eng{I prefer}.
\end{itemize}
\end{attention}

\courssub{Proposition de production et corrigé commenté}

\begin{exoresolu}[My ideal job and my preferences — production modèle]
Rédige un paragraphe d'environ huit lignes intitulé \eng{My ideal job and my preferences}, en utilisant au moins trois structures de préférence différentes et en justifiant tes choix.
\tcblower
\textbf{Production modèle :}

\eng{My ideal job is to work as a journalist for a television channel in Douala. I prefer working in a team to working alone, because I enjoy sharing ideas with other people. I would rather work in the field than stay in an office all day, as I like meeting people and discovering new places. I would prefer to have flexible working hours, since I am more productive in the evening. Finally, I like a moderate salary with interesting missions better than a high salary with boring tasks. That is why journalism is the perfect job for me.}

\textbf{Commentaire des choix de langue :}
\begin{itemize}
\item \eng{I prefer working in a team to working alone} : structure \eng{prefer + V-ing + to + V-ing} ; la préposition \eng{to} est correctement employée (et non \eng{than}).
\item \eng{I would rather work ... than stay ...} : après \eng{would rather}, les deux verbes sont à la base verbale (\eng{work}, \eng{stay}), sans \eng{to}.
\item \eng{I would prefer to have flexible working hours} : \eng{would prefer} est suivi de l'infinitif avec \eng{to}.
\item \eng{I like a moderate salary ... better than a high salary} : variante correcte de \eng{like ... better than}, appliquée ici à des noms et non à des verbes.
\item Les justifications sont introduites par \eng{because}, \eng{as}, \eng{since}, ce qui évite la répétition et enrichit le texte.
\item Le connecteur final \eng{That is why} conclut logiquement le paragraphe : compétence d'organisation du texte exigée en Terminale.
\end{itemize}
\end{exoresolu}

\courssub{Bilan de l'activité}

Cette activité t'a permis de mobiliser l'ensemble des savoir-faire de la leçon dans une situation de communication authentique et proche de ta future vie professionnelle. À l'oral, tu as appris à \textbf{poser des questions} sur les préférences et à \textbf{justifier tes choix}, deux compétences essentielles pour tout entretien d'embauche. À l'écrit, tu as produit un paragraphe structuré qui combine plusieurs structures de préférence sans les confondre.

Retiens les trois points critiques vérifiés pendant la correction collective : \eng{prefer} se construit avec \eng{to}, \eng{would rather} est suivi de la base verbale, et \eng{like ... better than} permet de comparer deux noms ou deux activités. Si tu maîtrises ces trois structures et que tu sais les alterner dans un même texte, tu es capable d'exprimer tes goûts et tes choix professionnels de manière claire, variée et correcte — exactement ce qu'attend un recruteur anglophone. Garde ton paragraphe dans ton portfolio : il pourra servir de base à une lettre de motivation ou à un entretien réel, au lycée comme au-delà.

\newpage

\coursec{Méthodes et exercices résolus}

\coursec{Lesson 20 — Grammar: Expressing preferences}

\begin{textebox}[Observation : Dialogue au carrefour Wada (Yaoundé)]
\eng{— “What kind of job would you like to do after university?”}
\eng{— “I \textbf{prefer working} for an international NGO. The salaries are good and you travel.”}
\eng{— “Really? I \textbf{would rather start} my own business. I \textbf{like being} my own boss \textbf{better than} taking orders.”}
\eng{— “My sister \textbf{would prefer to work} in a bank. She says it’s more stable.”}
\eng{— “We all have different priorities. But nobody \textbf{would rather stay} unemployed!”}
\end{textebox}

\begin{definition}[Structures clés pour exprimer la préférence]
\begin{itemize}
\item \textbf{Prefer + gérondif / nom} : préférence générale, goût durable. Ex. \eng{I prefer teaching to trading.}
\item \textbf{Would rather + base verbale (sans \eng{to})} : choix ponctuel, situation précise. Ex. \eng{I would rather stay in Douala.}
\item \textbf{Would prefer + to + base verbale} : choix ponctuel, plus formel. Ex. \eng{She would prefer to work in a bank.}
\item \textbf{Like ... better than} : préférence comparative, registre courant. Ex. \eng{I like farming better than trading.}
\end{itemize}
\end{definition}

\courssub{Savoir-faire 1 : Identifier et construire correctement les structures de préférence}

\begin{propriete}[Règle de base : Prefer (goût général)]
\begin{tabular}{|l|l|l|}
\hline
\rowcolor{popblueL} Structure & Forme verbale & Exemple \\
\hline
\eng{prefer + nom} & nom / pronom & \eng{I prefer tea.} / \eng{I prefer it.} \\
\hline
\eng{prefer + V-ing} & gérondif & \eng{I prefer drinking palm wine.} \\
\hline
\eng{prefer X to Y} & comparaison & \eng{I prefer tea \textbf{to} coffee.} (jamais \eng{than}) \\
\hline
\end{tabular}
\par\medskip
\emph{Exception :} \eng{prefer + to + infinitif} possible pour une préférence spécifique future (\eng{I prefer to leave early tomorrow}), mais au Bac, on attend \eng{prefer + V-ing} pour l'habitude.
\end{propriete}

\begin{propriete}[Règle de base : Would rather / Would prefer (choix ponctuel)]
\begin{tabular}{|l|l|l|}
\hline
\rowcolor{poporangeL} Structure & Forme verbale & Exemple \\
\hline
\eng{would rather + BV} & base verbale \textbf{sans \eng{to}} & \eng{I would rather stay.} \\
\hline
\eng{would rather X than Y} & comparaison & \eng{I would rather walk \textbf{than} run.} \\
\hline
\eng{would prefer + to + BV} & infinitif \textbf{avec \eng{to}} & \eng{I would prefer to stay.} \\
\hline
\eng{would prefer X to Y} & nom / gérondif & \eng{I would prefer tea to coffee.} \\
\hline
\end{tabular}
\par\medskip
\emph{Négation :} \eng{would rather not + BV} / \eng{would prefer not to + BV}. \emph{Interrogation :} inversion \eng{Would you rather ...?} / \eng{Would you prefer ...?}
\end{propriete}

\begin{methode}[Choisir la bonne structure de préférence]
Pour exprimer une préférence en anglais, suis cette démarche en quatre étapes :
\begin{enumerate}
\item \textbf{Repère le type de préférence} : est-ce un goût général et durable (préférence d'habitude) ou un choix ponctuel (ce que je veux faire maintenant / dans une situation donnée) ?
\item \textbf{Choisis la structure adaptée} :
\begin{itemize}
\item Goût général : \eng{prefer + nom / V-ing} ou \eng{like ... better than} ;
\item Choix ponctuel : \eng{would rather + base verbale} ou \eng{would prefer + to + base verbale}.
\end{itemize}
\item \textbf{Vérifie la construction grammaticale} : après \eng{prefer}, emploie le nom ou le gérondif (\eng{I prefer tea.} / \eng{I prefer drinking tea.}) ; après \eng{would rather}, \textbf{jamais de \eng{to}} (\eng{I would rather stay.}).
\item \textbf{Relis la phrase} : contrôle la préposition (\eng{prefer X to Y}, \eng{would rather X than Y}) et la place du sujet.
\end{enumerate}
Réflexe : \eng{prefer} = préférence générale ; \eng{would rather / would prefer} = choix du moment.
\end{methode}

\begin{exoresolu}[Identifier la structure correcte]
\textbf{Énoncé.} Choose the correct option to complete each sentence.
\begin{enumerate}
\item I prefer (drink / drinking / to drinking) palm wine to beer at village festivals.
\item She would rather (to stay / staying / stay) in Douala this weekend.
\item My father would prefer (retire / retiring / to retire) at sixty.
\end{enumerate}
\tcblower
\textbf{Solution détaillée.}
\begin{enumerate}
\item \eng{I prefer \textbf{drinking} palm wine to \textbf{drinking} beer at village festivals.} — Goût général $\rightarrow$ \eng{prefer + V-ing}. La comparaison exige le parallélisme : gérondif / gérondif (ou nom / nom). \eng{Drink} (base verbale) impossible après \eng{prefer} seul ; \eng{to drinking} calque fautif du français.
\item \eng{She would rather \textbf{stay} in Douala this weekend.} — Après \eng{would rather}, \textbf{toujours} la base verbale sans \eng{to}. \eng{To stay} et \eng{staying} incorrects.
\item \eng{My father would prefer \textbf{to retire} at sixty.} — \eng{Would prefer} se construit avec l'infinitif en \eng{to} : \eng{would prefer to + BV}. \eng{Retire} seul et \eng{retiring} fautifs ici.
\end{enumerate}
\textbf{Conclusion.} La clé est de distinguer préférence générale (\eng{prefer + V-ing}) et choix ponctuel (\eng{would rather + BV}, \eng{would prefer to + BV}).
\end{exoresolu}

\courssub{Savoir-faire 2 : Éviter les erreurs fréquentes des francophones}

\begin{attention}[Les 5 pièges qui coûtent des points au Bac]
\begin{enumerate}
\item \textbf{\eng{I prefer tea than coffee.}} — \textbf{Faux !} Avec \eng{prefer}, la comparaison exige \eng{to} : \eng{I prefer tea \textbf{to} coffee.} C'est l'erreur n° 1 des candidats francophones (calque de « préférer ... à »).
\item \textbf{\eng{I would rather to stay.}} — \textbf{Faux !} \textbf{Jamais} de \eng{to} après \eng{would rather} : \eng{I would rather \textbf{stay}.} En revanche, \eng{would prefer} \textbf{exige} le \eng{to} : \eng{I would prefer \textbf{to} stay.} Ne confonds pas les deux.
\item \textbf{\eng{I prefer to stay home} pour un goût général.} — Pour une habitude, utilise le gérondif : \eng{I prefer \textbf{staying} at home.} L'infinitif après \eng{prefer} seul est réservé à des cas particuliers (préférence future précise) ; au lycée, retiens : \eng{prefer + V-ing}.
\item \textbf{\eng{I would rather don't go.}} — \textbf{Faux !} La négation se construit \textbf{sans auxiliaire} : \eng{I would rather \textbf{not} go.} Le \eng{not} se place directement avant la base verbale.
\item \textbf{Oubli du sujet ou ordre des mots fautif} : \eng{Prefer football to handball} (sans sujet) ou \eng{I like better football} sont fautifs. Dis \eng{I prefer football to handball} ou \eng{I like football \textbf{better than} handball} — sujet obligatoire et \eng{better} placé \textbf{après} le nom comparé.
\end{enumerate}
\end{attention}

\begin{methode}[Repérer et corriger les calques du français]
\begin{enumerate}
\item \textbf{Méfie-toi de l'infinitif} : le français dit « je préfère rester », mais l'anglais exige \eng{I prefer \textbf{staying}} (goût général) ou \eng{I would prefer \textbf{to stay}} (choix ponctuel) — jamais \eng{I prefer to stay} pour une habitude dans un exercice de grammaire scolaire.
\item \textbf{Contrôle la préposition de comparaison} : \eng{prefer X \textbf{to} Y} (pas \eng{than}) ; \eng{would rather X \textbf{than} Y} (pas \eng{to}).
\item \textbf{N'oublie pas le sujet} : en anglais, chaque phrase a un sujet exprimé : \eng{I prefer tea}, jamais \eng{Prefer tea}.
\item \textbf{Attention à la négation} : \eng{I would rather \textbf{not} go.} (« je préférerais ne pas y aller ») — le \eng{not} se place avant la base verbale.
\item \textbf{Attention à l'interrogation} : \eng{Would you rather stay or leave?} — inversion de \eng{would}, comme avec un auxiliaire.
\end{enumerate}
\end{methode}

\begin{exoresolu}[Corriger les erreurs]
\textbf{Énoncé.} Each sentence contains one mistake. Correct it.
\begin{enumerate}
\item I prefer fufu than rice.
\item She would rather to sell tomatoes at Mokolo market.
\item Would you prefer staying in Buea or moving to Bamenda?
\item He prefers don't work on Sundays.
\end{enumerate}
\tcblower
\textbf{Solution détaillée.}
\begin{enumerate}
\item \eng{I prefer fufu \textbf{to} rice.} — Avec \eng{prefer}, la comparaison se fait avec la préposition \eng{to}, jamais \eng{than}. Calque du français « je préfère le foufou au riz ».
\item \eng{She would rather \textbf{sell} tomatoes at Mokolo market.} — Après \eng{would rather}, base verbale sans \eng{to}. Le \eng{to} est un calque de « elle préférerait vendre ».
\item \eng{Would you prefer \textbf{to stay} in Buea or \textbf{move} to Bamenda?} — \eng{Would prefer} exige l'infinitif en \eng{to} ; le gérondif \eng{staying} est fautif après \eng{would prefer} dans ce type d'exercice (choix ponctuel entre deux actions).
\item \eng{He prefers \textbf{not to work} on Sundays.} ou \eng{He would rather \textbf{not work} on Sundays.} — La négation se construit avec \eng{prefer not to + BV} ou \eng{would rather not + BV} ; on ne peut pas coller \eng{don't} après \eng{prefers}.
\end{enumerate}
\textbf{Conclusion.} Trois pièges dominent : la préposition (\eng{to}/\eng{than}), l'infinitif en \eng{to} après \eng{would rather}, et la place de la négation.
\end{exoresolu}

\courssub{Savoir-faire 3 : Utiliser les structures dans des phrases et de courts textes}

\begin{methode}[Rédiger un court texte exprimant des préférences (type Bac)]
Pour une production écrite de type « describe your preferences » (sujet fréquent au Bac) :
\begin{enumerate}
\item \textbf{Annonce le thème} : \eng{When it comes to jobs, people have different preferences.}
\item \textbf{Varie les structures} : utilise au moins trois formes différentes — \eng{prefer + V-ing}, \eng{would rather + BV}, \eng{like ... better than} — pour montrer ta maîtrise.
\item \textbf{Justifie tes choix} avec \eng{because}, \eng{as}, \eng{since} : \eng{I prefer teaching because I enjoy working with young people.}
\item \textbf{Nuance} avec \eng{however}, \eng{on the other hand} pour comparer deux options.
\item \textbf{Conclus} par une préférence claire : \eng{All in all, I would rather work in a hospital than in an office.}
\end{enumerate}
Structure-type à retenir : \eng{I prefer A to B because ... However, if I had to choose today, I would rather ...}
\end{methode}

\begin{exoresolu}[Production guidée]
\textbf{Énoncé.} Write a short paragraph (60–80 words) about your occupational preferences. Use at least three different structures of preference.
\tcblower
\textbf{Solution détaillée.} Démarche : on choisit deux métiers à comparer (par exemple enseignant et commerçant), on annonce le thème, on varie les structures, on justifie, on conclut.

\textbf{Modèle de réponse :}

\eng{Young Cameroonians have different occupational preferences. Personally, I prefer teaching to trading because I love sharing knowledge with students. I like working in a school better than sitting in a shop all day, as a classroom is a lively place. However, if I could not find a teaching job, I would rather become a nurse than sell goods in the market, because helping sick people is also rewarding. All in all, I would prefer to work in the public service, where jobs are more stable.}

(Traduction : « Les jeunes Camerounais ont des préférences professionnelles différentes. Personnellement, je préfère enseigner plutôt que commercer parce que j'aime partager le savoir avec les élèves... »)

\textbf{Points de grammaire vérifiés :} \eng{prefer teaching to trading} (gérondif + \eng{to}) ; \eng{like ... better than} ; \eng{would rather become ... than sell} (bases verbales) ; \eng{would prefer to work} (infinitif en \eng{to}). Quatre structures variées, justification avec \eng{because} et \eng{as}, contraste avec \eng{however}.
\end{exoresolu}

\courssub{Savoir-faire 4 : Transformer des phrases d'une structure à l'autre}

\begin{methode}[Transformer prefer, would rather, would prefer]
Exercice classique de transformation : réécrire une phrase sans changer le sens.
\begin{enumerate}
\item \textbf{Identifie la structure de départ} et sa signification (goût général ou choix ponctuel).
\item \textbf{Choisis la structure cible équivalente} :
\begin{itemize}
\item \eng{I prefer tea to coffee.} $\rightarrow$ \eng{I like tea better than coffee.}
\item \eng{I prefer walking.} $\rightarrow$ \eng{I would rather walk.} (si choix ponctuel)
\item \eng{I would rather stay.} $\rightarrow$ \eng{I would prefer to stay.}
\end{itemize}
\item \textbf{Adapte la forme verbale} : gérondif après \eng{prefer}, base verbale après \eng{would rather}, infinitif en \eng{to} après \eng{would prefer}.
\item \textbf{Adapte la préposition} : \eng{to} avec \eng{prefer}, \eng{than} avec \eng{would rather} et \eng{better}.
\item \textbf{Vérifie que le sens est identique} et que le registre (général/ponctuel) est respecté.
\end{enumerate}
\end{methode}

\begin{exoresolu}[Transformations]
\textbf{Énoncé.} Rewrite each sentence as indicated, keeping the same meaning.
\begin{enumerate}
\item I prefer travelling by bus to travelling by bendskin. (use \eng{would rather})
\item She would rather work in a bank. (use \eng{would prefer})
\item I like cocoa farming better than petty trading. (use \eng{prefer})
\end{enumerate}
\tcblower
\textbf{Solution détaillée.}
\begin{enumerate}
\item \eng{I would rather \textbf{travel} by bus \textbf{than} (travel) by bendskin.} — On passe de \eng{prefer + V-ing ... to} à \eng{would rather + BV ... than}. La base verbale \eng{travel} remplace le gérondif \eng{travelling}, et \eng{than} remplace \eng{to}. La répétition de \eng{travel} est facultative (ellipse).
\item \eng{She would prefer \textbf{to work} in a bank.} — Équivalence directe : \eng{would rather + BV} devient \eng{would prefer to + BV}. Ne pas oublier le \eng{to}.
\item \eng{I prefer cocoa farming \textbf{to} petty trading.} — \eng{Like X better than Y} se transforme en \eng{prefer X to Y} : le \eng{than} disparaît au profit de \eng{to}, et \eng{better} disparaît.
\end{enumerate}
\textbf{Conclusion.} Toute transformation repose sur deux conversions : la forme verbale (gérondif / base verbale / infinitif) et la préposition (\eng{to} / \eng{than}).
\end{exoresolu}

\courssub{Savoir-faire 5 : Comprendre et exploiter un texte contenant des préférences}

\begin{methode}[Repérer les préférences dans un texte de compréhension]
\begin{enumerate}
\item \textbf{Souligne les marqueurs} : \eng{prefer}, \eng{would rather}, \eng{would prefer}, \eng{like ... better}, \eng{favourite}, \eng{rather than}.
\item \textbf{Identifie qui préfère quoi} : note le sujet de chaque structure (les questions du Bac portent souvent là-dessus : \eng{Who prefers what?}).
\item \textbf{Distingue goût général et choix ponctuel} : \eng{prefers} indique une habitude, \eng{would rather} un choix dans une situation donnée.
\item \textbf{Reformule} pour vérifier ta compréhension : traduis mentalement la préférence en français avant de répondre.
\item \textbf{Réponds avec des mots du texte}, en phrases complètes et à la bonne personne.
\end{enumerate}
\end{methode}

\begin{exoresolu}[Compréhension de l'écrit]
\textbf{Énoncé.} Read the text and answer the questions.

\begin{textebox}[Extrait]
\eng{“In Yaoundé, many young graduates face a difficult choice. Some prefer working for the government because civil servants enjoy job security. Others would rather start their own business than wait for years for a public service job. My sister, for instance, likes trading better than office work: she opened a small shop near Carrefour Wada. As for me, I would prefer to work for a private company, where salaries are often higher.”}
\end{textebox}

\begin{enumerate}
\item Why do some graduates prefer government jobs?
\item What would other graduates rather do?
\item What does the narrator's sister prefer?
\end{enumerate}
\tcblower
\textbf{Solution détaillée.}
\begin{enumerate}
\item \eng{Some graduates prefer government jobs because civil servants enjoy job security.} — La réponse se trouve après \eng{because} ; on reprend la structure \eng{prefer working} du texte.
\item \eng{They would rather start their own business than wait for years for a public service job.} — Marqueur : \eng{would rather ... than}. On remplace \eng{others} par \eng{they} et on garde la base verbale \eng{start}.
\item \eng{She prefers trading to office work.} ou \eng{She likes trading better than office work.} — Le texte dit \eng{likes trading better than office work} ; on peut citer la structure ou la transformer en \eng{prefer ... to} (compétence de reformulation).
\end{enumerate}
\textbf{Conclusion.} Les trois structures du texte (\eng{prefer working}, \eng{would rather start}, \eng{likes ... better than}) correspondent exactement aux trois points de grammaire de la leçon : le texte sert de modèle pour ta propre production.
\end{exoresolu}

\begin{savaistu}[Préférences et culture professionnelle au Cameroun]
Au Cameroun anglophone (Northwest, Southwest), l'expression \eng{“I’d rather”} (contraction orale de \eng{I would rather}) est très courante dans les discussions sur l'emploi : \eng{“I’d rather do business than work for somebody.”} L'esprit d'entreprise (\eng{entrepreneurship}) est valorisé. En zone francophone, le secteur public (\eng{civil service}) reste très prisé pour la sécurité de l'emploi (\eng{job security}). Au Bac, un texte peut opposer ces deux mentalités : utilise \eng{prefer} pour les tendances générales, \eng{would rather} pour les choix individuels face à une offre concrète.
\end{savaistu}

\begin{exercice}[Entraînement complet — Expressing preferences]
\begin{enumerate}
\item \textbf{Complète} avec la forme correcte du verbe entre parenthèses.
\begin{enumerate}
\item I prefer \_\_\_\_\_\_ (read) novels to \_\_\_\_\_\_ (watch) TV.
\item If you ask me, I would rather \_\_\_\_\_\_ (go) by bendskin than \_\_\_\_\_\_ (wait) for a bus.
\item My uncle would prefer \_\_\_\_\_\_ (invest) in cocoa farming.
\end{enumerate}
\item \textbf{Transforme} les phrases suivantes en utilisant le mot donné sans le changer.
\begin{enumerate}
\item I like plantains better than yams. (prefer)
\item They would rather stay in Buea. (would prefer)
\item She prefers working alone. (would rather)
\end{enumerate}
\item \textbf{Corrige} les erreurs dans ces phrases d'élèves.
\begin{enumerate}
\item I prefer to eat fufu than rice.
\item Would you rather to be a doctor or a teacher?
\item He prefers not working on Sundays.
\item We like better football than handball.
\end{enumerate}
\item \textbf{Production écrite :} Écris un paragraphe de 70–90 mots sur tes préférences pour un job d'été. Utilise au moins \textbf{quatre} structures de préférence différentes et justifie tes choix.
\end{enumerate}
\end{exercice}

\begin{corrige}[Exercice — Expressing preferences]
\begin{enumerate}
\item \textbf{Complète :}
\begin{enumerate}
\item I prefer \textbf{reading} novels to \textbf{watching} TV. (gérondifs parallèles après \eng{prefer ... to})
\item If you ask me, I would rather \textbf{go} by bendskin than \textbf{wait} for a bus. (bases verbales après \eng{would rather ... than})
\item My uncle would prefer \textbf{to invest} in cocoa farming. (infinitif \eng{to} après \eng{would prefer})
\end{enumerate}
\item \textbf{Transforme :}
\begin{enumerate}
\item I prefer plantains to yams.
\item They would prefer to stay in Buea.
\item She would rather work alone. (choix ponctuel implicite)
\end{enumerate}
\item \textbf{Corrige :}
\begin{enumerate}
\item I prefer \textbf{eating} fufu \textbf{to} rice. (ou \eng{I prefer fufu to rice.})
\item Would you rather \textbf{be} a doctor or a teacher? (pas de \eng{to})
\item He prefers \textbf{not working} on Sundays. (ou \eng{He would rather not work on Sundays.})
\item We \textbf{like football better than} handball. (\eng{better} après le nom)
\end{enumerate}
\item \textbf{Production écrite :} (Exemple de réponse)
\eng{During the long holidays, I want to earn some money. I prefer working in a restaurant to selling in the streets because the income is regular. I like serving customers better than carrying heavy loads, as it is less tiring. However, if I find an internship in an office, I would rather take it than stay in a restaurant, because it gives professional experience. All in all, I would prefer to work in a place where I can learn new skills.}
\end{enumerate}
\end{corrige}

\newpage

\coursec{Exercices}

\courssub{Je vérifie mes connaissances}

\begin{exercice}[QCM : la bonne structure]
Pour chaque phrase, choisis la bonne réponse (a, b, c ou d).
\begin{enumerate}
\item I'd rather \_\_\_\_\_ at home tonight. \quad a) to stay \quad b) stay \quad c) staying \quad d) stayed
\item She prefers teaching \_\_\_\_\_ trading. \quad a) than \quad b) to \quad c) rather \quad d) from
\item I would prefer \_\_\_\_\_ a nurse rather than a trader. \quad a) be \quad b) being \quad c) to be \quad d) been
\item My father likes farming \_\_\_\_\_ than office work. \quad a) most \quad b) more \quad c) better \quad d) best
\item “I'd rather you \_\_\_\_\_ your homework first.” \quad a) do \quad b) did \quad c) doing \quad d) to do
\end{enumerate}
\end{exercice}

\begin{exercice}[Vrai ou faux ? Justifie]
Dis si les affirmations suivantes sont vraies (V) ou fausses (F), puis justifie ta réponse en citant la règle.
\begin{enumerate}
\item After “would rather”, we use the infinitive with “to”.
\item “Prefer” can be followed by a noun, an infinitive with “to”, or a gerund (-ing).
\item In “I prefer cocoa to coffee”, the word “to” is a preposition, so we can also say “I prefer growing cocoa to selling coffee”.
\item “I'd rather you came tomorrow” uses a past tense to talk about a present or future situation.
\item “Would rather” and “had better” are followed by exactly the same verb form.
\end{enumerate}
\end{exercice}

\begin{exercice}[Vocabulaire : les mots des préférences]
Complète chaque définition avec le mot anglais qui convient, choisi dans la liste : \emph{preference, choice, option, keen on, fond of, would rather}.
\begin{enumerate}
\item A strong liking for one thing more than another : \_\_\_\_\_\_\_\_\_\_\_\_
\item The act of deciding between two or more possibilities : \_\_\_\_\_\_\_\_\_\_\_\_
\item One of the possibilities you can choose : \_\_\_\_\_\_\_\_\_\_\_\_
\item To be very interested in something (adjective phrase) : to be \_\_\_\_\_\_\_\_\_\_\_\_
\item To like someone or something with affection : to be \_\_\_\_\_\_\_\_\_\_\_\_
\item Expression used to say what you want to do instead of something else : \_\_\_\_\_\_\_\_\_\_\_\_
\end{enumerate}
\end{exercice}

\begin{exercice}[Conjugaison et formes à compléter]
Complète les phrases avec la forme correcte du verbe entre parenthèses.
\begin{enumerate}
\item I prefer \_\_\_\_\_\_\_\_ (study) economics to \_\_\_\_\_\_\_\_ (study) literature.
\item She would rather \_\_\_\_\_\_\_\_ (work) in a bank than in a market.
\item They would prefer \_\_\_\_\_\_\_\_ (open) a small business in Buea.
\item I'd rather you \_\_\_\_\_\_\_\_ (not / spend) all your money at the market.
\item We like \_\_\_\_\_\_\_\_ (buy) local products better than imported ones.
\end{enumerate}
\end{exercice}

\courssub{Je m'entraîne}

\begin{exercice}[Traduction]
Traduis les phrases suivantes en anglais en utilisant la structure indiquée entre parenthèses.
\begin{enumerate}
\item Je préfère l'agriculture au commerce. (prefer ... to ...)
\item J'aimerais mieux devenir ingénieur que chauffeur de taxi. (would rather ... than ...)
\item Mon frère aime mieux travailler à Douala qu'à Yaoundé. (like ... better than ...)
\item Elle préférerait étudier la comptabilité. (would prefer + to)
\item Je préférerais que tu choisisses une carrière dans la santé. (would rather + sujet + prétérit)
\end{enumerate}
\end{exercice}

\begin{exercice}[Transformations type Bac]
Réécris chaque phrase selon la consigne, sans changer le sens.
\begin{enumerate}
\item I prefer farming to trading. $\rightarrow$ Réécris avec \emph{would rather}.
\item I prefer farming to trading. $\rightarrow$ Réécris avec \emph{like ... better than}.
\item I would rather stay in the village. $\rightarrow$ Réécris avec \emph{prefer} + gerund.
\item She likes teaching better than selling. $\rightarrow$ Réécris avec \emph{would prefer} + infinitive.
\item I prefer to work for myself rather than for a company. $\rightarrow$ Réécris avec \emph{would rather ... than}.
\end{enumerate}
\end{exercice}

\begin{exercice}[Texte à trous : l'économie camerounaise]
Complète le texte suivant avec les mots : \emph{to, than, rather, better, prefer, would} (chaque mot peut être utilisé plusieurs fois).
\begin{textebox}[Choosing a job in Cameroon]
Many young Cameroonians \_\_\_\_\_\_\_ working in the informal sector \_\_\_\_\_\_\_ looking for a government job, because it is faster to start. However, others \_\_\_\_\_\_\_ rather have a stable salary. “I like trading \_\_\_\_\_\_\_ than farming,” says Alain, a shopkeeper in Douala, “because I earn money every day.” His sister, on the other hand, would \_\_\_\_\_\_\_ become a civil servant. “I would prefer \_\_\_\_\_\_\_ work for the state,” she explains. “It is safer \_\_\_\_\_\_\_ running a small business.”
\end{textebox}
\end{exercice}

\begin{exercice}[Appariements et correction d'erreurs]
\textbf{A. Appariement.} Associe chaque début de phrase (1--5) à sa fin logique (a--e).
\begin{enumerate}
\item I prefer rice
\item I'd rather take a bendskin
\item She would prefer to study
\item We like local food better
\item I'd rather you
\end{enumerate}
\begin{itemize}
\item[a)] than imported food.
\item[b)] medicine at the university.
\item[c)] to fufu corn.
\item[d)] told me the truth about your results.
\item[e)] than wait for a taxi in the rain.
\end{itemize}
\textbf{B. Correction d'erreurs.} Le texte suivant contient \textbf{cinq fautes} typiques de francophones dans l'expression des préférences. Relève-les et corrige-les.
\begin{textebox}[Career choices]
I prefer to be doctor than to be a teacher. I would rather to work in a hospital because I like better helping people than selling things. My parents prefer that I choose this career, and I am agree with them.
\end{textebox}
\end{exercice}

\courssub{Je me prépare au Bac}

\begin{exercice}[Compréhension et langue — type Bac]
Lis le texte suivant, puis réponds aux questions.
\begin{textebox}[Farming or trading? A young graduate's dilemma]
When Manka finished secondary school in Bamenda, everyone had an opinion about her future. Her father, a cocoa farmer, wanted her to study agriculture. “I would rather you took over the family farm,” he told her. “Farming feeds the nation.” Her mother, who sells second-hand clothes at the market, disagreed. “I prefer trading to farming,” she said. “You make money faster, and you meet people every day.”

Manka listened to both of them, but she had her own ideas. “I like working with computers better than working in the fields or in a shop,” she explained. “I would prefer to study computer science in Yaoundé rather than stay here.” In the end, her parents accepted her choice. “We would rather see you happy in a career you love than unhappy in one we chose for you,” her father finally admitted.
\end{textebox}
\textbf{A. Comprehension.} Answer in English, in complete sentences.
\begin{enumerate}
\item What did Manka's father want her to do? Why?
\item Why does Manka's mother prefer trading to farming?
\item What career does Manka prefer, and where does she want to study?
\item Do her parents finally accept her decision? Justify with a sentence from the text.
\end{enumerate}
\textbf{B. Language work.}
\begin{enumerate}
\item Find in the text two different structures used to express a preference and copy them.
\item Rewrite with \emph{would rather}: “I prefer trading to farming.”
\item Rewrite with \emph{prefer ... to ...}: “I like working with computers better than working in the fields.”
\item Explain why the father says “I would rather you \emph{took} over” and not “you \emph{take} over”.
\end{enumerate}
\end{exercice}

\begin{exercice}[Traduction — type Bac]
Traduis en anglais le passage suivant.
\begin{textebox}[Texte à traduire]
Beaucoup de jeunes Camerounais préfèrent chercher du travail en ville plutôt que de rester au village. Ils aiment mieux les métiers modernes que l'agriculture traditionnelle. Pourtant, certains préféreraient devenir agriculteurs, car ils pensent que la terre peut les rendre riches. « J'aimerais mieux que mes enfants choisissent un métier qu'ils aiment », dit un père de famille à Bafoussam.
\end{textebox}
\end{exercice}

\begin{exercice}[Expression écrite — type Bac]
\textbf{Composition.} Write a paragraph of about \textbf{100 words} on the following topic:
\begin{textebox}[Topic]
“My career preferences” — Say which career you prefer and which one you would not like to do. Give your reasons.
\end{textebox}
Consignes obligatoires :
\begin{itemize}
\item Utilise \textbf{au moins trois structures différentes} pour exprimer tes préférences : \emph{prefer ... to ...}, \emph{would rather ... than ...}, \emph{would prefer to ...}, \emph{like ... better than ...}.
\item Situe ta réflexion dans le contexte camerounais (métiers, villes, réalités locales).
\item Soigne la grammaire, l'orthographe et la ponctuation ; souligne les trois structures de préférence utilisées.
\end{itemize}
\end{exercice}

\newpage

\coursec{Corrigés des exercices}

\coursec{Corrigés des exercices de la leçon 20}

\begin{corrige}[Exercice 1 — QCM : la bonne structure]
\begin{enumerate}
\item \textbf{b) stay} — Après \eng{would rather}, on utilise la base verbale (infinitif sans \eng{to}) : \eng{I'd rather stay at home tonight.} « Je préférerais rester à la maison ce soir. »
\item \textbf{b) to} — La structure est \eng{prefer + -ing + to + -ing} : \eng{She prefers teaching to trading.} « Elle préfère enseigner au commerce. » Attention : ce \eng{to} est une préposition, pas la marque de l'infinitif.
\item \textbf{c) to be} — Après \eng{would prefer}, on utilise l'infinitif avec \eng{to} : \eng{I would prefer to be a nurse rather than a trader.} « Je préférerais être infirmière plutôt que commerçante. »
\item \textbf{c) better} — La structure est \eng{like ... better than ...} : \eng{My father likes farming better than office work.} « Mon père aime mieux l'agriculture que le travail de bureau. » \eng{Better} est le comparatif de \eng{well}.
\item \textbf{b) did} — Après \eng{would rather + sujet}, le verbe se met au prétérit pour parler du présent ou du futur : \eng{I'd rather you did your homework first.} « Je préférerais que tu fasses d'abord tes devoirs. »
\end{enumerate}
\textbf{Point de vigilance :} ne jamais écrire \eng{would rather to stay} (faute n° 1 des francophones) ni \eng{prefer ... than} au lieu de \eng{prefer ... to}.
\end{corrige}

\begin{corrige}[Exercice 2 — Vrai ou faux ? Justifie]
\begin{enumerate}
\item \textbf{F} — After \eng{would rather}, we use the infinitive \emph{without} \eng{to} (the bare infinitive): \eng{I'd rather stay}, not \eng{I'd rather to stay}.
\item \textbf{V} — \eng{Prefer} can be followed by a noun (\eng{I prefer tea}), an infinitive with \eng{to} (\eng{I prefer to walk}) or a gerund (\eng{I prefer walking}).
\item \textbf{V} — In \eng{I prefer cocoa to coffee}, \eng{to} is a preposition, so a verb after it takes the \eng{-ing} form: \eng{I prefer growing cocoa to selling coffee.} Both verbs must be parallel (gerund + gerund).
\item \textbf{V} — After \eng{would rather + subject}, the past tense expresses a present or future wish: \eng{I'd rather you came tomorrow} = « je préférerais que tu viennes demain ».
\item \textbf{V} — Both \eng{would rather} and \eng{had better} are followed by the bare infinitive: \eng{I'd rather go} / \eng{You'd better go}. (Mais attention : le sens diffère — préférence contre conseil.)
\end{enumerate}
\end{corrige}

\begin{corrige}[Exercice 3 — Vocabulaire : les mots des préférences]
\begin{enumerate}
\item \textbf{preference} — \eng{a preference} (nom) : « une préférence ».
\item \textbf{choice} — \eng{a choice} (nom) : « un choix ». Attention au faux ami : \eng{choice} ne se prononce pas comme « choix » mais « tchoïce » (/tʃɔɪs/).
\item \textbf{option} — \eng{an option} (nom) : « une possibilité, une option ».
\item \textbf{keen on} — \eng{to be keen on + -ing} : « être passionné par, aimer beaucoup » : \eng{She is keen on computing.}
\item \textbf{fond of} — \eng{to be fond of + nom/-ing} : « aimer (avec affection) » : \eng{He is fond of his grandmother.}
\item \textbf{would rather} — \eng{I'd rather stay home.} : « je préférerais / j'aimerais mieux ».
\end{enumerate}
\textbf{Point de vigilance :} \eng{keen on} et \eng{fond of} sont suivis d'un nom ou d'un gérondif, jamais de l'infinitif : \eng{fond of reading}, pas \eng{fond of to read}.
\end{corrige}

\begin{corrige}[Exercice 4 — Conjugaison et formes à compléter]
\begin{enumerate}
\item \eng{I prefer \textbf{studying} economics to \textbf{studying} literature.} — Gérondif des deux côtés (parallélisme obligatoire). « Je préfère étudier l'économie à la littérature. »
\item \eng{She would rather \textbf{work} in a bank than in a market.} — Base verbale après \eng{would rather}. « Elle préférerait travailler dans une banque que dans un marché. »
\item \eng{They would prefer \textbf{to open} a small business in Buea.} — Infinitif avec \eng{to} après \eng{would prefer}. « Ils préféreraient ouvrir une petite entreprise à Buea. »
\item \eng{I'd rather you \textbf{didn't spend} all your money at the market.} — Prétérit négatif après \eng{would rather + sujet}. « Je préférerais que tu ne dépenses pas tout ton argent au marché. »
\item \eng{We like \textbf{buying} local products better than imported ones.} — Gérondif après \eng{like}. « Nous aimons mieux acheter des produits locaux que des produits importés. »
\end{enumerate}
\end{corrige}

\begin{corrige}[Exercice 5 — Traduction]
\begin{enumerate}
\item \eng{I prefer farming to trading.} — Structure \eng{prefer + nom + to + nom}. Ne jamais écrire \eng{than} ici.
\item \eng{I would rather become an engineer than a taxi driver.} — Base verbale après \eng{would rather} ; \eng{than} introduit la seconde option. Attention à l'article : \eng{an engineer}.
\item \eng{My brother likes working in Douala better than working in Yaoundé.} — \eng{like + -ing ... better than + -ing}. On peut aussi dire \eng{My brother likes Douala better than Yaoundé} si l'on compare les lieux.
\item \eng{She would prefer to study accounting.} — Infinitif avec \eng{to} après \eng{would prefer}. « Accounting » = « la comptabilité ».
\item \eng{I would rather you chose a career in health.} — Prétérit (\eng{chose}) après \eng{would rather + sujet} pour traduire le subjonctif français « que tu choisisses ».
\end{enumerate}
\textbf{Point de vigilance :} le subjonctif français après « j'aimerais mieux que » se rend en anglais par \eng{would rather + sujet + prétérit}, jamais par un subjonctif.
\end{corrige}

\begin{corrige}[Exercice 6 — Transformations type Bac]
\begin{enumerate}
\item \eng{I prefer farming to trading.} $\rightarrow$ \eng{I would rather farm than trade.} (ou \eng{I would rather be a farmer than a trader.}) — Base verbale des deux côtés de \eng{than}.
\item \eng{I prefer farming to trading.} $\rightarrow$ \eng{I like farming better than trading.} — Gérondif après \eng{like} et après \eng{than}.
\item \eng{I would rather stay in the village.} $\rightarrow$ \eng{I prefer staying in the village.} (ou \eng{I prefer to stay in the village.}) — \eng{Prefer} + gérondif ou infinitif.
\item \eng{She likes teaching better than selling.} $\rightarrow$ \eng{She would prefer to teach rather than sell.} — Infinitif avec \eng{to} après \eng{would prefer}, puis base verbale après \eng{rather than}.
\item \eng{I prefer to work for myself rather than for a company.} $\rightarrow$ \eng{I would rather work for myself than for a company.} — \eng{Would rather + base verbale + than}.
\end{enumerate}
\textbf{Point de vigilance :} dans une transformation, vérifier que le sens reste identique et que la forme verbale correspond exactement à la structure imposée : c'est la forme qui est notée au Bac.
\end{corrige}

\begin{corrige}[Exercice 7 — Texte à trous : l'économie camerounaise]
\textbf{Texte corrigé :}
\par\medskip\noindent
\textit{Many young Cameroonians \textbf{prefer} working in the informal sector \textbf{to} looking for a government job, because it is faster to start. However, others \textbf{would} rather have a stable salary. “I like trading \textbf{better} than farming,” says Alain, a shopkeeper in Douala, “because I earn money every day.” His sister, on the other hand, would \textbf{rather} become a civil servant. “I would prefer \textbf{to} work for the state,” she explains. “It is safer \textbf{than} running a small business.”}
\par\medskip
\textbf{Justifications :}
\begin{itemize}
\item \eng{prefer ... to ...} : préposition \eng{to} entre deux gérondifs (\eng{working} / \eng{looking}).
\item \eng{would rather have} : base verbale \eng{have} après \eng{would rather}.
\item \eng{like ... better than ...} : comparatif \eng{better}.
\item \eng{would rather become} : base verbale \eng{become}.
\item \eng{would prefer to work} : infinitif avec \eng{to}.
\item \eng{safer than} : \eng{than} suit toujours un comparatif.
\end{itemize}
\end{corrige}

\begin{corrige}[Exercice 8 — Appariements et correction d'erreurs]
\textbf{A. Appariement.}
\begin{center}
\begin{tabularx}{\linewidth}{|X|X|}
\hline
\rowcolor{popblueL} Début & Fin \\
\hline
1. I prefer rice & c) to fufu corn. \\
\hline
2. I'd rather take a bendskin & e) than wait for a taxi in the rain. \\
\hline
3. She would prefer to study & b) medicine at the university. \\
\hline
4. We like local food better & a) than imported food. \\
\hline
5. I'd rather you & d) told me the truth about your results. \\
\hline
\end{tabularx}
\end{center}
Vérification des structures : 1 = \eng{prefer + nom + to + nom} ; 2 = \eng{would rather + base verbale ... than + base verbale} ; 3 = \eng{would prefer + to + infinitif} ; 4 = \eng{like ... better than ...} ; 5 = \eng{would rather + sujet + prétérit} (\eng{told}).
\medskip

\textbf{B. Correction d'erreurs.} Les cinq fautes :
\begin{enumerate}
\item \eng{I prefer to be doctor} $\rightarrow$ \eng{I prefer to be \textbf{a} doctor} — article indéfini obligatoire devant un métier au singulier.
\item \eng{to be doctor \textbf{than} to be a teacher} $\rightarrow$ \eng{... \textbf{rather than} to be a teacher} — avec \eng{prefer to}, la seconde option s'introduit par \eng{rather than}, pas \eng{than} seul.
\item \eng{I would rather \textbf{to work}} $\rightarrow$ \eng{I would rather \textbf{work}} — base verbale sans \eng{to} après \eng{would rather}.
\item \eng{I \textbf{like better helping} people} $\rightarrow$ \eng{I \textbf{like helping} people \textbf{better}} — ordre des mots : \eng{like + -ing + better than ...}.
\item \eng{I \textbf{am agree} with them} $\rightarrow$ \eng{I \textbf{agree} with them} — faux ami : « être d'accord » = \eng{agree} (verbe), sans \eng{be}.
\end{enumerate}
\textbf{Texte corrigé :}
\par\medskip\noindent
\textit{I prefer to be a doctor rather than to be a teacher. I would rather work in a hospital because I like helping people better than selling things. My parents prefer that I choose this career, and I agree with them.}
\par\medskip
\end{corrige}

\begin{corrige}[Exercice 9 — Compréhension et langue — type Bac]
\textbf{Barème indicatif (sur 20) :} A. Comprehension : 8 pts (2 pts par réponse : 1 pt sens + 1 pt langue) ; B. Language work : 12 pts (4 + 2 + 2 + 4).
\medskip

\textbf{A. Comprehension — réponses modèles :}
\begin{enumerate}
\item \eng{Manka's father wanted her to study agriculture and take over the family farm, because he believes that farming feeds the nation.}
\item \eng{She prefers trading to farming because you make money faster and you meet people every day.}
\item \eng{Manka prefers a career in computer science, and she wants to study in Yaoundé.}
\item \eng{Yes, they finally accept her decision. Her father says: “We would rather see you happy in a career you love than unhappy in one we chose for you.”}
\end{enumerate}
\textbf{B. Language work — réponses modèles :}
\begin{enumerate}
\item Two structures: \eng{“I prefer trading to farming”} (\eng{prefer ... to ...}) and \eng{“I would prefer to study computer science ... rather than stay here”} (\eng{would prefer to ... rather than ...}). On accepte aussi \eng{“I like working with computers better than ...”} et \eng{“I would rather you took over ...”}.
\item \eng{I would rather trade than farm.} (ou \eng{I would rather be a trader than a farmer.})
\item \eng{I prefer working with computers to working in the fields.}
\item After \eng{would rather + subject}, English uses a past tense (\eng{took}) to talk about a present or future wish. The father is not talking about the past; he is expressing what he wants his daughter to do now or in the future.
\end{enumerate}
\textbf{Points de vigilance :} en compréhension, répondre par des phrases complètes sans recopier tout le texte ; en question B.4, ne pas dire que \eng{took} exprime le passé — c'est un prétérit de souhait (subjonctif français : « que tu reprennes »).
\end{corrige}

\begin{corrige}[Exercice 10 — Traduction — type Bac]
\textbf{Barème indicatif (sur 10) :} 2 pts par phrase (1 pt sens, 1 pt grammaire/lexique).
\medskip

\textbf{Traduction modèle :}
\par\medskip\noindent
\textit{Many young Cameroonians prefer looking for a job in town to staying in the village. They like modern jobs better than traditional farming. However, some would rather become farmers, because they think that the land can make them rich. “I would rather my children chose a job they love,” says a father in Bafoussam.}
\par\medskip
\textbf{Justifications grammaticales :}
\begin{itemize}
\item « préfèrent chercher ... plutôt que de rester » $\rightarrow$ \eng{prefer looking ... to staying} : gérondifs parallèles, préposition \eng{to} (jamais \eng{than}).
\item « aiment mieux ... que » $\rightarrow$ \eng{like ... better than ...}.
\item « préféreraient devenir » $\rightarrow$ \eng{would rather become} (base verbale) ou \eng{would prefer to become}.
\item « J'aimerais mieux que mes enfants choisissent » $\rightarrow$ \eng{I would rather my children chose} : prétérit après \eng{would rather + sujet} pour rendre le subjonctif.
\end{itemize}
\textbf{Points de vigilance :} \eng{town} (« en ville ») sans article ; \eng{the village} avec article ; ne pas traduire « la terre » par \eng{the earth} (qui signifie « la planète ») mais par \eng{the land}.
\end{corrige}

\begin{corrige}[Exercice 11 — Expression écrite — type Bac]
\textbf{Barème indicatif (sur 20) :} respect du sujet et de la longueur (4 pts) ; utilisation correcte d'au moins trois structures de préférence (6 pts) ; richesse du vocabulaire lié aux métiers (4 pts) ; grammaire, orthographe, ponctuation (4 pts) ; organisation et cohérence (2 pts).
\medskip

\textbf{Proposition de rédaction modèle (environ 110 mots) :}
\par\medskip\noindent
\textbf{Modèle :}
\par\medskip\noindent
\textit{Like many young Cameroonians, I have thought carefully about my future career. \textbf{I prefer working with my head to working with my hands}, so \textbf{I would rather become a computer engineer than a mechanic}. New technologies are developing fast in cities like Yaoundé and Douala, and good programmers are needed everywhere. \textbf{I would prefer to work for a private company} rather than wait for a government job, because I want to be independent. Of course, I respect farmers and traders: our country needs them. But \textbf{I like solving problems on a computer better than selling goods in a market}. In short, \textbf{I would rather choose a career I love than follow a path chosen by others}.}
\par\medskip
\textbf{Structures de préférence utilisées (à souligner) :} \eng{prefer working ... to working ...} ; \eng{would rather become ... than ...} ; \eng{would prefer to work ...} ; \eng{like ... better than ...} ; \eng{would rather choose ... than follow ...}.
\medskip

\textbf{Points de vigilance pour l'élève :}
\begin{itemize}
\item Vérifier chaque structure : base verbale après \eng{would rather}, infinitif avec \eng{to} après \eng{would prefer}, préposition \eng{to} (et non \eng{than}) après \eng{prefer}.
\item Article obligatoire devant les métiers : \eng{a computer engineer}, \eng{a mechanic}.
\item Éviter les calques du français : « j'aime mieux » $\rightarrow$ \eng{I like ... better}, jamais \eng{I more like}.
\item Rester dans la limite des 100 mots environ et relire pour éliminer les fautes d'accord (\eng{he wants}, pas \eng{he want}).
\end{itemize}
\end{corrige}

\newpage

\coursec{Fiche bilan}

\begin{popfigure}
\centering
\begin{tikzpicture}[
    mindmap,
    grow cyclic,
    every node/.style={concept, circular drop shadow, execute at begin node=\hskip0pt, font=\small\bfseries},
    root concept/.append style={concept color=popblue, font=\large\bfseries, text width=3.5cm, align=center},
    level 1 concept/.append style={level distance=4.2cm, sibling angle=60, text width=2.8cm, align=center, font=\small\bfseries},
    concept color=popblue!70
]
\node[root concept, align=center]{Expressing\\ Preferences}
    child[concept color=popblue] { node[concept, align=center]{Prefer\\ + \eng{-ing} / \eng{to} / noun} }
    child[concept color=poporange] { node[concept, align=center]{Would rather\\ + bare inf. / past} }
    child[concept color=popgreen] { node[concept, align=center]{Would prefer\\ + \eng{to}-inf. / noun} }
    child[concept color=poppurple] { node[concept, align=center]{Like / Love\\ + better than} }
    child[concept color=poppink] { node[concept, align=center]{Had better\\ (strong advice)} }
    child[concept color=popgold] { node[concept, align=center]{Traps FR/EN\\ \& Transformations} };
\end{tikzpicture}
\legende{Carte mentale : Les structures pour exprimer la préférence en anglais}
\end{popfigure}

\begin{bilan}
\end{bilan}

\courssub{Lexique clé (Vocabulary)}

\begin{center}
\begin{tabularx}{\linewidth}{|X|X|X|}
\hline
\rowcolor{popblueL} \textbf{English} & \textbf{Français} & \textbf{Remarque / Collocation} \\
\hline
\eng{preference} & préférence & \eng{to have a preference for} \\
\eng{rather} & plutôt & adv. ; \eng{would rather} = préférer \\
\eng{prefer} & préférer & v. transitif direct (pas \eng{prefer \textbf{to}} au sens de « préférer ») \\
\eng{better} & mieux / meilleur & compar. de \eng{well/good} ; \eng{like ... better} \\
\eng{fancy} & avoir envie de / aimer bien & informel, surtout UK ; + \eng{-ing} \\
\eng{favour} (US \eng{favor}) & favoriser / préférer & formel ; \eng{to favour doing sth} \\
\eng{opt for} & opter pour & + gérondif ou nom \\
\hline
\end{tabularx}
\end{center}

\courssub{Règles de grammaire essentielles (à connaître par cœur)}

\begin{center}
\begin{tabularx}{\linewidth}{|>{\bfseries}l|X|X|}
\hline
\rowcolor{popblueL} \textbf{Structure} & \textbf{Forme / Règle} & \textbf{Exemple + Traduction} \\
\hline
\eng{Prefer} & \eng{prefer + -ing} (préférence générale) & \eng{I prefer teaching to marking.} « Je préfère enseigner à corriger. » \\
& \eng{prefer + to-inf} (choix précis, futur) & \eng{I prefer to teach today.} « Je préfère enseigner aujourd'hui. » \\
& \eng{prefer + noun/pronoun} & \eng{I prefer coffee.} « Je préfère le café. » \\
& \textbf{Comparaison} : \eng{prefer A to B} (pas \eng{than}) & \eng{I prefer tea to coffee.} « Je préfère le thé au café. » \\
\hline
\eng{Would rather} (\eng{'d rather}) & \textbf{Même sujet} : \eng{would rather + bare infinitive} & \eng{I'd rather go now.} « Je préférerais partir maintenant. » \\
& \textbf{Sujet différent} : \eng{would rather + sujet + prétérit (subjonctif)} & \eng{I'd rather you went now.} « Je préférerais que tu partes maintenant. » \\
& \textbf{Négation} : \eng{would rather not + bare inf.} & \eng{I'd rather not say.} « Je préférerais ne rien dire. » \\
\hline
\eng{Would prefer} (\eng{'d prefer}) & \eng{would prefer + to-inf} (choix précis) & \eng{I'd prefer to stay home.} « Je préférerais rester à la maison. » \\
& \eng{would prefer + noun} & \eng{I'd prefer the blue one.} « Je préférerais le bleu. » \\
& \textbf{Comparaison} : \eng{would prefer A rather than B} & \eng{I'd prefer to walk rather than drive.} « Je préférerais marcher plutôt que conduire. » \\
\hline
\eng{Like / Love / Hate ... better than} & \eng{like/love/hate + noun/-ing + better than + noun/-ing} & \eng{I like reading better than watching TV.} « J'aime lire mieux que regarder la télé. » \\
\hline
\eng{Had better} (\eng{'d better}) & \textbf{Conseil fort / Avertissement} : \eng{had better + bare inf.} (pas \eng{would}) & \eng{You'd better revise.} « Tu ferais mieux de réviser. » \\
& Négation : \eng{had better not + bare inf.} & \eng{You'd better not be late.} « Tu ferais mieux de ne pas être en retard. » \\
\hline
\end{tabularx}
\end{center}

\courssub{Méthodes express}

\begin{itemize}
    \item \textbf{Choisir la bonne structure :} Général/ Habitude $\rightarrow$ \eng{prefer + -ing} / \eng{like ... better}. Choix précis / Futur $\rightarrow$ \eng{would prefer / would rather + inf}. Conseil $\rightarrow$ \eng{had better}.
    \item \textbf{Transformer \eng{Prefer} en \eng{Would rather} :} \eng{I prefer tea to coffee} $\rightarrow$ \eng{I'd rather have tea than coffee}. Attention : \eng{to} devient \eng{than}, \eng{-ing} devient \eng{bare inf}.
    \item \textbf{Gérer le changement de sujet (Would rather) :} Si le sujet change, passer au \textbf{prétérit} (valeur subjonctive). \eng{I'd rather you \textbf{came}} (pas \eng{come}).
    \item \textbf{Négation :} \eng{Prefer not to} / \eng{Would rather not} / \eng{Had better not} (jamais \eng{don't prefer}, \eng{don't rather}).
\end{itemize}


\begin{aretenir}[Checklist avant le Bac]
\begin{itemize}
    \item $\square$ Je sais conjuguer \eng{prefer} (régulier) et utiliser \eng{would rather / would prefer / had better} (modaux/semi-modaux invariables).
    \item $\square$ Je distingue \textbf{préférence générale} (\eng{prefer + -ing}) et \textbf{choix ponctuel} (\eng{would prefer / would rather + inf}).
    \item $\square$ Je maîtrise la préposition de comparaison : \eng{prefer A \textbf{to} B} (jamais \eng{than}).
    \item $\square$ Je construis correctement \eng{Would rather} avec \textbf{changement de sujet} : \eng{S1 would rather S2 + \textbf{prétérit}}.
    \item $\square$ Je place la négation au bon endroit : \eng{would rather \textbf{not} + bare inf} / \eng{had better \textbf{not} + bare inf}.
    \item $\square$ Je ne confonds pas \eng{Would rather} (préférence) et \eng{Had better} (conseil/obligation morale).
    \item $\square$ Je sais transformer : \eng{Prefer A to B} $\leftrightarrow$ \eng{Would rather A than B} / \eng{Would prefer A rather than B}.
    \item $\square$ J'évite l'erreur \eng{*I prefer to go than to stay} $\rightarrow$ \eng{I prefer going to staying} ou \eng{I'd rather go than stay}.
    \item $\square$ J'utilise \eng{Like ... better than} pour les goûts (équivalent \eng{prefer} général) avec parallélisme (\eng{-ing} vs \eng{-ing}).
    \item $\square$ Je reconnais \eng{Fancy + -ing} (registre familier UK) et \eng{Favour + -ing} (registre formel).
    \item $\square$ Je sais exprimer « Je n'ai pas de préférence » : \eng{I don't mind} / \eng{I have no preference} / \eng{It's all the same to me}.
    \item $\square$ Je relis mes phrases pour vérifier l'ordre des mots (sujet + modal + base verbale) et l'absence de \eng{to} après \eng{rather / better}.
\end{itemize}
\end{aretenir}

\findecours

\end{document}
